%%%%%%%%%%%%%%%%%%%%%%%%%%%%%%%%%%%%%%%%%
% Beamer Presentation
% LaTeX Template
% Version 1.0 (10/11/12)
%
% This template has been downloaded from:
% http://www.LaTeXTemplates.com
%
% License:
% CC BY-NC-SA 3.0 (http://creativecommons.org/licenses/by-nc-sa/3.0/)
%
%%%%%%%%%%%%%%%%%%%%%%%%%%%%%%%%%%%%%%%%%

%----------------------------------------------------------------------------------------
%	PACKAGES AND THEMES
%----------------------------------------------------------------------------------------

\documentclass[aspectratio=169]{beamer}

\mode<presentation> {


\definecolor{LightBlue}{RGB}{93,111,122}
\definecolor{TitleCopper}{RGB}{0,108,92}
\definecolor{emphasis}{RGB}{212,0,0}
\definecolor{darkgreen}{rgb}{0,0.5,0}
% The Beamer class comes with a number of default slide themes
% which change the colors and layouts of slides. Below this is a list
% of all the themes, uncomment each in turn to see what they look like.

\usetheme{Madrid}
% Change base colour beamer@blendedblue (originally RGB: 0.2,0.2,0.7)
\colorlet{beamer@blendedblue}{blue!80!black}

% As well as themes, the Beamer class has a number of color themes
% for any slide theme. Uncomment each of these in turn to see how it
% changes the colors of your current slide theme.

%\setbeamertemplate{footline} % To remove the footer line in all slides uncomment this line
%\setbeamertemplate{footline}[page number] % To replace the footer line in all slides with a simple slide count uncomment this line

%\setbeamertemplate{navigation symbols}{} % To remove the navigation symbols from the bottom of all slides uncomment this line
}

\newcommand\Fontvi{\fontsize{10}{10}\selectfont}
\newcommand{\norm}[1]{\left\lVert#1\right\rVert_2}
\newcommand{\abs}[1]{\left|#1\right|}
\newcommand\emphasis[1]{{\textbf{\color{emphasis}#1}}}

% Toggle animations on or off for faster build times
\newif\ifshowanimations
% \showanimationstrue
\showanimationsfalse

% Toggle animations on or off for faster build times
\newif\ifshowcitations
\showcitationstrue
% \showcitationsfalse

% for zero-padded image numbers inside animations
\newcommand\zeropad[2]{%
  \ifnum#2<0\relax%
    {\ensuremath-}\zeropadA{#1}{\the\numexpr#2*-1\relax}%
  \else%
    \zeropadA{#1}{#2}%
  \fi%
}
\def\zeropadA#1#2{%
  \ifnum1#2<1#1
    \zeropadA{#1}{0#2}%
  \else%
    #2%
  \fi%
}

%==============================================================================%
%% Bibliography settings
\ifshowcitations
	\usepackage[backend=biber,autocite=footnote,citestyle=authoryear,style=mla]{biblatex}

	\newcommand*{\volnumpunct}{\addcomma\space}
	\renewcommand*{\newunitpunct}{\addcomma\space}
	\renewcommand*{\subtitlepunct}{\nopunct\space}
	\renewcommand*{\intitlepunct}{\nopunct\addspace}
	\renewcommand*{\relatedpunct}{\addsemicolon\space}
	\renewcommand*{\finentrypunct}{} 

	\renewcommand*{\bibfootnotewrapper}[1]{\bibsentence #1} % remove final period

	\DeclareCiteCommand{\footpartcite}[\mkbibfootnote]{
		{
			\boolfalse{citetracker} % not sure if this helps
			\usebibmacro{prenote}
		}%
		{\usebibmacro{citeindex}%
			\setunit{\addcomma}
			% \printnames{labelname}%
			% \setunit{\labelnamepunct}
			\printnames[default]{author} % default ordering first name, last name
			\newunit
			\printfield{year}
			\newunit
			{\color{emphasis}\printfield{journaltitle}}
		}
		{\usebibmacro{postnote}}
		\vspace{-1.5em}
	}

	\AtEveryCitekey{\iffootnote{\color{emphasis}\tiny}{\color{blue}}}

	\makeatletter
	\renewcommand\@makefnmark{\textsuperscript{\@thefnmark}} % suppress ref mark
	\renewcommand\@makefntext[1]{\tiny\noindent{[\@thefnmark]}#1} % smaller font size
	\makeatother

	\addbibresource{bibliography.bib}
	% \bibliography{\jobname.bib}
	% \renewcommand*{\thefootnote}{\alph{footnote}} % alpha footnotes
\fi
%==============================================================================%
%% Footer formatting

\setbeamertemplate{footline}
{
	\leavevmode%
	\hbox{%
	\begin{beamercolorbox}[wd=.2\paperwidth,ht=2.25ex,dp=1ex,center]{author in head/foot}%
	\usebeamerfont{author in head/foot}\insertshortauthor
	\end{beamercolorbox}%
	\begin{beamercolorbox}[wd=.6\paperwidth,ht=2.25ex,dp=1ex,center]{title in head/foot}%
	\usebeamerfont{title in head/foot}{\href{https://khbalhandawi.github.io/}{https://khbalhandawi.github.io/}}
	\end{beamercolorbox}%
	\begin{beamercolorbox}[wd=.2\paperwidth,ht=2.25ex,dp=1ex,center]{date in head/foot}%
	\insertdate~~\insertframenumber{} / \inserttotalframenumber\hspace*{1ex}
	\end{beamercolorbox}}%
	\vskip0pt%
}

\usepackage{textpos}
% \addtobeamertemplate{frametitle}{}{%
% 	\begin{textblock*}{100mm}(.94\textwidth,-1.1cm)
% 		\colorbox{white}{\includegraphics[height=1.2cm]{logos/mcgill_logo.pdf}}
% 	\end{textblock*}
% }

%==============================================================================%
%% Reduce TOC spacing

\makeatletter
\patchcmd{\beamer@sectionintoc}{\vskip1.5em}{\vskip0.5em}{}{}
\makeatother

%==============================================================================%
%% Good behaving itemize env with top aligned columns

\makeatletter
\newenvironment{myitemize}{%
   \setlength{\topsep}{0pt}
   \setlength{\partopsep}{0pt}
   \renewcommand*{\@listi}{\leftmargin\leftmargini \parsep\z@ \topsep\z@ \itemsep\z@}
   \let\@listI\@listi
   \itemize
}{\enditemize}
\makeatother  

%==============================================================================%
%% Imported packages

\usepackage{changepage} % adjust slide margins
\usepackage{graphicx} % Allows including images
\usepackage{animate} % For animated images
\graphicspath{ {images/} }
\usepackage{tikz}
\usepackage{tikz-dimline}
\usetikzlibrary{calc,fadings,positioning,arrows,shapes,backgrounds,shapes.geometric,arrows.meta}
\tikzfading[name=fade r,right color=transparent!100,left color=transparent!0]
\tikzstyle{every picture}+=[remember picture]
\tikzstyle{na} = [baseline=-.5ex]
% this "frames" a rectangle node
\newcommand\framenode[2][10pt]{
    \fill[white,path fading=fade r] (#2.south west) rectangle ($(#2.north west)+( #1,0)$);
}
\newcommand{\fixedpicture}[4]{ % height, content
    \node (TRSschematic) at ( [xshift = #1, yshift = #2] current page.center) %or: (current page.west
	{\includegraphics[height=#3]{#4}};
}
\newcommand{\tikzmark}[3][]{\tikz[remember picture,baseline] \node [anchor=base,#1](#2) {$#3$};}

\usepackage{algorithmic}
\usepackage[linesnumbered,ruled,vlined]{algorithm2e} % for algorithms


 % code snippets
\usepackage{listings}
\definecolor{codegreen}{rgb}{0,0.6,0}
\definecolor{codegray}{rgb}{0.5,0.5,0.5}
\definecolor{codepurple}{rgb}{0.58,0,0.82}
\definecolor{backcolour}{rgb}{0.95,0.95,0.92}

\lstdefinestyle{mystyle}{
    backgroundcolor=\color{backcolour},   
    commentstyle=\color{codegreen},
    keywordstyle=\color{magenta},
    numberstyle=\tiny\color{codegray},
    stringstyle=\color{codepurple},
    basicstyle=\ttfamily\tiny,
    breakatwhitespace=false,         
    breaklines=true,                 
    captionpos=b,                    
    keepspaces=true,                 
    numbers=left,                    
    numbersep=1pt,                  
    showspaces=false,                
    showstringspaces=false,
    showtabs=false,                  
    tabsize=2,
	aboveskip=0.0em,
    belowskip=0.0em
}
\lstset{style=mystyle}

\usepackage{tabularx}
\usepackage{multirow}% tables
\usepackage{booktabs} % Allows the use of \toprule, \midrule and \bottomrule in tables
\newcolumntype{L}[1]{>{\raggedright\let\newline\\\arraybackslash\hspace{0pt}}m{#1}}
\newcolumntype{C}[1]{>{\centering\let\newline\\\arraybackslash\hspace{0pt}}m{#1}}
\newcolumntype{R}[1]{>{\raggedleft\let\newline\\\arraybackslash\hspace{0pt}}m{#1}}

\usepackage{lmodern} % load a font with all the characters
\usepackage{amssymb}
\usepackage{amsmath}
%\usepackage{unicode-math}
\usepackage{cmbright}
\usepackage{kbordermatrix}% http://www.hss.caltech.edu/~kcb/TeX/kbordermatrix.sty
%\SetSymbolFont{largesymbols}{normal}{OMX}{iwona}{m}{n}

\usepackage{cleveref}% reference multiple equations
\usepackage{caption}
\usepackage{subcaption}
\usepackage{enumitem} % tick and cross bullet points
\usepackage{amsfonts}
\usepackage{ragged2e} % justify text
\usepackage{multido} % loop over i images
\usepackage{array} % for tables 
\usepackage{varwidth} % for bullets inside tikz nodes
\usepackage{fancybox} % for rounded boxes
\usepackage{xstring} % for string lookup functions 

\usepackage{pifont} % for ding commands
\newcommand{\cmark}{\ding{51}} % for check marks
\newcommand{\xmark}{\ding{55}} % for cross marks
\usepackage{bbding,enumitem} % tick and cross bullet points

\newlength{\mytextsize} % fontsize length for animations inline
%This part fails.... 
\makeatletter
	\setlength{\mytextsize}{\f@size pt}
\makeatother

\setitemize{label=\usebeamerfont*{itemize item}% Restore original beamer style
\usebeamercolor[fg]{itemize item}
\usebeamertemplate{itemize item}}

\setbeamerfont{block body}{size=\scriptsize}
\setbeamerfont{block body alerted}{size=\scriptsize}
\setbeamerfont{block body example}{size=\scriptsize}
\captionsetup{textfont={scriptsize,it},labelformat=empty,labelsep=none}
\usepackage{multicol} % multiple columns for itemize lists

% INTEGRATED PROBLEM BLOCK %
\setbeamercolor{block title}{fg=white,bg=blue}
\setbeamercolor{block body}{fg=black, bg=blue!20!white}

% INTEGRATED PROBLEM BLOCK %
\setbeamercolor{block title}{fg=white,bg=blue}
\setbeamercolor{block body}{fg=black, bg=blue!20!white}

% INTEGRATED PROBLEM BLOCK %
\newenvironment<>{blueblock}[1]{%
\setbeamercolor{itemize item}{fg=blue}
\setbeamercolor{block title}{fg=white,bg=blue}
\setbeamercolor{block body}{fg=black, bg=blue!20!white}
	\begin{block}#2{#1}}{\end{block}}

% TITLE BLOCK %
\newenvironment<>{titleblock}[1]{%
\setbeamercolor{block title}{fg=black,bg=red!20!white}
\setbeamercolor{block body}{fg=black, bg=red!20!white}
	\begin{block}#2{#1}}{\end{block}}

% well behaved centering environment
\newenvironment{nscenter}
 {\parskip=0pt\par\nopagebreak\centering}
 {\par\noindent\ignorespacesafterend}

% Image scales and sizes
\newlength\imagewidth
\newlength\imagescale
\pgfmathsetlength{\imagewidth}{6cm}
\pgfmathsetlength{\imagescale}{\imagewidth/600}

% https://tex.stackexchange.com/a/28633
\makeatletter
% Use Euclid's Algorithm to find the greatest 
% common divisor of two integers.
\def\gcd#1#2{{% #1 = a, #2 = b
    \ifnum#2=0 \edef\next{#1}\else
        \@tempcnta=#1 \@tempcntb=#2 \divide\@tempcnta by\@tempcntb
        \multiply\@tempcnta by\@tempcntb  % q*b
        \@tempcntb=#1
        \advance\@tempcntb by-\@tempcnta % remainder in \@tempcntb
        \ifnum\@tempcntb=0
            \@tempcnta=#2
            \ifnum\@tempcnta < 0 \@tempcnta=-\@tempcnta\fi
            \xdef\gcd@next{\noexpand%
                \def\noexpand\thegcd{\the\@tempcnta}}%
        \else
            \xdef\gcd@next{\noexpand\gcd{#2}{\the\@tempcntb}}%
        \fi
    \fi}\gcd@next
}
\newcommand\reduceFrac[2]
{%
   \gcd{#1}{#2}{\@tempcnta=#1 \divide\@tempcnta by\thegcd
   \@tempcntb=#2 \divide\@tempcntb by\thegcd
   \ifnum\@tempcntb<0\relax
      \@tempcntb=-\@tempcntb
        \@tempcnta=-\@tempcnta
    \fi
    \xdef\rfNumer{\the\@tempcnta}
    \xdef\rfDenom{\the\@tempcntb}}%
}
\makeatother

\newcommand*{\fracReduced}[2]{%
    \reduceFrac{#1}{#2}%
    \ensuremath{%
        \ifnum\rfDenom=1
            \rfNumer
        \else
            \frac{\rfNumer}{\rfDenom}%
        \fi
    }%
}%


\newcommand{\addfraction}[2]{% 1N, 1D, 2N, 2D
    \pgfmathsetmacro{\newnumerator}{#1}
    \pgfmathsetmacro{\newdenominator}{#2}
    \pgfmathtruncatemacro{\newnumeratorTrunc}{\newnumerator}
    \pgfmathtruncatemacro{\newdenominatorTrunc}{\newdenominator}
    \pgfkeys{/pgf/number format/.cd,int detect,precision=2}
    \fracReduced{\newnumeratorTrunc}{\newdenominatorTrunc}
}
%----------------------------------------------------------------------------------------
%	TITLE PAGE
%----------------------------------------------------------------------------------------

\title[]{Overview of my research projects} % The short title appears at the bottom of every slide, the full title is only on the title page

\author[K. Al Handawi]{Khalil Al Handawi, PhD} % Your name

\institute[McGill]% Your institution as it will appear on the bottom of every slide, may be shorthand to save space
{
	~\\
	\textit{\small
	Technical presentation\\} % Your institution for the title page
}
\date{\today} % Date, can be changed to a custom date

\begin{document}

\begin{frame}
	% \begin{tikzpicture}[remember picture,overlay]
	% \node[opacity=0.9] (image) at
	% 	([xshift=-0.125\textwidth, yshift =-1cm]current page.east) 	%or: (current page.center)
	% 	{\includegraphics[height=1.6\textheight]{logos/mcgill_logo.pdf}};
	% 	\framenode[600pt]{image}
	% \end{tikzpicture}

	\titlepage % Print the title page as the first slide
\end{frame}

%----------------------------------------------------------------------------------------
% \begin{frame}
% \frametitle{Overview} % Table of contents slide, comment this block out to remove it
% \tableofcontents[hideallsubsections] % Throughout your presentation, if you choose to use \section{} and \subsection{} commands, these will automatically be printed on this slide as an overview of your presentation
% \end{frame}
%----------------------------------------------------------------------------------------

%----------------------------------------------------------------------------------------
%	PRESENTATION SLIDES
%----------------------------------------------------------------------------------------

% Section 00: Introduction
%%%%%%%%%%%%%%%%%%%%%%%%%%%%%%%%%%%%%%%%%%%%%%%%%%%%%%%
%%                   Introduction                    %%
%%%%%%%%%%%%%%%%%%%%%%%%%%%%%%%%%%%%%%%%%%%%%%%%%%%%%%%
%------------------------------------------------
\section{About me}
%------------------------------------------------
\begin{frame}
\frametitle{About me}
\begin{columns}[t] % The "c" option specifies centered vertical alignment while the "t" option is used for top vertical alignment

	\begin{column}{.5\textwidth} % Left column and width
		\begin{itemize}
			\item Khalil Al Handawi (\textit{he/him/il})
			\begin{itemize}
				\item Experience with simulation based design
				\item Experience with scientific computing
				\item Experience with optimization and simulation research
			\end{itemize}
			\item Born and raised in Abu Dhabi, United Arab Emirates
			\item Moved to Montr\'{e}al in 2017
			\item Obtained a doctorate in mechanical engineering from McGill University
		\end{itemize}
	\end{column}

	\begin{column}{.5\textwidth} % Left column and width

		\begin{center}
			\begin{tikzpicture}[x=\imagescale,y=-\imagescale]
				\clip (600/2, 567/2) circle (567/2);
				\node[anchor=north west, inner sep=0pt, outer sep=0pt] at (0,0) {\includegraphics[width=\imagewidth]{profilepic.jpg}};
			\end{tikzpicture}
		\end{center}
			
	\end{column}
\end{columns}

\end{frame}
%------------------------------------------------

% Section 01: scalability
%%%%%%%%%%%%%%%%%%%%%%%%%%%%%%%%%%%%%%%%%%%%%%%%%%%%%%%
%%                   Introduction                    %%
%%%%%%%%%%%%%%%%%%%%%%%%%%%%%%%%%%%%%%%%%%%%%%%%%%%%%%%
%------------------------------------------------
\section{Introduction}
%------------------------------------------------
\begin{frame}[t]
	\frametitle{The challenge}
	
	\begin{itemize}
		\uncover<1->{\item This is a Turbine Rear Structure (TRS) under thermal loading}
		\uncover<2->{\item Over its lifetime, design requirements can change}
	\end{itemize}
	
	\vspace{-2em}
	\begin{columns}[t] % The "c" option specifies centered vertical alignment while the "t" option is used for top vertical alignment
			
		\begin{column}{.4\textwidth} % Left column and width
			\begin{figure}
				\only<0>{
					\includegraphics[height=5cm]{scalability/TRS_overview_pictures/1a_TRS_isometric_1.pdf}
				}%
				\only<1>{
					\includegraphics[height=5cm]{scalability/TRS_overview_pictures/1a_TRS_isometric_2.pdf}
				}%
				\only<2->{
					\includegraphics[height=5cm]{scalability/TRS_overview_pictures/1a_TRS_isometric_3.pdf}
				}%
			\end{figure}
			\vspace{-3em}
		\end{column}
	
		\begin{column}{.5\textwidth} % Left column and width
	
			\begin{figure}
				\centering\includegraphics[height=5cm]{scalability/trent_900.pdf}
				\caption{An aeroengine}
			\end{figure}%
				
		\end{column}
	\end{columns}
\end{frame}
%------------------------------------------------
\begin{frame}[t]
	\frametitle{The opportunity}
	
	\begin{itemize}
		\item Disposal or redesign can be sub-optimal options
		\uncover<2->{\item Additive remanufacturing may be a better option to extend useful lifetime}
	\end{itemize}
	
	\vspace{-2em}
	\begin{columns}[t] % The "c" option specifies centered vertical alignment while the "t" option is used for top vertical alignment
			
		\begin{column}{.4\textwidth} % Left column and width
			\begin{figure}
				\includegraphics[height=5cm]{scalability/TRS_overview_pictures/1a_TRS_isometric_3.pdf}
			\end{figure}
			\vspace{-3em}
		\end{column}
	
		\begin{column}{.5\textwidth} % Left column and width
	
			\begin{figure}
				\only<1>{
					\vspace{2em}
					\centering\includegraphics[height=4cm]{scalability/dispose_guy.jpg}
				}%
				\only<2>{
					\centering
					\begin{animateinline}[autoplay,height=5cm]{5}
						\ifshowanimations
							\multiframe{30}{i=1+6}{%
								\includegraphics{scalability/TEMP_GIF/frame_\i.png}
							}
						\else
							\multiframe{1}{i=180+0}{%
								\includegraphics{scalability/TEMP_GIF/frame_\i.png}
							}
						\fi
					\end{animateinline}
					\caption{AM Remanufacturing of TRS}
				}%
			\end{figure}%
				
		\end{column}
	\end{columns}
\end{frame}
%------------------------------------------------
% https://tex.stackexchange.com/a/136166
\tikzset{
  transparent/.style={opacity=0.4},
  untransparent on/.style={alt={#1{}{transparent}}},
  alt/.code args={<#1>#2#3}{%
    \alt<#1>{\pgfkeysalso{#2}}{\pgfkeysalso{#3}} % \pgfkeysalso doesn't change the path
  },
  invisible/.style={opacity=0.0},
  visible on/.style={alt={#1{}{invisible}}},
  alt/.code args={<#1>#2#3}{%
    \alt<#1>{\pgfkeysalso{#2}}{\pgfkeysalso{#3}} % \pgfkeysalso doesn't change the path
  },
}
\begin{frame}[t]
	\newcommand{\yshift}{0.0cm}
	\frametitle{Change absorption through scalability}
    %
    \tikzstyle{background grid}=[draw, black!50,step=.5cm]
    \begin{tikzpicture}[remember picture, overlay] %show background grid, 
        \node [inner sep=0pt,above right, opacity=1.0]  at (7cm,-5cm+\yshift) (original) 
        {
            \includegraphics[height=3.5cm]{scalability/intro_definitions/TRS_isometric_A.pdf}
        };
        \node [inner sep=0pt,above right, opacity=1.0]  at (7cm,-5cm+\yshift) (loads) 
        {
            \includegraphics[height=3.5cm]{scalability/intro_definitions/arrows_1.pdf}
        };
		\node [inner sep=0pt,above right, opacity=1.0, visible on=<2-4>]  at (7cm,-5cm+\yshift) (loads2) 
        {
            \includegraphics[height=3.5cm]{scalability/intro_definitions/arrows_2.pdf}
        };
        \only<4->{
            \node [inner sep=0pt,above right, untransparent on=<5>]  at (12cm,-3cm+\yshift) (trsA) 
            {
                \includegraphics[height=3.5cm]{scalability/intro_definitions/TRS_isometric_B.pdf}
            };
            \node [inner sep=0pt,above right, untransparent on=<5>]  at (12cm,-7cm+\yshift) (trsB)
            {
                \includegraphics[height=3.5cm]{scalability/intro_definitions/TRS_isometric_C.pdf}
            };%
            \draw[,->,thick,untransparent on=<5->] (original.east) to [in=180,out=0] (trsA.west);
            \draw[,->,thick,untransparent on=<5->] (original.east) to [in=180,out=0] (trsB.west);
        }%
		\node[%
        visible on=<5>,
            draw=green,
            fill=green!10,
            cloud,
            line width=1mm,
            font=\fontfamily{ppl}\fontsize{0.5cm}{0.5cm}\selectfont
        ] at (11cm,-5.25cm+\yshift) {\cmark};
        \node[%
            visible on=<5>,
            draw=red,
            fill=red!10,
            cloud,
            line width=1mm,
            font=\fontfamily{ppl}\fontsize{0.5cm}{0.5cm}\selectfont
        ] at (11cm,-1.25cm+\yshift) {\xmark};
        % show origin
        % \fill (0,0) circle (2pt);
    \end{tikzpicture}%
    %
	\begin{columns}[c]
		\begin{column}{.5\textwidth} % Left column and width
			~\\
			~\\
			~\\
           	When the applied \only<1>{loads change}\only<2->{\emphasis{loads change}}:\\%
			%
			\uncover<3->{A design can be \emphasis{changed}:}
			%
			\begin{itemize}
				\item<4-> \textit{Scalability}: the ability to change to accommodate stricter requirements \ifshowcitations\footnotemark[1]\fi
				\item<5-> Successful \textit{remanufacturing}: enabled by scalability of components \ifshowcitations\footnotemark[2]\fi
			\end{itemize}
			%
		\end{column}
		\begin{column}{.5\textwidth} % Left column and width
		\end{column}
	\end{columns}
	%
	{\color{white}
	\ifshowcitations
		\only<4->{\footpartcite{Ross2008}[1]}
		\only<5->{\footpartcite{Xing2007}[2]}
	\fi
	}
\end{frame}
\addtocounter{footnote}{-2}

%%%%%%%%%%%%%%%%%%%%%%%%%%%%%%%%%%%%%%%%%%%%%%%%%%%%%%%
%%                   Design problem                  %%
%%%%%%%%%%%%%%%%%%%%%%%%%%%%%%%%%%%%%%%%%%%%%%%%%%%%%%%
%------------------------------------------------
\section{Parametric optimization problem formulation}
% \frame{\tableofcontents[
% 	currentsection, 
% 	hideallsubsections
% 	]}
%------------------------------------------------
\begin{frame}[c,noframenumbering]
	\centering
	% \setlength\fboxsep{0pt}
	\begin{titleblock}{}
		~\\%
		{\centering\LARGE Change absorption through scalability and remanufacturing \ifshowcitations\footnotemark[1]\fi\\}%
		~\\%
	\end{titleblock}
	{\color{white}\ifshowcitations\footpartcite{Alhandawi2020}[1]\fi}
\end{frame}
\addtocounter{footnote}{-1}
%------------------------------------------------
\subsection{Obtain set of parametric optimal designs using surrogate}
%------------------------------------------------
\begin{frame}[t]
\frametitle{Parametric design optimization problem formulation}
%
Remanufacturing variables are included in the design optimization problems and changing requirements are modeled by parameters $\mathbf{p}$ \\~\\

\vspace{-0.5em}
\begin{columns}[t] % The "c" option specifies centered vertical alignment while the "t" option is used for top vertical alignment
	\column{.4\textwidth} % Left column and width
		\vspace{-1em}
		% Optimization problem
		\begin{alertblock}{Design optimization problem}
			\vspace{-1.5em}
			\begin{equation*}
				\begin{aligned}
					& \underset{\mathbf{x}}{\text{minimize}}
					& & {f}(\mathbf{x};\mathbf{p}) = -n_{\textrm{safety}}(\mathbf{x};P_{\textrm{load}})\\
					& \text{subject to}
					& & {g_1}(\mathbf{x};\mathbf{p}) = x_3 + x_1 - W_{\textrm{total}} \le 0\\
					&&& {g_2}(\mathbf{x};\mathbf{p}) = T_m - T_{\textrm{deposit}} \le 0\\
				\end{aligned}
			\end{equation*}
		\end{alertblock}
		% Parameters
		\vspace{-1em}
		\only<1->{
		\begin{exampleblock}{Changing requirements}
			\vspace{-1.5em}
			\begin{equation*}
				\begin{aligned}
					& & & W_{\textrm{total}} = \textrm{width of base part}\\
					& & & T_{m} = \textrm{Melting temperature of deposit}\\
					& & & P_{\textrm{load}} = \textrm{Pressure load on outer casing}\\
				\end{aligned}
			\end{equation*}
		\end{exampleblock}
		}%
		
	\column{.6\textwidth} % Right column and width
		
		\begin{figure}
			\includegraphics[height=5cm]{scalability/TRS_overview_pictures/1b_TRS_overview_3.pdf}
			\caption{TRS cross-section}
		\end{figure}
	\end{columns}

\end{frame}

%%%%%%%%%%%%%%%%%%%%%%%%%%%%%%%%%%%%%%%%%%%%%%%%%%%%%%%
%%              Parametric optimization              %%
%%%%%%%%%%%%%%%%%%%%%%%%%%%%%%%%%%%%%%%%%%%%%%%%%%%%%%%
%------------------------------------------------
\section{Scalability Assessment}
% \frame{\tableofcontents[
% 	currentsection, 
% 	hideallsubsections
% 	]}
%------------------------------------------------
\subsection{Conditions for scalability}
%------------------------------------------------
\begin{frame}[t]
\frametitle{Conditions for scalability}
%
\vspace{-0.5em}
\begin{columns}[t] % The "c" option specifies centered vertical alignment while the "t" option is used for top vertical alignment
	\column{.4\textwidth} % Left column and width
		\vspace{-1em}
		% Optimization problem
		\begin{alertblock}{Design optimization problem}
			\vspace{-1.5em}
			\begin{equation*}
				\begin{aligned}
					& \underset{\mathbf{x}}{\text{minimize}}
					& & {f}(\mathbf{x};\mathbf{p}) = -n_{\textrm{safety}}(\mathbf{x};{P_{\textrm{load}}})\\
					& \text{subject to}
					& & {g_1}(\mathbf{x};\mathbf{p}) = x_3 + x_1 - {W_{\textrm{total}}} \le 0\\
					&&& {g_2}(\mathbf{x};\mathbf{p}) = {T_m} - T_{\textrm{deposit}} \le 0\\
				\end{aligned}
			\end{equation*}
		\end{alertblock}
		% Parameters
		\vspace{-1em}
		\begin{exampleblock}{Changing requirements}
			\vspace{-1.5em}
			\begin{equation*}
				\begin{aligned}
					& & & W_{\textrm{total}} = \textrm{width of base part}\\
					& & & T_{m} = \textrm{Melting temperature of deposit}\\
					& & & P_{\textrm{load}} = \textrm{Pressure load on outer casing}\\
				\end{aligned}
			\end{equation*}
		\end{exampleblock}
		
	\column{.6\textwidth} % Right column and width
		Let us consider one variable %({\color{red}$x_3$}) 
  and two parameters 
  %({\color{red}$p_2$} and {\color{red}$p_3$})
  \\
		For an optimal design to be scalable, the product of 
		\begin{itemize}
			\item<2-> the \emphasis{direction of change} 
			\item<5-> the \emphasis{change effect}, and 
			\item<6-> the variable \emphasis{monotonicity} 
		\end{itemize}
 \uncover<6->{must be greater or equal to zero}
		\uncover<13->{~\\Our full example: {\color{red}$4$} variables and {\color{red}$3$} parameters}

		\newcommand{\xshift}{5.6cm}
		\newcommand{\yshift}{-2cm}

		\begin{tikzpicture}[remember picture,overlay]
				
			\only<6>{
				\node (textscalable) at ( [xshift = \xshift, yshift = -0.1cm] current page.center)
				{\textbf{Scalable design!}};
			}%
			\only<12>{
				\node (textscalable) at ( [xshift = \xshift, yshift = -0.1cm] current page.center)
				{\textbf{{\color{red} Not} a Scalable design!}};
			}%

			\only<1>{
				\fixedpicture{\xshift}{\yshift}{3.6cm}{scalability/scalability_example/11_gradient_demo_1.pdf}
			}% empty contours
			\only<2>{
				\fixedpicture{\xshift}{\yshift}{3.6cm}{scalability/scalability_example/11_gradient_demo_2.pdf}
			}% gradient vector only
			\only<3>{
				\fixedpicture{\xshift}{\yshift}{3.6cm}{scalability/scalability_example/11_gradient_demo_3.pdf}
			}% gradient vector + component A
			\only<4>{
				\fixedpicture{\xshift}{\yshift}{3.6cm}{scalability/scalability_example/11_gradient_demo_4.pdf}
			}% gradient vector + component A + B
			\only<5-6>{
				\fixedpicture{\xshift}{\yshift}{3.6cm}{scalability/scalability_example/11_gradient_demo_5.pdf}
			}% gradient vector + component A + B + change quadrant (scalable)
			\only<7>{
				\fixedpicture{\xshift}{\yshift}{3.6cm}{scalability/scalability_example/11_gradient_demo_6.pdf}
			}% empty contours
			\only<8>{
				\fixedpicture{\xshift}{\yshift}{3.6cm}{scalability/scalability_example/11_gradient_demo_7.pdf}
			}% gradient vector only
			\only<9>{
				\fixedpicture{\xshift}{\yshift}{3.6cm}{scalability/scalability_example/11_gradient_demo_8.pdf}
			}% gradient vector + component A
			\only<10>{
				\fixedpicture{\xshift}{\yshift}{3.6cm}{scalability/scalability_example/11_gradient_demo_9.pdf}
			}% gradient vector + component A + B
			\only<11-12>{
				\fixedpicture{\xshift}{\yshift}{3.6cm}{scalability/scalability_example/11_gradient_demo_10.pdf}
			}% gradient vector + component A + B + change quadrant (non scalable)
		\end{tikzpicture}


	\end{columns}
	\only<-12>{~\\~\\}%
	$
		\small
		\only<-12>{
			\uncover<5->{
				\overbrace{
					\begin{bmatrix}
						\color<3->{red}-1 & 0\\ 
						0 & \color<3->{red}-1
					\end{bmatrix}
				}^{\mathbf{N}=\textrm{Change effect}}
				\times
			}%
		}% example matrix
		\only<13->{
			\overbrace{
				\begin{bmatrix}
					\color{red}n_1 &  0 & 0\\ 
					0 & \color{red}n_2 & 0\\ 
					0 &  0 & \color{red}n_3
				\end{bmatrix}
			}^{\mathbf{N}=\textrm{Change effect}}
			\times
		} % generalized matrix
		%---------------------------------
		\only<-12>{
			\uncover<2->{
				\overbrace{
					\begin{bmatrix}
						\uncover<3->{{^{\partial x_3^*}/_{\partial p_2}}} \\ 
						\uncover<4->{{^{\partial x_3^*}/_{\partial p_3}}}
					\end{bmatrix}
				}^{\mathbf{J}=\textrm{Direction of change}}
			}
		}% example matrix
		\only<13->{
			\uncover<2->{
				\overbrace{
					\begin{bmatrix}
						{^{\partial x_1^*}/_{\partial p_1}} & {^{\partial x_2^*}/_{\partial p_1}} & {^{\partial x_3^*}/_{\partial p_1}} & {^{\partial x_4^*}/_{\partial p_1}} \\ 
						{^{\partial x_1^*}/_{\partial p_2}} & {^{\partial x_2^*}/_{\partial p_2}} & {^{\partial x_3^*}/_{\partial p_2}} & {^{\partial x_4^*}/_{\partial p_2}} \\ 
						{^{\partial x_1^*}/_{\partial p_3}} & {^{\partial x_2^*}/_{\partial p_3}} & {^{\partial x_3^*}/_{\partial p_3}} & {^{\partial x_4^*}/_{\partial p_3}}
					\end{bmatrix}
				}^{\mathbf{J}=\textrm{Direction of change}}
			}
		}% generalized matrix
		%---------------------------------
		\only<-12>{
			\uncover<6->{
				\times
				\overbrace{
					\begin{bmatrix}
						\color<4->{red}1
					\end{bmatrix} 
				}^{\mathbf{M}=\textrm{Monotonicity}}
				\ge \mathbf{0}
			}%
		}% example matrix
		\only<13->{
			\times
			\overbrace{
				\begin{bmatrix}
					\color{red}m_1 & 0 & 0 & 0 \\ 
					0 & \color{red}m_2 & 0 & 0 \\ 
					0 & 0 & \color{red}m_3 & 0 \\ 
					0 & 0 & 0 & \color{red}m_4
				\end{bmatrix} 
			}^{\mathbf{M}=\textrm{Monotonicity}}
			\ge \mathbf{0}
		}% generalized matrix
	$
\end{frame}

%%%%%%%%%%%%%%%%%%%%%%%%%%%%%%%%%%%%%%%%%%%%%%%%%%%%%%%
%%               Scalability assessment              %%
%%%%%%%%%%%%%%%%%%%%%%%%%%%%%%%%%%%%%%%%%%%%%%%%%%%%%%%
%------------------------------------------------
\section{Practical scalability results}
% \frame{\tableofcontents[
% 	currentsection, 
% 	hideallsubsections
% 	]}
%------------------------------------------------
\subsection{Response surface of parameter space to obtain scalable optimal designs}
%------------------------------------------------
\begin{frame}[t]
\frametitle{Response surface of parametric optimal solutions and scalable designs}
\begin{itemize}
	\item Computing the components of $\mathbf{J}(\mathbf{p})$ can be prohibitively expensive 
 %in engineering design
	\item<2-> Response surface $\hat{\mathbf{x}^*}(\mathbf{p})$ can be obtained from parametric optimal designs $\mathbf{X}^*= \left\{\mathbf{x}^*(\mathbf{p}_1)~\mathbf{x}^*(\mathbf{p}_2)~\cdots~\mathbf{x}^*(\mathbf{p}_k)\right\}$ \\
	(we use kernel smoothing)
 %is used for the response surface
	\item<3-> Jacobian $\mathbf{J}(\mathbf{p})$ of $\hat{\mathbf{x}^*}(\mathbf{p})$ estimated from the derivatives of the kernel basis functions
	\item<4-> Scalable optimal designs given by $\mathbf{N}\mathbf{J}(\mathbf{p})^{\mathrm T}\mathbf{M} \ge \mathbf{0}$. %This is the \emphasis{set of designs} we are looking for.
\end{itemize}

\begin{figure}
	\only<2-3>{\includegraphics[width=0.9\textwidth]{scalability/parameter_space/8_PS_f_1.pdf}}%
	\only<4->{\includegraphics[width=0.9\textwidth]{scalability/parameter_space/8_PS_f_3.pdf}}%
	%\caption{\only<2->{2D projections of the parameter space  \only<5>{with {\color{red}sescalable regions} overlayed}}
\end{figure}

\end{frame}

%NEW
%------------------------------------------------
\subsection{Scalable optimal designs in the design space}
%------------------------------------------------
\begin{frame}[t]
\frametitle{Map scalable optimal design set to design space}
% \only<1>{
% 	Scalable optimal designs mapped back to design space using response surface
% }%
% \only<2->{This is the result we are looking for! \only<4>{\color{red}A manageable set of scalable solutions}}%
% \only<1>{
% \begin{tikzpicture}[remember picture,overlay]
% 	\node (SBDresult) at ([xshift=0.5\textwidth, yshift =-0.05\textwidth]current page.west) %or: (current page.center)
% 		{\includegraphics[width=0.9\textwidth]{scalability/SBD_results/6_SBD_results_f_9.pdf}};
% \end{tikzpicture}
% }%
% \only<2>{
\begin{tikzpicture}[remember picture,overlay]
	\node (SBDresult) at ([xshift=0.5\textwidth, yshift =-0.05\textwidth]current page.west) %or: (current page.center)
		{\includegraphics[width=0.9\textwidth]{scalability/SBD_results/6_SBD_results_f_10.pdf} };
\end{tikzpicture}
%}%
% \only<3>{
% \begin{tikzpicture}[remember picture,overlay]
% 	\node (SBDresult) at ([xshift=0.5\textwidth, yshift =-0.05\textwidth]current page.west) %or: (current page.center)
% 		{\includegraphics[width=0.9\textwidth]{scalability/SBD_results/6_SBD_results_f_11.pdf}};
% \end{tikzpicture}
% }%
% \only<4>{
% \begin{tikzpicture}[remember picture,overlay]
% 	\node (SBDresult) at ([xshift=0.5\textwidth, yshift =-0.05\textwidth]current page.west) %or: (current page.center)
% 		{\includegraphics[width=0.9\textwidth]{scalability/SBD_results/6_SBD_results_f_12.pdf}};
% \end{tikzpicture}
% }%

\end{frame}



% Section 02: design margins
%%%%%%%%%%%%%%%%%%%%%%%%%%%%%%%%%%%%%%%%%%%%%%%%%%%%%%%
%%        ANSWERING THE RESEARCH QUESTIONS           %%
%%%%%%%%%%%%%%%%%%%%%%%%%%%%%%%%%%%%%%%%%%%%%%%%%%%%%%%
\section{Change absorption through excess and buffer}
%------------------------------------------------
\begin{frame}[c,noframenumbering]
	\centering
	% \setlength\fboxsep{0pt}
	\begin{titleblock}{}
		~\\%
		{\centering\LARGE Using design margins to absorb change in requirements \ifshowcitations\footnotemark[1]\fi\\}%
		~\\%
	\end{titleblock}
	{\color{white}\ifshowcitations\footpartcite{Alhandawi2021}[1]\fi}
\end{frame}
\addtocounter{footnote}{-1}
%------------------------------------------------
\begin{frame}[t,label=ap2_2]
	\frametitle{Quantifying capability, buffer, and excess}
	\vspace{-1em}
	\uncover<2->{We focus on one 2D projection}%
	~\\
	\begin{columns}[t] % The "c" option specifies centered vertical alignment while the "t" option is used for top vertical alignment
		\begin{column}{.5\textwidth} % Left column and width
			\small
			\vspace{-1em}
			\begin{itemize}
				\item<1-> \textit{Failure criterion}: Factor of safety $n_{\textrm{safety}}(\mathbf{T})$ \uncover<2->{represented using isocontours}
				\item<3-> \textit{Capability}: Loads that \emphasis{can} be satisfied
				\item<4-> \textit{Requirements}: Loads that \emphasis{must} be satisfied
				\item<5-> \textit{Buffer}: Portion of requirement \emphasis{satisfied}
				\item<6-> \textit{Danger}: Portion of requirement \emphasis{not satisfied}
				\item<7-> \textit{Excess}: \emphasis{Unused} capability not needed for satisfying requirement
				\item<8-> \textit{Reliability}: $\smash{\mathbb{P}(\mathbf{T} \in C)}$ estimated using Monte-Carlo integration
			\end{itemize}
		\end{column}
	
		\begin{column}{.45\textwidth} % Left column and width
			\begin{figure}
				\only<1>{\includegraphics[height=5.0cm]{design_margins/DM_geometry/TRS_isometric_failure_domain.pdf}}%
				\only<2>{\includegraphics[height=5.6cm]{design_margins/DM_parameter_space/109_req_1_thermal_out_nominal_RS_2D_pi_0.pdf}}%
				\only<3>{\includegraphics[height=5.6cm]{design_margins/DM_parameter_space/109_req_1_thermal_out_nominal_RS_2D_pi_1.pdf}}%
				\only<4>{\includegraphics[height=5.6cm]{design_margins/DM_parameter_space/109_req_1_thermal_out_nominal_RS_2D_pi_2.pdf}}%
				\only<5>{\includegraphics[height=5.6cm]{design_margins/DM_parameter_space/109_req_1_thermal_out_nominal_RS_2D_pi_3.pdf}}%
				\only<6>{\includegraphics[height=5.6cm]{design_margins/DM_parameter_space/109_req_1_thermal_out_nominal_RS_2D_pi_4.pdf}}%
				\only<7>{\includegraphics[height=5.6cm]{design_margins/DM_parameter_space/109_req_1_thermal_out_nominal_RS_2D_pi_5.pdf}}%
				\only<8>{\includegraphics[height=5.6cm]{design_margins/DM_parameter_space/109_req_1_thermal_out_nominal_RS_2D_pi_6.pdf}}%
				\only<1>{
					\vspace{-0.2em}
					\caption{Failure domain}
				}%
			\end{figure}
		\end{column}
		
	\end{columns}
	
\end{frame}
%------------------------------------------------
\begin{frame}[t,label=ap2_3]
	\frametitle{Considering redesign when buffer is insufficient}
	\vspace{-1em}
	\only<1>{Going back to our baseline design we can calculate these metrics as follows:}%
	\only<2>{Consider depositing more material}%
	\only<3>{Reliability increases}%
	\only<4->{However, {\color{red} requirements can change as well!}}%
	~\\
	\begin{columns}[t] % The "c" option specifies centered vertical alignment while the "t" option is used for top vertical alignment
		\begin{column}{.45\textwidth} % Left column and width
			\only<1>{$c=1$,$\mathbf{D}=\left[1,2,4\right]$}%
			\only<2->{$c=1$,$\mathbf{D}=\left[1,2,4,{\color{red}0}\right]$}%
			\begin{figure}
				\only<1-2>{\includegraphics[height=5.0cm]{design_margins/DM_geometry/TRS_isometric_7.pdf}}%
				\only<3->{\includegraphics[height=5.0cm]{design_margins/DM_geometry/TRS_isometric_8.pdf}}%
				\vspace{-0.2em}
				\caption{Turbine Rear Structure (TRS)}
			\end{figure}
			\vspace{-2em}
			\begin{table}[h!]
				\centering
				% \renewcommand{\arraystretch}{1.0}% Wider
				\footnotesize\addtolength{\tabcolsep}{-2pt}
				\begin{tabular}{ccc}
					\bf Reliability & \bf Excess & \bf Weight\\
					\only<1-2>{0.3089 & 0.529 & 13.9\\}%
					\only<3>{0.759({\color{darkgreen}$\uparrow 145\%$}) & 0.856({\color{red}$\uparrow 62\%$}) & 18.5($\color{red}\uparrow 33\%$)\\}
					\only<4>{0.943({\color{darkgreen}$\uparrow 24\%$}) & 0.754({\color{darkgreen}$\downarrow 12\%$}) & 18.5\\}
					\only<5>{1.00({\color{darkgreen}$\uparrow 6\%$}) & 0.871({\color{red}$\uparrow 15\%$}) & 18.5\\}
				\end{tabular}
			\end{table}
		\end{column}
		
		\begin{column}{.45\textwidth} % Left column and width
			\centering
			~\\%
			~\\%
			\vspace{-1em}
			\begin{figure}
				\only<1-2>{\includegraphics[height=5.6cm]{design_margins/DM_parameter_space/109_req_1_thermal_out_nominal_RS_2D_pi_6.pdf}}%
				\only<3>{\includegraphics[height=5.6cm]{design_margins/DM_parameter_space/110_req_1_thermal_out_nominal_RS_2D_pi_6.pdf}}%
				\only<4>{\includegraphics[height=5.6cm]{design_margins/DM_parameter_space/110_reqs_2D/110_req_15_thermal_out_RS_2D_pi_6.pdf}}%
				\only<5>{\includegraphics[height=5.6cm]{design_margins/DM_parameter_space/110_reqs_2D/110_req_11_thermal_out_RS_2D_pi_6.pdf}}%
				\label{fig:2DPspace}
			\end{figure}
		\end{column}
		
	\end{columns}
	
\end{frame}

%%%%%%%%%%%%%%%%%%%%%%%%%%%%%%%%%%%%%%%%%%%%%%%%%%%%%%%
%%         APPLYING THE RESEARCH QUESTIONS           %%
%%%%%%%%%%%%%%%%%%%%%%%%%%%%%%%%%%%%%%%%%%%%%%%%%%%%%%%

%------------------------------------------------
\subsection{Optimizing when and how much change is needed}
%------------------------------------------------
\begin{frame}[t]
	\frametitle{Optimizing when and how much change is needed}
	\vspace{-0.5em}
	Formulate optimization problem to minimize total excess while being reliable
	
	\begin{columns}[t] % The "c" option specifies centered vertical alignment while the "t" option is used for top vertical alignment
		
		\begin{column}{.42\textwidth} % Left column and width
			\vspace{-1.5em}
			% Optimization problem
			\begin{alertblock}{Objective and constraints}
				\vspace{-1.5em}
				\begin{equation*}
					\begin{aligned}
						& {\text{minimize}}
						& & \text{excess}(c,\mathbf{D};\text{requirement arc})\\
						& \text{subject to}
						& & reliability(c,\mathbf{D};\text{requirement arc})\\
						& \text{over}
						& & \text{The set of feasible design decisions}\\
						& & & \text{(404 possible combinations)}
					\end{aligned}
				\end{equation*}
			\end{alertblock}
			\vspace{-0.5em}
			\uncover<3->{
				% Variables
				\begin{exampleblock}{Design variables}
					\vspace{-0.5em}
					\begin{equation*}
						\begin{aligned}
							& & & c \in \left\{"wavy","hatched","tubular"\right\}\\
							& & & \mathbf{D} = \left[\text{decision}_1,\text{decision}_2,\cdots\right]\\
							& & & \text{decision} \in \left\{-1,0,1,2,3,4\right\}\\
						\end{aligned}
					\end{equation*}
				\end{exampleblock}
			}
			\vspace{-0.5em}
		\end{column}

		\begin{column}{.5\textwidth} % Right column and width
			\vspace{-1em}
			\tikzstyle{background grid}=[draw, black!50,step=.5cm]
			
			\begin{tikzpicture}%[show background grid]
			% Put the graphic inside a node. This makes it easy to place the
			% graphic and to draw on top of it. 
			% The above right option is used to place the lower left corner
			% of the image at the (0,0) coordinate. 
				\uncover<-5>{
					\node [inner sep=0pt,above right, opacity=0.6]  at (0, 1) (objective)
						{
							\only<2->{\includegraphics[height=5.0cm]{design_margins/DM_parameter_space/110_reqs_2D/110_req_11_thermal_out_RS_2D_pi_6.pdf}}%
						};
					\node [inner sep=0pt,above right, opacity=0.7, fill=white]  at (4.0, -.75) (constraint)
						{
							\only<3->{\includegraphics[height=4.0cm]{design_margins/DM_geometry/TRS_isometric_8.pdf}}%
						};
				}
				% show origin
				% \fill (0,0) circle (2pt);
				% define destination coordinates
			\end{tikzpicture}

		\end{column}

	\end{columns}
\end{frame}

% Section 03: Design space exploration
%==============================================================================%
%% Mathematical symbols
\definecolor{design}{HTML}{D0A31D}
\definecolor{spec}{HTML}{FF7575}
\definecolor{margin}{HTML}{CC99FF}
\definecolor{perf}{HTML}{87C061}
\newcommand\design[1]{{\color{design}#1}}
\newcommand\spec[1]{{\color{spec}#1}}
\newcommand\margin[1]{{\color{margin}#1}}
\newcommand\perf[1]{{\color{perf}#1}}
%==============================================================================%
%------------------------------------------------
\section{Ongoing work}
%------------------------------------------------
\begin{frame}[c,noframenumbering]
	\centering
	% \setlength\fboxsep{0pt}
	\begin{titleblock}{}
		~\\%
		{\centering\LARGE Ongoing work: Design space exploration\\}%
		~\\%
	\end{titleblock}
\end{frame}
%------------------------------------------------
\subsection{Introduction on DSE}
%------------------------------------------------
\begin{frame}
	\frametitle{Design space exploration at the conceptual level}
	
	\only<1>{Consider a functional model (FM) of a product (conceptual design)}%
	\only<2>{\emph{hyperlinks}\textsuperscript{1} used in multi-domain modelling}%
	\only<3>{to link FM $\to$ margin value analysis\textsuperscript{2}}%
	\only<4->{different concepts have different preliminary designs}%
	{\color{white}\ifshowcitations\footpartcite{Ahmad2013}\textsuperscript{,}\footpartcite{Brahma2020}\fi}%
	
	\begin{columns}[t] % The "c" option specifies centered vertical alignment while the "t" option is used for top vertical alignment
			
		\begin{column}{.5\textwidth} % Left column and width
	
			\centering\textbf{conceptual design}
			\begin{figure}
				\only<1>{\includegraphics[height=5cm]{design_space_exploration/Intro/efm_example_1.pdf}}%
				\only<2>{\includegraphics[height=5cm]{design_space_exploration/Intro/efm_example_2.pdf}}%
				\only<3-5>{\includegraphics[height=5cm]{design_space_exploration/Intro/efm_example_3.pdf}}%
				\only<6->{\includegraphics[height=5cm]{design_space_exploration/Intro/efm_example_4.pdf}}%
			\end{figure}
			%
			\tikzstyle{background grid}=[draw, black!50,step=.5cm]
			\begin{tikzpicture}[remember picture, overlay, shift={(2,6)}] %show background grid, 
				\node [draw=black, line width=0.5mm, inner sep=0pt,above left, opacity=1.0]  at (-0.35\textwidth,-0.45\textheight) (trsA) 
					{
						\only<3->{\includegraphics[width=0.15\textwidth]{design_space_exploration/Intro/TRS_alt_1.pdf}}%
					};
				\node [draw=black, line width=0.5mm, inner sep=0pt,above left, opacity=1.0]  at (-0.15\textwidth,-0.45\textheight) (trsB) 
					{
						\only<6->{\includegraphics[width=0.15\textwidth]{design_space_exploration/Intro/TRS_alt_3.pdf}}%
					};
				% show origin
				% \fill (0,0) circle (2pt);
			\end{tikzpicture}%
			%
		\end{column}
	
		\vrule{}
	
		\begin{column}{.5\textwidth} % Left column and width
	
			\centering\textbf{preliminary design}
	
			\begin{figure}
				\uncover<2->{
					\only<-2>{\includegraphics[height=5cm]{design_space_exploration/Intro/man_example_1.pdf}}%
				}%
				\only<3-5>{\includegraphics[height=5cm]{design_space_exploration/Intro/man_example_2.pdf}}%
				\only<6->{\includegraphics[height=5cm]{design_space_exploration/Intro/man_example_3.pdf}}%
			\end{figure}
			%
			\tikzstyle{background grid}=[draw, black!50,step=.5cm]
			\begin{tikzpicture}[remember picture, overlay, shift={(2,6)}] %show background grid, 
				\node [draw=black, line width=0.1mm, fill=white!20, inner sep=0pt,above left, opacity=1.0]  at (-0.4\textwidth,-0.65\textheight) (loadA) 
					{
						\only<5->{\includegraphics[width=0.35\textwidth]{design_space_exploration/Intro/load_cases_1.pdf}}%
					};
				\node [draw=black, line width=0.1mm, fill=white!20, inner sep=0pt,above left, opacity=1.0]  at (-0.4\textwidth,-0.15\textheight) (loadB) 
					{
						\only<7->{\includegraphics[width=0.35\textwidth]{design_space_exploration/Intro/load_cases_2.pdf}}%
					};
				\node [draw=black, line width=0.1mm, fill=white!20, inner sep=1pt,above left, opacity=1.0]  at (0.22\textwidth,-0.65\textheight) (manuA) 
					{
						\only<4->{\includegraphics[width=0.35\textwidth]{design_space_exploration/TEMP_GIF/frame_180.png}}%
					};
				% show origin
				% \fill (0,0) circle (2pt);

				\path (-3,-5) coordinate (B1 A)
					  (-4.5,-8) coordinate (B3 A);

			\end{tikzpicture}%
			%
		\end{column}
	\end{columns}
	
	\vspace{-2em}
\end{frame}
\addtocounter{footnote}{-2}
%------------------------------------------------
\begin{frame}
	\frametitle{Design space exploration at the conceptual level}
	
	Multi-domain modelling (MDM) matrix for change propagation management can be used to represent hyperlinks \footnotemark[1]
	\begin{figure}
		\includegraphics[height=5cm]{design_space_exploration/Intro/MDM_example.pdf}
	\end{figure}
	
	\vspace{-2em}
	{\color{white}
		\footpartcite{Koh2012}[1]
		\addtocounter{footnote}{-1}
	}
	
\end{frame}
%------------------------------------------------
\subsection{Example problem: Design of TRS strut}
%------------------------------------------------
\begin{frame}[t,label=ap2_2]
	\frametitle{Overview of turbine rear frame design}
	\vspace{-1em}
	\only<1-2>{The turbine rear structure (TRS) sits aft of the turbine}%
	\only<3>{We consider the design of a single strut}%
	\only<4>{Variable thermal loads are expected}%
	\only<5-7>{Which cause thermal expansion \uncover<6-7>{and stress on the strut}}%
	\only<8>{The entire problem can be represented compactly as an MDM}%
	~\\
	\begin{columns}[t] % The "c" option specifies centered vertical alignment while the "t" option is used for top vertical alignment
			
		\begin{column}{.45\textwidth} % Left column and width
			\footnotesize
			Design parameters:
			\design{
				\uncover<3->{
				\begin{itemize}\itemsep0em 
					\item \only<8>{$D1$} vane width $w$
					\item \only<8>{$D2$} vane height $h$
					\item \only<8>{$D3$} lean angle $\theta$
					\item \only<8>{$D4$} Material - yield stress $\sigma_y$ (off-the shelf)
				\end{itemize}
			}}%
			Input Specifications:
			\spec{
				\uncover<5->{
				\begin{itemize}\itemsep0em 
					\item \only<8>{$S1$} Nacelle temperature $T_1$ ($\downarrow$)
					\item \only<8>{$S2$} Gas surface temperature $T_2$ ($\uparrow$)
				\end{itemize}
			}}%
			Design margins:
			\margin{
				\uncover<6->{
				\begin{itemize}\itemsep0em 
					\item \only<8>{$E1$} compression $\sigma_y - \sigma_a \ge 0$
					\item \only<8>{$E2$} bending $\sigma_y - \sigma_m \ge 0$
					\item \only<8>{$E3$} buckling $F_\text{buckling} - F_a \ge 0$
				\end{itemize}
			}}%
			Performance parameters:
			\perf{
				\uncover<7->{
				\begin{itemize}\itemsep0em 
					\item \only<8>{$P1$}  Weight
				\end{itemize}
			}}%
		\end{column}
	
		\begin{column}{.4\textwidth} % Left column and width
			\begin{figure}
				\only<1>{
					\includegraphics[width=6.0cm]{design_space_exploration/TRS/engine.png}
					\vspace{-0.2em}
					\caption{Trent 900 aeroengine}
				}%
				\only<2>{
					\includegraphics[width=5.0cm]{design_space_exploration/TRS/rear_view.pdf}
					\vspace{-0.2em}
					\caption{Turbine rear structure}
				}%
				\only<3>{
					\includegraphics[width=7.0cm]{design_space_exploration/TRS/cross_section_view.pdf}
					\vspace{-0.2em}
					\caption{Turbine rear structure: cross-sectional view}
				}%
				\only<4>{
					\includegraphics[width=6.0cm]{design_space_exploration/TRS/expansion_1.pdf}
					\vspace{-0.2em}
					\caption{Thermal expansion}
				}%
				\only<5>{
					\includegraphics[width=6.0cm]{design_space_exploration/TRS/expansion_2.pdf}
					\vspace{-0.2em}
					\caption{Thermal expansion}
				}%
				\only<6-7>{
					\includegraphics[width=7.0cm]{design_space_exploration/TRS/loads.pdf}
					\vspace{-0.2em}
					\caption{Load cases}
				}%
				\only<8>{
					\vspace{-1.2em}
					\includegraphics[width=6.7cm]{design_space_exploration/TRS/MDM_strut.pdf}
					\vspace{-0.2em}
					\caption{Mulit-domain matrix of strut example}
				}%
			\end{figure}
		\end{column}
		
	\end{columns}
	
\end{frame}
%------------------------------------------------
\subsection{Applying the margin value method}
%------------------------------------------------
\subsubsection{Calculating impact on performance}
%------------------------------------------------
\begin{frame}[t,label=ap2_2]
	\frametitle{Calculating impact of excess margin on performance}
	\vspace{-1em}
	Need to construct a response surface%
	~\\
	\begin{columns}[t] % The "c" option specifies centered vertical alignment while the "t" option is used for top vertical alignment
			
		\begin{column}{.45\textwidth} % Left column and width
			\footnotesize
			Design parameters:
			\design{
				\begin{itemize}\itemsep0em 
					\item vane width $w$
					\item vane height $h$
					\item lean angle $\theta$
					\item Material - yield stress $\sigma_y$ (off-the shelf)
				\end{itemize}
			}%
			Input Specifications:
			\spec{
				\begin{itemize}\itemsep0em 
					\item Nacelle temperature $T_1$ ($\downarrow$)
					\item Gas surface temperature $T_2$ ($\uparrow$)
				\end{itemize}
			}%
			Design margins:
			\margin{
				\begin{itemize}\itemsep0em 
					\item compression $\sigma_y - \sigma_a \ge 0$
					\item bending $\sigma_y - \sigma_m \ge 0$
					\item buckling $F_\text{buckling} - F_a \ge 0$
				\end{itemize}
			}%
			Performance parameters:
			\perf{
				\begin{itemize}\itemsep0em 
					\item Weight
				\end{itemize}
			}%
		\end{column}
	
		\begin{column}{.4\textwidth} % Left column and width
			\footnotesize
			Need to find a function $f$ that translates \margin{excess} to \perf{performance}
			\renewcommand\baselinestretch{0.5}
			\only<1>{
				\small
				\begin{equation*}
					\begin{aligned}
						& \hat{\mathbf{p}} = f(\mathbf{e})\\
						& \mathbf{p} = \left[p_1,p_2,\cdots,p_j\right]^{\mathit{T}}\\
						& \mathbf{e} = \left[e_1,e_2,\cdots,e_m\right]^{\mathit{T}}
					\end{aligned}
				\end{equation*}
			}%
			\only<2>{
				\small
				\begin{equation*}
					\begin{aligned}
						& \hat{\mathbf{p}} = f(\mathbf{e})\\
						& \mathbf{p}_1^\mathrm{threshold} = \left[p_1,p_2,\cdots,p_j\right]^{\mathit{T}}\\
						& \mathbf{e} = \left[{\color{red}0},e_2,\cdots,e_m\right]^{\mathit{T}}
					\end{aligned}
				\end{equation*}
			}%
			\only<3>{
				\small
				\begin{equation*}
					\begin{aligned}
						& \hat{\mathbf{p}} = f(\mathbf{e})\\
						& \mathbf{p}_2^\mathrm{threshold} = \left[p_1,p_2,\cdots,p_j\right]^{\mathit{T}}\\
						& \mathbf{e} = \left[e_1,{\color{red}0},\cdots,e_m\right]^{\mathit{T}}
					\end{aligned}
				\end{equation*}
			}%
			\only<4>{
				\renewcommand{\kbldelim}{[}% Left delimiter
				\renewcommand{\kbrdelim}{]}% Right delimiter
				\[
				  \text{Impact} = \kbordermatrix{
										& \text{weight} \\
					\text{buckling} 	& 0.545  	\\
					\text{compression} 	& 0.111 	\\
					\text{bending} 		& -1.011
				  }
				\]
			}%

			\vspace{-1em}
			\begin{figure}
				\only<1>{
					\includegraphics[width=5.0cm]{design_space_exploration/MVM/IOP_1.pdf}
				}%
				\only<2>{
					\includegraphics[width=5.0cm]{design_space_exploration/MVM/IOP_2.pdf}
				}%
				\only<3>{
					\includegraphics[width=5.0cm]{design_space_exploration/MVM/IOP_3.pdf}
				}%
			\end{figure}
		\end{column}
		
	\end{columns}
	
\end{frame}
%------------------------------------------------
\subsubsection{Calculating change absorption capability}
%------------------------------------------------
\begin{frame}[t,label=ap2_2]
	\frametitle{Calculating deterioration (allowable change in specifications)}
	\vspace{-1em}
	We calculate the maximum allowable deterioration in the \spec{specifications}%
	~\\
	\begin{columns}[t] % The "c" option specifies centered vertical alignment while the "t" option is used for top vertical alignment
			
		\begin{column}{.45\textwidth} % Left column and width
			\footnotesize
			Design parameters:
			\design{
				\begin{itemize}\itemsep0em 
					\item vane width $w$
					\item vane height $h$
					\item lean angle $\theta$
					\item Material - yield stress $\sigma_y$ (off-the shelf)
				\end{itemize}
			}%
			Input Specifications:
			\spec{
				\begin{itemize}\itemsep0em 
					\item Nacelle temperature $T_1$ ($\uparrow$)
					\item Gas surface temperature $T_2$ ($\downarrow$)
				\end{itemize}
			}%
			Design margins:
			\margin{
				\begin{itemize}\itemsep0em 
					\item compression $\sigma_y - \sigma_a \ge 0$
					\item bending $\sigma_y - \sigma_m \ge 0$
					\item buckling $F_\text{buckling} - F_a \ge 0$
				\end{itemize}
			}%
			Performance parameters:
			\perf{
				\begin{itemize}\itemsep0em 
					\item Weight
				\end{itemize}
			}%
		\end{column}
	
		\begin{column}{.4\textwidth} % Left column and width
			Gradually change each input specifications until any margin is consumed
			\renewcommand\baselinestretch{0.5}
			\renewcommand{\kbldelim}{[}% Left delimiter
			\renewcommand{\kbrdelim}{]}% Right delimiter
			\only<1>{
				\[
				  \mathbf{s} = \kbordermatrix{
							&		\\
					T_1 	& 450  	\\
					T_2 	& 425 	\\
				  }
				\]
			}%
			\only<2>{
				\[
				  \mathbf{s} = \kbordermatrix{
					  						&			\\
					\color{red}T_1\uparrow 	& 454.5   	\\
					T_2						& 425 		\\
				  }
				\]
			}%
			\only<3>{
				\[
				  \mathbf{s} = \kbordermatrix{
					  							&			\\
					T_1 						& 450   	\\
					\color{red}T_2\downarrow	& 420.75 	\\
				  }
				\]
			}%
			\only<4>{
				\[
				  \mathbf{s}^\text{max} = \kbordermatrix{
					  				&			\\
					T_1\uparrow 	& 454.5   	\\
					T_2\downarrow	& 420.75 	\\
				  }
				\]
			}%
			\only<5>{
				\[
				  \text{deterioration} = \kbordermatrix{
					  				&			\\
					T_1				& 0.01  	\\
					T_2				& 0.01 		\\
				  }
				\]
			}%
			\only<1>{
				\[
				  \mathbf{e} = \kbordermatrix{
										&				\\
					\text{buckling} 	& 597.440		\\
					\text{compression} 	& 72.6904696	\\
					\text{bending} 		& 451.9918707	\\
				  }
				\]
			}%
			\only<2-3>{
				\[
				  \mathbf{e} = \kbordermatrix{
										&				\\
					\text{buckling} 	& 488.404		\\
					\text{compression} 	& \color{red}0	\\
					\text{bending} 		& 450.4889005	\\
				  }
				\]
			}%
		\end{column}
		
	\end{columns}
	
\end{frame}
%------------------------------------------------
\begin{frame}[t,label=ap2_2]
	\frametitle{Calculating deterioration (allowable change in specifications)}
	\vspace{-1em}
	\only<1>{We calculate the change absorption capability as a function of deterioration}%
	~\\
	\begin{columns}[t] % The "c" option specifies centered vertical alignment while the "t" option is used for top vertical alignment
			
		\begin{column}{.45\textwidth} % Left column and width
			\footnotesize
			Design parameters:
			\design{
				\begin{itemize}\itemsep0em 
					\item vane width $w$
					\item vane height $h$
					\item lean angle $\theta$
					\item Material - yield stress $\sigma_y$ (off-the shelf)
				\end{itemize}
			}%
			Input Specifications:
			\spec{
				\begin{itemize}\itemsep0em 
					\item Nacelle temperature $T_1$ ($\downarrow$)
					\item Gas surface temperature $T_2$ ($\uparrow$)
				\end{itemize}
			}%
			Design margins:
			\margin{
				\begin{itemize}\itemsep0em 
					\item compression $\sigma_y - \sigma_a \ge 0$
					\item bending $\sigma_y - \sigma_m \ge 0$
					\item buckling $F_\text{buckling} - F_a \ge 0$
				\end{itemize}
			}%
			Performance parameters:
			\perf{
				\begin{itemize}\itemsep0em 
					\item Weight
				\end{itemize}
			}%
		\end{column}
	
		\begin{column}{.4\textwidth} % Left column and width
			Change absorption is calculated using the following equation:
			\renewcommand\baselinestretch{0.5}
			\begin{equation*}
				\text{absorption}_{mi} = \dfrac{t^\text{new}_{mi} - t^\text{nominal}_{mi}}{t^\text{nominal}_{mi}\times\text{deterioration}_i}
			\end{equation*}
			for each margin node $m$ and specification $i$
			\renewcommand{\kbldelim}{[}% Left delimiter
			\renewcommand{\kbrdelim}{]}% Right delimiter
			\only<1>{
				\[
				  \text{absorption} = \kbordermatrix{
										& T_1	& T_2	\\
					\text{buckling} 	&2.79	& 3.87	\\
					\text{compression} 	&1.98	& 3.04	\\
					\text{bending} 		&2.58	& 3.66	\\
				  }
				\]
			}%
		\end{column}
		
	\end{columns}
	
\end{frame}
%------------------------------------------------
\subsubsection{Calculating change absorption capability}
%------------------------------------------------
\begin{frame}[t,label=ap2_2]
	\frametitle{Constructing the margin value plot (MVP)}
	\vspace{-1em}
	{Using ``impact'' and ``absorption'' we can construct a tradespace (MVP)}%
	~\\
	\begin{columns}[t] % The "c" option specifies centered vertical alignment while the "t" option is used for top vertical alignment
			
		\begin{column}{.45\textwidth} % Left column and width
			\footnotesize
			Design parameters:
			\design{
				\begin{itemize}\itemsep0em 
					\only<1-2>{
						\item vane width
						\item vane height
						\item lean angle
					}%
					\only<3>{
						\item vane width $w = 6.78$
						\item vane height $h = 125.60$
						\item lean angle $\theta = 14.6$
					}%
					\only<4>{
						\item vane width $w = 11.5$
						\item vane height $h = 113.33$
						\item lean angle $\theta = 18.0$
					}%
					\only<5->{
						\item vane width $w = 5.42$
						\item vane height $h = 106.14$
						\item lean angle $\theta = 15.7$
					}%
					\item Material - yield stress $\sigma_y$ (off-the shelf)
				\end{itemize}
			}%
			Input Specifications:
			\spec{
				\begin{itemize}\itemsep0em 
					\item Nacelle temperature $T_1$ ($\downarrow$)
					\item Gas surface temperature $T_2$ ($\uparrow$)
				\end{itemize}
			}%
			Design margins:
			\margin{
				\begin{itemize}\itemsep0em 
					\item compression $\sigma_y - \sigma_a \ge 0$
					\item bending $\sigma_y - \sigma_m \ge 0$
					\item buckling $F_\text{buckling} - F_a \ge 0$
				\end{itemize}
			}%
			Performance parameters:
			\perf{
				\begin{itemize}\itemsep0em 
					\item Weight
				\end{itemize}
			}%
		\end{column}
	
		\begin{column}{.4\textwidth} % Left column and width
			\only<-2>{\uncover<2->{We aggregate across each \perf{performance} and \spec{specification}}}%
			\only<3->{%
				Design A: $d = 2.70$\\
				\uncover<4->{Design B: $d = 3.13$}\\
				\uncover<5->{Design C: $d = -3.77$}
			}%
			\renewcommand\baselinestretch{0.5}
			\renewcommand{\kbldelim}{[}% Left delimiter
			\renewcommand{\kbrdelim}{]}% Right delimiter
			\only<1>{
				\[
				  	\text{Impact} = \kbordermatrix{
										& \text{weight} \\
					\text{buckling} 	& 0.545  		\\
					\text{compression} 	& 0.111 		\\
					\text{bending} 		& -1.011		\\
				  }
				\]
				\[
				  	\text{absorption} = \kbordermatrix{
										& T_1	& T_2	\\
					\text{buckling} 	&2.79	& 3.87	\\
					\text{compression} 	&1.98	& 3.04	\\
					\text{bending} 		&2.58	& 3.66	\\
				  }
				\]
			}%
			\only<2>{
				\[
				  	\text{Impact} = \kbordermatrix{
										&  			\\
					\text{buckling} 	& 0.545		\\
					\text{compression} 	& 0.111		\\
					\text{bending} 		& -1.011	\\
				  }
				\]
				\[
				  	\text{absorption} = \kbordermatrix{
										& 		\\
					\text{buckling} 	&3.33	\\
					\text{compression} 	&2.51	\\
					\text{bending} 		&3.12	\\
				  }
				\]
			}%
			\only<3->{
				\begin{figure}
					\only<3>{
						\includegraphics[width=6.0cm]{design_space_exploration/MVP/MVP_1.pdf}
					}%
					\only<4>{
						\includegraphics[width=6.0cm]{design_space_exploration/MVP/MVP_2.pdf}
					}%
					\only<5->{
						\includegraphics[width=6.0cm]{design_space_exploration/MVP/MVP_3.pdf}
					}%
				\end{figure}
			}%
		\end{column}
		
	\end{columns}
\end{frame}
%------------------------------------------------
\subsection{Implementation}
%------------------------------------------------
\begin{frame}[fragile,t]
	\frametitle{Pseudo-code of MVM}
    %
	\vspace{-1em}
	\begin{algorithm}[H] % must use [H]
		\tiny
		\DontPrintSemicolon % Some LaTeX compilers require you to use \dontprintsemicolon instead
		\KwIn{
			set of possible choices $\mathcal{D}$, chosen design $\mathbf{c}\in\mathcal{D}$,
			the MAN given by analytic functions, models, and surrogates,
			$n_\mathrm{samples}$, Joint PDF $F_\mathbf{S}(\mathbf{s})$
		}
		\KwOut{estimated mean impact and change absorption $\bar{\bar{\mathbf{a}}}$, $\bar{\bar{\mathbf{i}}}$}
		%
		\textbf{[0] Initialization}: 
			set counter $i_\mathrm{sample} \gets 0$, initialize aggregate absorption and impact $\bar{\bar{\mathbf{a}}} \in \mathbb{R}^e$ and $\bar{\bar{\mathbf{i}}} \in \mathbb{R}^e$\;
		%
		\textbf{[1] Monte-carlo simulation}: 
			draw sample input specifications $\mathbf{S}\in\mathbb{R}^n \sim F_\mathbf{S}(\mathbf{s})$\;
		%
		\textbf{[2] Calculate threshold, decided, excess, and performance values}: 
			$\mathbf{t} = \mathbf{f}(\mathbf{S},\mathbf{c})$, 
			$\mathbf{d} = \mathbf{g}(\mathbf{c})$, 
			$\mathbf{e} = \left(\mathbf{t} - \mathbf{d}\right){\mathbf{i}^\mathrm{excess}}^\mathsf{T}$, 
			$\mathbf{p} =  \mathbf{h}(\mathbf{S},\mathbf{c})$\;
		%
		\textbf{[3] Calculate allowable deterioration}\;
		\qquad loop over input specifications \For{$k = 1, 2, ..., n$} {
			\qquad perform a line search along the $k$th direction to identify the root of: 
				$\min{((\mathbf{f}(\mathbf{x},\mathbf{c}) - \mathbf{d}){\mathbf{i}^\mathrm{excess}}^\mathsf{T})} = \mathbf{0}$, $s_k' \gets x^*$, 
				$s_k^\mathrm{limit} \gets s_k'$\;
		}
		%
		\textbf{[4] Calculate margin value metrics}\;
		\qquad initialization tensors and vectors, 
			$\mathbf{I} \in \mathbb{R}^{m\times e} \gets \mathbf{0}$ and $\mathbf{A} \in \mathbb{R}^{n\times e} \gets \mathbf{0}$, 
			$\mathbf{v} \in \mathbb{R}^{n} \gets \mathbf{0}$, $\bar{\mathbf{a}} \in \mathbb{R}^{e} \gets \mathbf{0}$, and $\bar{\mathbf{i}} \in \mathbb{R}^{e} \gets \mathbf{0}$\;
		\qquad loop over margin nodes \For{$i = 1, 2, ..., e$} {  
			\qquad \textbf{[a] Calculate impact on performance}\;
			\qquad $\mathbf{d}^\mathrm{threshold} \gets \mathbf{d}$, $d_i^\mathrm{threshold} \gets t_i$, 
				$\mathbf{c}_i^\mathrm{threshold} = \mathbf{g}^{-1}(\mathbf{d}^\mathrm{threshold})$, 
				$\mathbf{p}_i^\mathrm{threshold} = \mathbf{h}(\mathbf{S},\mathbf{c}_i^\mathrm{threshold})$\;
			\qquad loop over performance parameters \For{$j = 1, 2, ..., m$} {
				\qquad $i_{i,j} \gets (p_j - p_{j,i}^\mathrm{threshold})i_j^p/{p_{j,i}^\mathrm{threshold}}$\;
			}
			\qquad \textbf{[b] Calculate change absorption capability}\;
			\qquad loop over input specifications \For{$k = 1, 2, ..., n$} {
				\qquad calculate deterioration $v_k=(s_k^\mathrm{limit} - s_k)/{s_k}$\;
				\qquad $a_{i,k} \gets (f_i(\mathbf{s}_k^\mathrm{limit},\mathbf{c}) - f_i(\mathbf{S},\mathbf{c}))/(f_i(\mathbf{S},\mathbf{c})v_k)$\;
			}
			\qquad \textbf{[c] Aggregate metrics}: 
				$\bar{a}_i \gets \frac{1}{n}\sum_{k=1}^{n} a_{i,k}$, 
				$\bar{i}_i \gets \frac{1}{m}\sum_{k=1}^{m} i_{i,j}$
		}
		%
		\textbf{[5] sum of monte-carlo samples}: 
			$\bar{\bar{\mathbf{a}}} \gets \bar{\bar{\mathbf{a}}} + \bar{\mathbf{a}}$, 
			$\bar{\bar{\mathbf{i}}} \gets \bar{\bar{\mathbf{i}}} + \bar{\mathbf{i}}$\;
		%
		\textbf{[6] Termination}, 
		\If{$i_\mathrm{sample} < n_\mathrm{samples}$} {
			set $i_\mathrm{sample} \gets i_\mathrm{sample}+1$, go to \textbf{[1]}\;
			otherwise $\bar{\bar{\mathbf{a}}} \gets \bar{\bar{\mathbf{a}}}/n_\mathrm{samples}$ and $\bar{\bar{\mathbf{i}}} \gets \bar{\bar{\mathbf{i}}}/n_\mathrm{samples}$, \textbf{stop}
		}
	\end{algorithm}
\end{frame}
%------------------------------------------------
\begin{frame}[t,noframenumbering]
	\frametitle{UML diagram of \texttt{mvmlib}\ifshowcitations\footpartcite{Alhandawi2022}\fi}
    %
	\begin{figure}
        \centering
		\only<1>{\includegraphics[width=1.0\textwidth]{design_space_exploration/classes_mvmlib.pdf}}%
		\only<2>{\includegraphics[width=0.6\textwidth]{design_space_exploration/packages_mvmlib.pdf}}%
    \end{figure}%
\end{frame}
\addtocounter{footnote}{-1}
%------------------------------------------------

% Section 04: closing remarks
%%%%%%%%%%%%%%%%%%%%%%%%%%%%%%%%%%%%%%%%%%%%%%%%%%%%%%%
%%                    Conclusion                     %%
%%%%%%%%%%%%%%%%%%%%%%%%%%%%%%%%%%%%%%%%%%%%%%%%%%%%%%%
%------------------------------------------------
\begin{frame}[t,noframenumbering]
	\setbeamercovered{transparent}
	\frametitle{Conclusion and summary}
    %
	We demonstrated how a change in design requirements can be absorbed using:
	%
	\begin{columns}[t]
		\begin{column}{.5\textwidth} % Left column and width
           	~\\%
			%
			\uncover<2->{
				(a) Design \emphasis{changeability}:
				%
				\begin{itemize}
					\item Using a special kind of flexibility known as \textit{scalability},
					\item which is an enabler for \textit{remanufacturing}
				\end{itemize}
			}
			%
			\uncover<3->{
				(b) Design \emphasis{margins}:
				%
				\begin{itemize}
					\item By identifying \textit{excess} and \textit{buffer},
					\item and optimizing away excess
				\end{itemize}
			}
		\end{column}
		
		\begin{column}{.5\textwidth} % Left column and width
		\end{column}
	\end{columns}
\end{frame}

%------------------------------------------------
\begin{frame}
	\centering
	\Huge \emphasis{Thank you for your time}
\end{frame}
%------------------------------------------------

% Section 05: IATA aviation
%%%%%%%%%%%%%%%%%%%%%%%%%%%%%%%%%%%%%%%%%%%%%%%%%%%%%%%
%%                    The web app                    %%
%%%%%%%%%%%%%%%%%%%%%%%%%%%%%%%%%%%%%%%%%%%%%%%%%%%%%%%
%------------------------------------------------
\section{Web application for determining scalability}
% \frame{\tableofcontents[
% 	currentsection, 
% 	hideallsubsections
% 	]}
%------------------------------------------------
\begin{frame}[c,noframenumbering]
	\centering
	% \setlength\fboxsep{0pt}
	\begin{titleblock}{}
		~\\%
		{\centering\LARGE Software tools for assisting designers\\}%
		~\\%
	\end{titleblock}
\end{frame}
%------------------------------------------------
\begin{frame}[t]
\frametitle{Scale AM: an open-source tool for scalability assessment \ifshowcitations\textsuperscript{[1,2]}\fi}

\begin{columns}
    \begin{column}{.4\textwidth} % Left column and width
        \vspace{-2em}
        \begin{itemize}
            \item<2-> Takes as input DoE data 
            %from a design of experiments (DOE) run (e.g., Latin Hypercube, Full factorial, etc.)
            \item<3-> Provides a UI for providing \textit{change effect} and \textit{monotonicity} vectors %based on problem dimensionality
            \item<4-> Plots selected 2D projections of the parameter space and kernel smoothing model
            \item<5-> Displays Jacobian components \uncover<6->{\color{red}and their intersection}
        \end{itemize}
    \end{column}
    
    \begin{column}{.5\textwidth} % Left column and width
        \vspace{-2em}
        \begin{figure}
            \only<1>{\includegraphics[height=7cm]{scalability/web_app/0_example_trs.png}}%
            \only<2>{\includegraphics[height=7cm]{scalability/web_app/1_example_trs.png}}%
            \only<3>{\includegraphics[height=7cm]{scalability/web_app/2_example_trs.png}}%
            \only<4>{\includegraphics[height=7cm]{scalability/web_app/3_example_trs.png}}%
            \only<5>{\includegraphics[height=7cm]{scalability/web_app/4_example_trs.png}}%
            \only<6>{\includegraphics[height=7cm]{scalability/web_app/5_example_trs.png}}%
        \end{figure}
    \end{column}
\end{columns}

\vspace{-3em}
{\color{white}
\ifshowcitations
    \footpartcite{scaleAMGIT}[1]
    \footpartcite{scaleAM}[2]
\fi
}
\end{frame}
\addtocounter{footnote}{-2}

% Section 05: IATA aviation
%%%%%%%%%%%%%%%%%%%%%%%%%%%%%%%%%%%%%%%%%%%%%%%%%%%%%%%
%%               Graph representation                %%
%%%%%%%%%%%%%%%%%%%%%%%%%%%%%%%%%%%%%%%%%%%%%%%%%%%%%%%
%------------------------------------------------
\section{Graph representation of aviation network}
% \frame{\tableofcontents[
% 	currentsection, 
% 	hideallsubsections
% 	]}
\definecolor{gold}{HTML}{FFD700}
\definecolor{grey}{HTML}{D3D3D3}
%------------------------------------------------
\begin{frame}[c,noframenumbering]
	\centering
	% \setlength\fboxsep{0pt}
	\begin{titleblock}{}
		~\\%
		{\centering\LARGE Other projects: Graph representation of aviation network\\}%
		~\\%
	\end{titleblock}
\end{frame}
%------------------------------------------------
\begin{frame}[t]
    \setbeamercovered{transparent}
	\newcommand{\yshift}{0.0cm}
	\frametitle{Overview of graphs}
    %        
	\tikzstyle{background grid}=[draw, black!50,step=.5cm]
	\begin{tikzpicture}[remember picture, overlay] %show background grid, 
		\node [inner sep=0pt,above right, opacity=1.0]  at (0.49\textwidth,-0.83\textheight+\yshift) (barchart) 
		{
			\only<1>{\includegraphics[height=8cm]{GRL/networks/net_example_1.pdf}}%
            \only<2>{\includegraphics[height=8cm]{GRL/networks/net_example_2.pdf}}%
            \only<3>{\includegraphics[height=8cm]{GRL/networks/net_example_3.pdf}}%
            \only<4>{\includegraphics[height=8cm]{GRL/networks/net_example_4.pdf}}%
            \only<5>{\includegraphics[height=8cm]{GRL/networks/net_2019.pdf}}%
            % \only<6>{\includegraphics[height=8cm]{GRL/networks/net_alliance_2019.pdf}}%
		};
		% show origin
		% \fill (0,0) circle (2pt);
	\end{tikzpicture}%
    %
    \vspace{-2em}
    \begin{columns}[t] % The "c" option specifies centered vertical alignment while the "t" option is used for top vertical alignment

        \begin{column}{.5\textwidth} % Left column and width
            \begin{itemize}
                \item<1-> {\color{gold} Nodes}: represent the airlines
                \item<2-> {\color{grey} Edges}: represent the codesharing
                \item<3-> Edge weights given by thickness of lines
                \item<4-> Node strength: sum of incoming edge weights at each node (size of nodes)
                \item<5-> We identify large communities
                % \item<6-> Compare them to \textbf{six} airline alliances
            \end{itemize}
            \uncover<6->{
                \begin{alertblock}{Finding optimal community structure}
                    \begin{equation*}
                        \begin{aligned} 
                            & \underset{\left\{\mathcal{A}_1, \cdots, \mathcal{A}_K\right\}\in\mathcal{V}^K}{\text{min}}
                            & & -Q\left(\mathcal{A}_1, \cdots, \mathcal{A}_K\right) \\
                            & \text{where}
                            & & K \in \mathbb{N},~\mathcal{V}: \text{Nodes} \\
                            & & & Q :\text{modularity}
                        \end{aligned}
                    \end{equation*}
                \end{alertblock}
            }
        \end{column}

        \begin{column}{.5\textwidth} % Left column and width
                
        \end{column}
    \end{columns}
\end{frame}

% Section 06: agent based model

%%%%%%%%%%%%%%%%%%%%%%%%%%%%%%%%%%%%%%%%%%%%%%%%%%%%%%%
%%                        ABMs                       %%
%%%%%%%%%%%%%%%%%%%%%%%%%%%%%%%%%%%%%%%%%%%%%%%%%%%%%%%
%------------------------------------------------
\section{Optimization of stochastic agent-based models}
%------------------------------------------------
\begin{frame}[c,noframenumbering]
	\centering
	% \setlength\fboxsep{0pt}
	\begin{titleblock}{}
		~\\%
		{\centering\LARGE Optimization of public health models using artificial life models\\}%
		~\\%
	\end{titleblock}
\end{frame}
%------------------------------------------------
\subsection{Public health policy-making problem}
%------------------------------------------------
\begin{frame}[t,label=abm_3]
	\frametitle{Public health policy-making problem formulation}
	\tikzstyle{background grid}=[draw, black!50,step=.5cm]
	%
	What is the \emphasis{cost} of public health interventions?\\
	%
	\begin{columns}[t] % The "c" option specifies centered vertical alignment while the "t" option is used for top vertical alignment
		\begin{column}{.5\textwidth} % Left column and width

			\begin{tikzpicture}[remember picture, overlay] %show background grid, 
				\node [inner sep=0pt,above right, opacity=1.0]  at (0.0\textwidth,-0.66\textheight) (abm) 
					{
						\only<-2>{
							\includegraphics[width=0.95\textwidth]{agent_based_models/render_2_no_restrictions/sim_10.pdf}
						}%	
						\only<3>{
							\begin{animateinline}[autoplay,width=0.95\textwidth]{70}
								\ifshowanimations
									\multiframe{77}{i=10+45}{%
										\includegraphics{agent_based_models/render_2_no_restrictions_anim/sim_\i.pdf}
									}
								\else
									\multiframe{1}{i=10+0}{%
										\includegraphics{agent_based_models/render_2_no_restrictions/sim_\i.pdf}
									}
								\fi
							\end{animateinline}%
						}%
						\only<4->{
							\includegraphics[width=0.95\textwidth]{agent_based_models/render_2_no_restrictions/sim_3500.pdf}
						}%
					};
				\only<1->{\node[inner sep=0pt,align=flush center,above=\belowcaptionskip of abm,text width=\linewidth]
					{\vspace{-1em}{\large No interventions applied}};}%
				%
			\end{tikzpicture}%

		\end{column}
		%
		\begin{column}{.5\textwidth} % Left column and width

			\tikzstyle{background grid}=[draw, black!50,step=.1cm]
			\begin{tikzpicture}[remember picture, overlay] %show background grid, 
				\node [inner sep=0pt,above right, opacity=1.0]  at (0.0\textwidth,-0.66\textheight) (abm) 
					{
						\only<2>{
							\includegraphics[width=0.95\textwidth]{agent_based_models/render_1_restrictions/sim_50.pdf}
						}%	
						\only<3>{
							\begin{animateinline}[autoplay,width=0.95\textwidth]{70}
								\ifshowanimations
									\multiframe{75}{i=50+45}{%
										\includegraphics{agent_based_models/render_1_restrictions_anim/sim_\i.pdf}
									}
								\else
									\multiframe{1}{i=50+0}{%
										\includegraphics{agent_based_models/render_1_restrictions/sim_\i.pdf}
									}
								\fi
							\end{animateinline}%
						}%
						\only<4->{
							\includegraphics[width=0.95\textwidth]{agent_based_models/render_1_restrictions/sim_3500.pdf}
						}%
					};
				\only<3-7>{\node[inner sep=0pt,align=flush center,above=\belowcaptionskip of abm,text width=\linewidth]
					{\vspace{-1em}{\large with intervention}};}%
				% show origin
				% \fill (0,0) circle (2pt);
			\end{tikzpicture}%
		\end{column}

	\end{columns}
	\vspace{-3em}
\end{frame}
%------------------------------------------------
\begin{frame}[t,label=abm_4]
	\frametitle{Optimization problem}
	\tikzstyle{background grid}=[draw, black!50,step=.5cm]
	%
	\uncover<2->{No gradient information available, blackbox is expensive and \only<3->{\emphasis}{noisy}}\\
	%
	\begin{columns}[t] % The "c" option specifies centered vertical alignment while the "t" option is used for top vertical alignment
		\begin{column}{.42\textwidth} % Left column and width
			\vspace{-1.2em}
			% Optimization problem
			\begin{exampleblock}{Objective and constraints}
				\only<1-3>{
					\begin{equation*}
						\begin{aligned}
							& \underset{\mathbf{x}}{\text{min}}
							& & f(\mathbf{x}) = -M^{T}\\
							& \text{subject to}
							& & \uncover<2->{{c}(\mathbf{x}) \equiv n_{I,\text{max}} - H_{\text{max}} \le 0}\\
							& \text{where}
							& & \mathbf{x}=\left[n_E,S_D,n_T\right]^\mathit{T}\\
						\end{aligned}
					\end{equation*}
				}%
				\only<4->{
					\vspace{-1.2em}
					\begin{equation*}
						\begin{aligned}
							& \underset{\mathbf{x}}{\text{min}}
							& & f(\mathbf{x}) = \mathbb{E}_{\Theta}\left[{f}_{\Theta}(\mathbf{x}) = -M^{T}\right]\\
							& \text{subject to}
							& & {c}(\mathbf{x}) = \mathbb{E}_{\Theta}\left[{c}_{\Theta}(\mathbf{x}) \equiv n_{I,\text{max}} - H_{\text{max}}\right] \le 0\\
							& \text{where}
							& & \mathbf{x}=\left[n_E,S_D,n_T\right]^\mathit{T},{\color{red}\Theta\mathrm{:realizations}}
						\end{aligned}
					\end{equation*}
				}%
			\end{exampleblock}
			\vspace{-0.5em}
			\uncover<1->{
				% Variables
				\begin{alertblock}{Design variables}
					\vspace{-0.0em}
						\begin{itemize}\itemsep0em
							\item $n_E:$ Number of essential workers
							\item $S_D:$ Social distancing factor
							\item $n_T:$ Number of tests daily
						\end{itemize}
				\end{alertblock}
			}%
			\vspace{-0.5em}
			\uncover<5->{
				% Parameters
				\begin{blueblock}{Randomly seeded parameters}
					\vspace{-0.0em}
					\begin{itemize}\itemsep0em
						\item Initial conditions
						\item Interactions, demographics
					\end{itemize}
				\end{blueblock}
			}%
		\end{column}
		%
		\begin{column}{.5\textwidth} % Left column and width
			\tikzstyle{background grid}=[draw, black!50,step=.5cm]
			\begin{tikzpicture}[remember picture, overlay]%[show background grid]
				\only<1-5>{
					\node [inner sep=0pt,above right, opacity=1.0] at (0, -4.5) (objective)
						{
							\only<1->{\includegraphics[width=0.7\textwidth]{agent_based_models/trajectory_demo/M_compare_opt_3.pdf}}%
						};
					\only<1->{
						\node [above right] at (3.2,-3.5) (a1) {};
						\node [above right] at (2.2,-1.5) (a2) {};
						\draw[-stealth,line width=2pt,black] (a1) to [out=90, in=310] (a2);
					}
					\node [inner sep=0pt,above right, opacity=0.7] at (2.5, -6.5) (constraint)
						{
							\only<2->{\includegraphics[width=0.7\textwidth]{agent_based_models/trajectory_demo/I_compare_opt_3.pdf}}%
						};
					\only<2->{
						\node [above right] at (4.7,-3.5) (a1) {};
						\node [above right] at (3.7,-5.5) (a2) {};
						\draw[-stealth,line width=2pt,black] (a1) to [out=210, in=90] (a2);
					}					

				}%
				% show origin
				% \fill (0,0) circle (2pt);
				% define destination coordinates
			\end{tikzpicture}%
		\end{column}
	
	\end{columns}
	\vspace{-3em}
\end{frame}

% % \includeonlyframes{abm_2}
% % Section 02: machine learning work
% %==============================================================================%
%% Mathematical symbols
\newcommand{\Ti}{\mbox{$T_1$}}
\newcommand{\nh}{\mbox{$n_\text{hidden}$}}
\newcommand{\Pd}{\mbox{$P_\text{dropout}$}}
\newcommand{\nla}{\mbox{$n_\text{layers}$}}
\newcommand{\Pt}{\mbox{$P_\text{teacher}$}}
\newcommand{\lr}{\mbox{$l_\text{rate}$}}
\newcommand{\bs}{\mbox{$b_\text{size}$}}
\newcommand{\R}{\mbox{$\lambda$}}
\newcommand{\e}{\mbox{$\epsilon$}}
\newcommand{\etol}{\mbox{$\epsilon_\text{tol}$}}
\newcommand{\fs}{\mbox{$f_\text{sample}$}}
\newcommand{\ff}{\mbox{$f_\text{features}$}}
\newcommand{\D}{\mbox{$D$}}
\newcommand{\ns}{\mbox{$n_\text{stages}$}}
\newcommand{\nd}{\mbox{$n_\text{degree}$}}

\newcommand{\To}{\mbox{$T_2$}}
\newcommand{\np}{\mbox{$n_\text{patience}$}}
\newcommand{\nepoch}{\mbox{$n_\text{epochs}$}}
\newcommand{\nbest}{\mbox{$n_\text{epochs}^\text{best}$}}
%==============================================================================%
%------------------------------------------------
\section{Hyperparameter tuning of machine learning problems}
%------------------------------------------------
\begin{frame}[c,noframenumbering]
	\centering
	% \setlength\fboxsep{0pt}
	\begin{titleblock}{}
		~\\%
		{\centering\LARGE COVID-19 incidence now-casting using machine learning\\}%
		~\\%
	\end{titleblock}
\end{frame}
%------------------------------------------------
\subsection{Overview of problem}
%------------------------------------------------
\begin{frame}[t]
	\frametitle{COVID-19 incidence forecasting}
	\tikzstyle{background grid}=[draw, black!50,step=.5cm]
	%
	How to \emph{design} a machine learning model for forecasting a time-series?\\
	%
	\tikzstyle{background grid}=[draw, black!50,step=.5cm]
	\begin{tikzpicture}[remember picture, overlay] %show background grid, 
		% Put the graphic inside a node. This makes it easy to place the
		% graphic and to draw on top of it. 
		% The above right option is used to place the lower left corner
		% of the image at the (0,0) coordinate. 
		\node [inner sep=0pt,above right, opacity=1.0]  at (0.0\textwidth,-0.72\textheight) (raw data) 
			{
				\only<1>{\includegraphics[width=0.43\textwidth]{machine_learning/raw_data/animation_0.pdf}}%
				\only<2>{\includegraphics[width=0.43\textwidth]{machine_learning/raw_data/animation_1.pdf}}%
				\only<3>{\includegraphics[width=0.43\textwidth]{machine_learning/raw_data/animation_2.pdf}}%
				\only<4->{\includegraphics[width=0.43\textwidth]{machine_learning/raw_data/animation_3.pdf}}%
			};
		\node [inner sep=0pt,above left, opacity=1.0]  at (0.99\textwidth,-0.65\textheight) (seq2seq) 
			{
				\only<5>{\includegraphics[width=0.55\textwidth]{machine_learning/model_1.pdf}}%
				\only<6>{\includegraphics[width=0.55\textwidth]{machine_learning/model_2.pdf}}%
				\only<7>{\includegraphics[width=0.55\textwidth]{machine_learning/model_3.pdf}}%
			};
		\only<5->{
			\node[inner sep=0pt,align=flush center,below=\belowcaptionskip of seq2seq,text width=\linewidth]
				{\vspace{-1em}{Seq2Seq model}};
		}%
		\node [inner sep=0pt, opacity=1.0]  at (0.43\textwidth,-0.15\textheight) (ct) {};%
		\node [inner sep=0pt, opacity=1.0]  at (0.43\textwidth,-0.50\textheight) (cases) {};%
		\only<-6>{
			\node [inner sep=0pt, opacity=1.0]  at (0.65\textwidth,-0.50\textheight) (input) {};%
		}%
		\only<7->{
			\node [inner sep=0pt, opacity=1.0]  at (0.60\textwidth,-0.60\textheight) (input) {};%
		}%
		% show origin
		% \fill (0,0) circle (2pt);
	\end{tikzpicture}%
	%
	\begin{tikzpicture}[overlay]
		\only<5->{\path[->,magenta,thick] (ct) edge [out=0, in=270] (input);}
		\only<5->{\path[->,magenta,thick] (cases) edge [out=0, in=270] (input);}
	\end{tikzpicture}%
	%
	\vspace{-3em}
\end{frame}
%------------------------------------------------
\subsection{Hyperparameter tuning}
%------------------------------------------------
\begin{frame}[t]
    \frametitle{Training the model}
	%
	There are several challenges associated with hyperparameter optimization\\
	%
    \begin{columns}[t] % The "c" option specifies centered vertical alignment while the "t" option is used for top vertical alignment
		
        \begin{column}{.5\textwidth} % Left column and width
        \only<1->{
			\begin{itemize}
				\item The number of epochs can be tuned using \textit{early stopping}%
				\item This is a form of \textit{regularization} to reduce over-fitting%
			\end{itemize}
        }%
        \only<3->{
			~~~However,\\
			\uncover<4->{
				~~~1) Other hyperparameters to tune
				\begin{itemize}
					\item Dropout%
					\item Training batch size%
					\item Activation function(s)%
				\end{itemize}
			}%
            \uncover<5->{~~~2) Training can be \emphasis{expensive}}\\
            \uncover<6->{~~~3) Backpropagation is \emphasis{stochastic}}
        }%
    
        \end{column}
    
        \begin{column}{.5\textwidth} % Left column and width
    
            \vspace{-1.5em}
            \begin{figure}
                \centering
                \only<1>{\includegraphics[width=0.9\textwidth]{machine_learning/training/epoch_0.pdf}}%
                \only<2>{
					\begin{animateinline}[autoplay,width=0.9\textwidth]{8}
						\ifshowanimations
							\multiframe{52}{i=0+5}{%
								\includegraphics{machine_learning/training_anim/epoch_\i.pdf}
							}
						\else
							\multiframe{1}{i=262+0}{%
								\includegraphics{machine_learning/training/epoch_\i.pdf}
							}
						\fi
					\end{animateinline}%
                }%
                \only<3->{\includegraphics[width=0.9\textwidth]{machine_learning/training/final_model.pdf}}%
                \vspace{-0.75em}
                \caption{Effect of number of epochs on testing loss}
            \end{figure}%
        \end{column}
    \end{columns}
        
\end{frame}
%------------------------------------------------
\begin{frame}[t]
	\frametitle{Hyperparameter tuning}
	\tikzstyle{background grid}=[draw, black!50,step=.5cm]
	%
	We can use StoMADS to solve such hyperparameter optimization problems \ifshowcitations\footpartcite{Khalil2021}\fi\\
	%
	\tikzstyle{background grid}=[draw, black!50,step=.5cm]
	\begin{tikzpicture}[remember picture, overlay] %show background grid, 
		% Put the graphic inside a node. This makes it easy to place the
		% graphic and to draw on top of it. 
		% The above right option is used to place the lower left corner
		% of the image at the (0,0) coordinate. 
		\node [inner sep=0pt,above right, opacity=1.0]  at (-0.01\textwidth,-0.7\textheight) (error) 
			{
				\only<1>{\includegraphics[width=0.5\textwidth]{machine_learning/box_plots/boxplot_input_dim_0.pdf}}%
				\only<2>{\includegraphics[width=0.5\textwidth]{machine_learning/box_plots/boxplot_input_dim_1.pdf}}%
				\only<3>{\includegraphics[width=0.5\textwidth]{machine_learning/box_plots/boxplot_input_dim_2.pdf}}%
				\only<4>{\includegraphics[width=0.5\textwidth]{machine_learning/box_plots/boxplot_input_dim_3.pdf}}%
				\only<5>{\includegraphics[width=0.5\textwidth]{machine_learning/box_plots/boxplot_input_dim_4.pdf}}%
				\only<6>{\includegraphics[width=0.5\textwidth]{machine_learning/box_plots/boxplot_input_dim_5.pdf}}%
				\only<7>{\includegraphics[width=0.5\textwidth]{machine_learning/box_plots/boxplot_input_dim_6.pdf}}%
				\only<8>{\includegraphics[width=0.5\textwidth]{machine_learning/box_plots/boxplot_input_dim_opt.pdf}}%
				\only<9>{\includegraphics[width=0.5\textwidth]{machine_learning/box_plots/boxplot_hid_dim_opt.pdf}}%
				\only<10>{\includegraphics[width=0.5\textwidth]{machine_learning/box_plots/boxplot_dropout_opt.pdf}}%
			};
		\node [inner sep=0pt,above left, opacity=1.0]  at (1.01\textwidth,-0.7\textheight) (prediction) 
			{
				\only<1>{\includegraphics[width=0.5\textwidth]{machine_learning/models/model_input_dim_0.pdf}}%
				\only<2>{\includegraphics[width=0.5\textwidth]{machine_learning/models/model_input_dim_1.pdf}}%
				\only<3>{\includegraphics[width=0.5\textwidth]{machine_learning/models/model_input_dim_2.pdf}}%
				\only<4>{\includegraphics[width=0.5\textwidth]{machine_learning/models/model_input_dim_3.pdf}}%
				\only<5>{\includegraphics[width=0.5\textwidth]{machine_learning/models/model_input_dim_4.pdf}}%
				\only<6>{\includegraphics[width=0.5\textwidth]{machine_learning/models/model_input_dim_5.pdf}}%
				\only<7>{\includegraphics[width=0.5\textwidth]{machine_learning/models/model_input_dim_6.pdf}}%
				\only<8>{\includegraphics[width=0.5\textwidth]{machine_learning/models/final_model_input_dim.pdf}}%
				\only<9>{\includegraphics[width=0.5\textwidth]{machine_learning/models/final_model_hid_dim.pdf}}%
				\only<10->{\includegraphics[width=0.5\textwidth]{machine_learning/models/final_model_dropout.pdf}}%
			};
		% show origin
		% \fill (0,0) circle (2pt);
	\end{tikzpicture}%
	%
	\begin{columns}[t] % The "c" option specifies centered vertical alignment while the "t" option is used for top vertical alignment
		\begin{column}{.42\textwidth} % Left column and width
			\vspace{-2.0em}
			% Optimization problem
			\uncover<11->{
				\begin{exampleblock}{Objective and constraints}
					\vspace{-0.0em}
					\begin{equation*}
						\begin{aligned}
							& \underset{\mathbf{x}}{\text{min}}
							& & f(\mathbf{x}) = \mathbb{E}_{\Theta}\left[{f}_{\Theta}(\mathbf{x}) = \mathrm{error}_\mathrm{CV}\right]\\
							& \text{where}
							& & \Theta\mathrm{:realizations}
						\end{aligned}
					\end{equation*}
				\end{exampleblock}
			}%
			\vspace{-0.5em}
			\uncover<11->{
				% Variables
				\begin{alertblock}{Design variables ($\mathbf{x}$)}
					\vspace{-0.0em}
						\begin{itemize}\itemsep0em
							\item $T_1:$ Input dimension
							\item $n_\text{hidden}:$ Number of hidden neurons
							\item $P_\text{dropout}:$ Probability of dropout, etc.
						\end{itemize}
				\end{alertblock}
			}%
			\vspace{-0.5em}
			\uncover<11->{
				% Parameters
				\begin{blueblock}{Randomly seeded parameters}
					\vspace{-0.0em}
					\begin{itemize}\itemsep0em
						\item Initial weights
						\item Gradient descent steps
					\end{itemize}
				\end{blueblock}
			}%
		\end{column}
		%
		\begin{column}{.5\textwidth} % Left column and width
		\end{column}
	
	\end{columns}
	%
	\vspace{-3em}
\end{frame}
\addtocounter{footnote}{-1}
%------------------------------------------------
\subsection{Results}
\subsubsection{Optimal hyperparameters}
%------------------------------------------------
\begin{frame}[t]
	\frametitle{Results: Optimal hyperparameters}
	\tikzstyle{background grid}=[draw, black!50,step=.5cm]
	%
	Optimal hyperparameters for the \emph{Seq2Seq} model:\\
	%
	\tikzstyle{background grid}=[draw, black!50,step=.5cm]
	\begin{tikzpicture}[remember picture, overlay] %show background grid, 
		% Put the graphic inside a node. This makes it easy to place the
		% graphic and to draw on top of it. 
		% The above right option is used to place the lower left corner
		% of the image at the (0,0) coordinate. 
		\node [inner sep=0pt,above left, opacity=1.0]  at (1.01\textwidth,-0.73\textheight) (prediction) 
			{\includegraphics[width=0.5\textwidth]{machine_learning/predictions/model_predictions_test_S2S_mean_Ct_daily_cases.pdf}};
		% show origin
		% \fill (0,0) circle (2pt);
	\end{tikzpicture}%
	%
	\begin{columns}[t] % The "c" option specifies centered vertical alignment while the "t" option is used for top vertical alignment
		\begin{column}{.42\textwidth} % Left column and width
			\vspace{-2.0em}
            % Column widths
            \newcommand{\ocwb}{3.7cm}
            \newcommand{\ocwc}{1cm}
            \newcommand{\ocwd}{1.2cm}
            \newcommand{\ocwe}{3cm}
            %
            \begin{table}[h!]
                \centering
                \footnotesize
                \renewcommand{\arraystretch}{1.5}% Wider
                \begin{tabular}{L{\ocwb}C{\ocwc}C{\ocwd}} \toprule
                    \multicolumn{2}{c}{\bf Hyperparameter}	& \bf Value         \\ \toprule
                    Sliding window size  		& \Ti		&	6				\\
                    Number of hidden neurons	& \nh		&	1500			\\
                    Probability of dropout		& \Pd		&	0.8				\\
                    Number of hidden layers		& \nh		&	2				\\
                    Teacher forcing probability	& \Pt		&	0.3				\\
                    Learning rate 				& \lr		&	$1\times10^{-4}$\\
                    batch size 					& \bs		&	32				\\\hline
                \end{tabular}
            \end{table}
        \end{column}
		%
		\begin{column}{.5\textwidth} % Left column and width
		\end{column}
	
	\end{columns}
	%
	\vspace{-3em}
\end{frame}
%------------------------------------------------
\begin{frame}[t]
	\frametitle{Results: Optimal hyperparameters}
	\tikzstyle{background grid}=[draw, black!50,step=.5cm]
	%
	Optimal hyperparameters for the \emph{support vector machine regression (SVR)} model:\\
	%
	\tikzstyle{background grid}=[draw, black!50,step=.5cm]
	\begin{tikzpicture}[remember picture, overlay] %show background grid, 
		% Put the graphic inside a node. This makes it easy to place the
		% graphic and to draw on top of it. 
		% The above right option is used to place the lower left corner
		% of the image at the (0,0) coordinate. 
		\node [inner sep=0pt,above left, opacity=1.0]  at (1.01\textwidth,-0.73\textheight) (prediction) 
			{\includegraphics[width=0.5\textwidth]{machine_learning/predictions/model_predictions_test_SVR_mean_Ct_daily_cases.pdf}};
		% show origin
		% \fill (0,0) circle (2pt);
	\end{tikzpicture}%
	%
	\begin{columns}[t] % The "c" option specifies centered vertical alignment while the "t" option is used for top vertical alignment
		\begin{column}{.42\textwidth} % Left column and width
			\vspace{-2.0em}
            % Column widths
            \newcommand{\ocwb}{3.7cm}
            \newcommand{\ocwc}{1cm}
            \newcommand{\ocwd}{1.2cm}
            \newcommand{\ocwe}{3cm}
            %
            \begin{table}[h!]
                \centering
                \footnotesize
                \renewcommand{\arraystretch}{1.5}% Wider
                \begin{tabular}{L{\ocwb}C{\ocwc}C{\ocwd}} \toprule
                    \multicolumn{2}{c}{\bf Hyperparameter}				    & \bf Value             \\ \toprule
                    Sliding window size 						& \Ti	    &	6					\\
                    Ridge factor								& \R	    &	$1\times10^{-4}$	\\
                    Margin of tolerance							& \e	    &	$1\times10^{-2}$	\\
                    Stopping criteria tolerance					& \etol	    &	0.1					\\
                    Learning rate 								& \lr	    &	$1\times10^{-5}$	\\ \hline
                \end{tabular}
            \end{table}
            \uncover<2->{Support vector machine models have \emphasis{deterministic} performance}
        \end{column}
		%
		\begin{column}{.5\textwidth} % Left column and width
		\end{column}
	
	\end{columns}
	%
	\vspace{-3em}
\end{frame}
%------------------------------------------------
\subsubsection{Prospective validation}
%------------------------------------------------
\begin{frame}[t]
	\frametitle{Results: Prospective validation}
	\tikzstyle{background grid}=[draw, black!50,step=.5cm]
	%
	Performance of models on \emphasis{unseen} data (first 4 months of 2021):\\
	%
	\tikzstyle{background grid}=[draw, black!50,step=.5cm]
	\begin{tikzpicture}[remember picture, overlay] %show background grid, 
		% Put the graphic inside a node. This makes it easy to place the
		% graphic and to draw on top of it. 
		% The above right option is used to place the lower left corner
		% of the image at the (0,0) coordinate. 
		\node [inner sep=0pt,above left, opacity=1.0]  at (1.01\textwidth,-0.73\textheight) (prediction) 
			{
                \only<1>{\includegraphics[width=0.5\textwidth]{machine_learning/predictions/model_predictions_unseen_S2S_mean_Ct_daily_cases_truth.pdf}}%
                \only<2>{\includegraphics[width=0.5\textwidth]{machine_learning/predictions/model_predictions_unseen_S2S_mean_Ct_daily_cases.pdf}}%
                \only<3>{\includegraphics[width=0.5\textwidth]{machine_learning/predictions/model_predictions_unseen_SVR_mean_Ct_daily_cases.pdf}}%
            };
		% show origin
		% \fill (0,0) circle (2pt);
	\end{tikzpicture}%
	%
	\begin{columns}[t] % The "c" option specifies centered vertical alignment while the "t" option is used for top vertical alignment
		\begin{column}{.42\textwidth} % Left column and width
			\vspace{-2.0em}
            % Column widths
            \newcommand{\ocwa}{5cm}
            \newcommand{\ocwb}{1.6cm}
            \newcommand{\ocwc}{1.6cm}
            \newcommand{\ocwd}{1.6cm}
            %
            \begin{table}[h!]
                \centering
                \footnotesize
                \renewcommand{\arraystretch}{1.5}% Wider
                \begin{tabular}{L{\ocwa}C{\ocwd}} \toprule
                    \multicolumn{1}{c}{\bf Model}               & \multicolumn{1}{c}{\bf Test error}    \\ \toprule
                    \only<2>{\emphasis}{Seq2Seq}		        & \only<2>{\emphasis}{0.571}            \\
                    Long short term memory (LSTM) cell          & 0.326                                 \\
                    feedforward neural network	                & 0.255                                 \\
                    \only<3>{\emphasis}{Support vector machine} & \only<3>{\emphasis}{0.168}            \\
                    Gradient boosting 		                    & 1.444                                 \\
                    Linear regression 		                    & 0.160                                 \\ \bottomrule
                \end{tabular}
            \end{table}
        \end{column}
        %
		\begin{column}{.5\textwidth} % Left column and width
		\end{column}
	
	\end{columns}
	%
	\vspace{-3em}
\end{frame}
%------------------------------------------------
\begin{frame}[t]
	\frametitle{Results: Prospective validation}
	\tikzstyle{background grid}=[draw, black!50,step=.5cm]
	%
	Effect of increasing number of training days (Adding 1 month of data):\\
	%
	\tikzstyle{background grid}=[draw, black!50,step=.5cm]
	\begin{tikzpicture}[remember picture, overlay] %show background grid, 
		% Put the graphic inside a node. This makes it easy to place the
		% graphic and to draw on top of it. 
		% The above right option is used to place the lower left corner
		% of the image at the (0,0) coordinate. 
		\node [inner sep=0pt,above left, opacity=1.0]  at (1.01\textwidth,-0.73\textheight) (prediction) 
			{
                \only<1>{\includegraphics[width=0.5\textwidth]{machine_learning/predictions/model_predictions_unseen_S2S_mean_Ct_daily_cases.pdf}}%
                \only<2>{\includegraphics[width=0.5\textwidth]{machine_learning/predictions/model_predictions_unseen_S2S_mean_Ct_daily_cases_G12.pdf}}%
                \only<3>{\includegraphics[width=0.5\textwidth]{machine_learning/predictions/model_predictions_unseen_SVR_mean_Ct_daily_cases_G12.pdf}}%
            };
		% show origin
		% \fill (0,0) circle (2pt);
	\end{tikzpicture}%
	%
	\begin{columns}[t] % The "c" option specifies centered vertical alignment while the "t" option is used for top vertical alignment
		\begin{column}{.42\textwidth} % Left column and width
			\vspace{-2.0em}
            % Column widths
            \newcommand{\ocwa}{5cm}
            \newcommand{\ocwb}{1.6cm}
            \newcommand{\ocwc}{1.6cm}
            \newcommand{\ocwd}{1.6cm}
            %
            \begin{table}[h!]
                \centering
                \footnotesize
                \renewcommand{\arraystretch}{1.5}% Wider
                \begin{tabular}{L{\ocwa}C{\ocwd}} \toprule
                    \multicolumn{1}{c}{\bf Model}               & \multicolumn{1}{c}{\bf Test error}                                \\ \toprule
                    \only<1-2>{\emphasis}{Seq2Seq}		        & \only<1>{\emphasis{0.571}}\only<2->{\bf\color{darkgreen}{0.106}}   \\
                    Long short term memory (LSTM) cell          & 0.326                                                             \\
                    feedforward neural network	                & 0.255                                                             \\
                    \only<3>{\emphasis}{Support vector machine} & \only<1-2>{0.168}\only<3>{\bf\color{darkgreen}{0.140}}            \\
                    Gradient boosting 		                    & 1.444                                                             \\
                    Linear regression 		                    & 0.160                                                             \\ \bottomrule
                \end{tabular}
            \end{table}
        \end{column}
        %
		\begin{column}{.5\textwidth} % Left column and width
		\end{column}
	
	\end{columns}
	%
	\vspace{-3em}
\end{frame}
%------------------------------------------------
\subsection{UML diagram}
%------------------------------------------------
\begin{frame}[t]
	\frametitle{UML diagram of code}
    %
	\begin{figure}
        \centering
		\only<1>{\includegraphics[width=1.0\textwidth]{machine_learning/classes_hpo.pdf}}%
		\only<2>{\includegraphics[width=1.0\textwidth]{machine_learning/packages_hpo.pdf}}%
    \end{figure}%
\end{frame}
%------------------------------------------------

% Supplementary slides
% %%%%%%%%%%%%%%%%%%%%%%%%%%%%%%%%%%%%%%%%%%%%%%%%%%%%%%%
%%      Supplementary Content: Design margins        %%
%%%%%%%%%%%%%%%%%%%%%%%%%%%%%%%%%%%%%%%%%%%%%%%%%%%%%%%
%------------------------------------------------
\begin{frame}[c,noframenumbering]
	\centering
	% \setlength\fboxsep{0pt}
	\begin{titleblock}{}
		~\\%
		{\centering\LARGE Supplementary content: Design margin allocation\\}%
		~\\%
	\end{titleblock}
\end{frame}
%------------------------------------------------
\subsection{Thermal loadcase explained}
%------------------------------------------------
\begin{frame}[t,label=ap2_1]
	\frametitle{Thermal loadcase explained}
	\vspace{-1em}
	\only<1>{Let us revisit the previous example of the TRS}%
	\only<2>{Consider four thermal loads $\mathbf{T} = \left[T_1,T_2,T_3,T_4\right]$ that can change}%
	\only<3-8>{A stiffener pattern can be deposited on the outercasing}%
	~\\
	\begin{columns}[t] % The "c" option specifies centered vertical alignment while the "t" option is used for top vertical alignment
			
		\begin{column}{.45\textwidth} % Left column and width
			\only<1-3>{\uncover<3>{"wavy" concept: $c=0${\color{white},$\mathbf{D}=\left[1\right]$}}}%
			\only<4>{"hatched" concept: $c=1${\color{white},$\mathbf{D}=\left[1\right]$}}%
			\only<5>{"tubular" concept:$c=2${\color{white},$\mathbf{D}=\left[1\right]$}}%
			\only<6>{$c=1$, Desposit steps: $\mathbf{D}=\left[1\right]$}%
			\only<7>{$c=1$, Desposit steps: $\mathbf{D}=\left[1,2\right]$}%
			\only<8->{$c=1$, Desposit steps: $\mathbf{D}=\left[1,2,4\right]$}%
			\begin{figure}
				\only<1-2>{\includegraphics[height=5.0cm]{design_margins/DM_geometry/TRS_isometric_1.pdf}}%
				\only<3>{\includegraphics[height=5.0cm]{design_margins/DM_geometry/TRS_isometric_2.pdf}}%
				\only<4>{\includegraphics[height=5.0cm]{design_margins/DM_geometry/TRS_isometric_3.pdf}}%
				\only<5>{\includegraphics[height=5.0cm]{design_margins/DM_geometry/TRS_isometric_4.pdf}}%
				\only<6>{\includegraphics[height=5.0cm]{design_margins/DM_geometry/TRS_isometric_5.pdf}}%
				\only<7>{\includegraphics[height=5.0cm]{design_margins/DM_geometry/TRS_isometric_6.pdf}}%
				\only<8->{\includegraphics[height=5.0cm]{design_margins/DM_geometry/TRS_isometric_7.pdf}}%
				\vspace{-0.2em}
				\caption{Turbine Rear Structure (TRS)}
			\end{figure}
		\end{column}
	
		\begin{column}{.4\textwidth} % Left column and width
			\begin{figure}
				\only<2->{
					\includegraphics[height=5.0cm]{design_margins/DM_geometry/thermal_loadcase.pdf}
					\vspace{-0.2em}
					\caption{Thermal loads}
				}%
			\end{figure}
		\end{column}
		
	\end{columns}
	
\end{frame}
%------------------------------------------------
\subsection{Dynamically allocating margins}
%------------------------------------------------
\begin{frame}[t,label=ap2_4]
	\frametitle{Changing requirements}
	\vspace{-1em}
	\only<-10>{Let us consider 6 extremely different requirement scenarios\\}%
	\only<11>{The {\color{blue}blue} decisions exceed the required thresholds\\}%
	\only<12>{The {\color{red}red} decisions violate the thresholds but have less excess relative to the {\color{blue}blue} decisions\\}%
	\only<13>{The {\color{darkgreen}green} decisions have more reliability but a lot of excess\\}%
	\uncover<3-10>{Let us also consider 2 different design decisions}%
	\begin{columns}[t] % The "c" option specifies centered vertical alignment while the "t" option is used for top vertical alignment
		
		\begin{column}{.45\textwidth} % Left column and width
			\only<1-10>{
				\uncover<4->{
					$c=1$,$\mathbf{D}=\left[\uncover<4->{{\color<4>{red} 2}},%
					\uncover<5->{{\color<5>{red} 1}},%
					\uncover<6->{\only<-6>{{\color<6>{red} -1}}\only<7->{{\color<7>{red} 0}}},%
					\uncover<8->{{\color<8>{red} 4}},\uncover<9->{{\color<9>{red} -1}},%
					\uncover<10->{{\color<10>{red}3}}\right]$
				}
			}%
			\only<11->{
				\only<11>{$c=1$,$\mathbf{D}=\left[2, 1, 0, 4, -1, 3\right]$}%
				\only<12>{$c=1$,$\mathbf{D}=\left[2, 1, -1, -1, 0, -1\right]$}%
				\only<13>{$c=1$,$\mathbf{D}=\left[4, 1, 0, 2, -1, 3\right]$}%
			}%
			\uncover<2-10>{$\mathbf{R}=\left[{\color<4>{red} Req_{1}},%
										  {\color<5>{red} Req_{2}},%
										  {\color<6-7>{red} Req_{3}},%
										  {\color<8>{red} Req_{4}},%
										  {\color<9>{red} Req_{5}},%
										  {\color<10>{red} Req_{6}}\right]$}%
			% \vspace{-1em}
				\begin{figure}
						\only<1>{\includegraphics[height=5.0cm]{design_margins/Stagespace/2D_stagespace_res_th_only_no_line.pdf}}%	
						\only<2>{\includegraphics[height=5.0cm]{design_margins/Stagespace/2D_stagespace_res_th_only.pdf}}%	
						\only<3>{\includegraphics[height=5.0cm]{design_margins/Stagespace/D1_stagespace_res_1.pdf}}%
						\only<4>{\includegraphics[height=5.0cm]{design_margins/Stagespace/D1_stagespace_res_2.pdf}}%
						\only<5>{\includegraphics[height=5.0cm]{design_margins/Stagespace/D1_stagespace_res_3.pdf}}%
						\only<6>{\includegraphics[height=5.0cm]{design_margins/Stagespace/D1_stagespace_res_4.pdf}}%
						\only<7>{\includegraphics[height=5.0cm]{design_margins/Stagespace/D3_stagespace_res_4.pdf}}%
						\only<8>{\includegraphics[height=5.0cm]{design_margins/Stagespace/D3_stagespace_res_5.pdf}}%
						\only<9>{\includegraphics[height=5.0cm]{design_margins/Stagespace/D3_stagespace_res_6.pdf}}%
						\only<10>{\includegraphics[height=5.0cm]{design_margins/Stagespace/D3_stagespace_res_7.pdf}}%
						\only<11-12>{\includegraphics[height=5.0cm]{design_margins/Stagespace/2D_stagespace_res.pdf}}%
						\only<13>{\includegraphics[height=5.0cm]{design_margins/Stagespace/3D_stagespace_res.pdf}}%
					\label{fig:stagespace} 
					\vspace{-1.5em}
					\caption{Reliability at each epoch}
				\end{figure}
		\end{column}
	
		\begin{column}{.45\textwidth} % Left column and width
			\centering
			\vspace{-2.4em}
			\only<-10>{
					\color{white}{
					$\boldsymbol{\mu}=\left[ 0.5 ,~ 0.5 ,~ 0.5 ,~ 0.5 \right]$\\%
					$\boldsymbol{\sigma}=\left[ 0.1875375 ,~ 0.125 ,~ 0.125 ,~ 0.1875 \right]$%
				}
			}
			\only<11->{
				~\\%
				Total excess = \only<11>{4.58}\only<12>{{\color{darkgreen} 3.91 $\downarrow$}}%
							   \only<13>{{\color{red} 5.16 $\uparrow$}}\\%
				~\\%
				Reliability: \only<11>{{\color{darkgreen} Above threshold}}%
							 \only<12>{{\color{red} Below threshold}}%
							 \only<13>{{\color{darkgreen} Above threshold}}%
			}
			\begin{figure}
				\only<4>{\includegraphics[height=5.6cm]{design_margins/DM_parameter_space/146_reqs_2D/146_req_36_thermal_out_RS_2D_pi_6.pdf}}%
				\only<5>{\includegraphics[height=5.6cm]{design_margins/DM_parameter_space/163_reqs_2D/163_req_50_thermal_out_RS_2D_pi_6.pdf}}%
				\only<6>{\includegraphics[height=5.6cm]{design_margins/DM_parameter_space/163_reqs_2D/163_req_1_thermal_out_RS_2D_pi_6.pdf}}%
				\only<7>{\includegraphics[height=5.6cm]{design_margins/DM_parameter_space/164_reqs_2D/164_req_1_thermal_out_RS_2D_pi_6.pdf}}%
				\only<8>{\includegraphics[height=5.6cm]{design_margins/DM_parameter_space/167_reqs_2D/167_req_46_thermal_out_RS_2D_pi_6.pdf}}%
				\only<9>{\includegraphics[height=5.6cm]{design_margins/DM_parameter_space/167_reqs_2D/167_req_13_thermal_out_RS_2D_pi_6.pdf}}%
				\only<10>{\includegraphics[height=5.6cm]{design_margins/DM_parameter_space/168_reqs_2D/168_req_31_thermal_out_RS_2D_pi_6.pdf}}%
				\only<11>{\includegraphics[height=5.0cm]{design_margins/Stagespace/D3_stagespace_obj.png}}%
				\only<12>{\includegraphics[height=5.0cm]{design_margins/Stagespace/D1_stagespace_obj.png}}%
				\only<13>{\includegraphics[height=5.0cm]{design_margins/Stagespace/D2_stagespace_obj.png}}%
				\label{fig:stagespace_obj}
				\only<11->{
					\vspace{-1.5em}
					\caption{Excess at each epoch}
				}
			\end{figure}
		\end{column}
		
	\end{columns}
	
\end{frame}
%------------------------------------------------
\subsection{Optimizing when and how much change is needed}
%------------------------------------------------
\begin{frame}[t]
	\frametitle{Optimizing when and how much change is needed}
	\vspace{-0.5em}
	Formulate optimization problem to minimize total excess while being reliable
	
	\begin{columns}[t] % The "c" option specifies centered vertical alignment while the "t" option is used for top vertical alignment
		
		\begin{column}{.42\textwidth} % Left column and width
			\vspace{-1.5em}
			% Optimization problem
			\begin{alertblock}{Objective and constraints}
				\vspace{-1.5em}
				\begin{equation*}
					\begin{aligned}
						\uncover<2->{
							& {\text{minimize}}
							& & \text{excess}(c,\mathbf{D};\text{requirement arc})\\
						}
						\uncover<3->{
							& \text{subject to}
							& & reliability(c,\mathbf{D};\text{requirement arc})\\
						}	
						\uncover<4->{
							& \text{over}
							& & \text{The set of feasible design decisions}\\
							& & & \text{(404 possible combinations)}
						}
					\end{aligned}
				\end{equation*}
			\end{alertblock}
			\vspace{-0.5em}
			\uncover<4->{
				% Variables
				\begin{exampleblock}{Design variables}
					\vspace{-0.5em}
					\begin{equation*}
						\begin{aligned}
							& & & c \in \left\{"wavy","hatched","tubular"\right\}\\
							& & & \mathbf{D} = \left[\text{decision}_1,\text{decision}_2,\cdots\right]\\
							& & & \text{decision} \in \left\{-1,0,1,2,3,4\right\}\\
						\end{aligned}
					\end{equation*}
				\end{exampleblock}
			}
			\vspace{-0.5em}
			\uncover<6->{
				% Parameters
				\begin{block}{Requirement arcs}
					\vspace{-0.5em}
					\begin{equation*}
						\begin{aligned}
							& & & \text{We can randomly generate}\\
							& & & \text{10000 different requirement arcs}
						\end{aligned}
					\end{equation*}
				\end{block}
			}
		\end{column}

		\begin{column}{.5\textwidth} % Right column and width
			\vspace{-1em}
			\tikzstyle{background grid}=[draw, black!50,step=.5cm]
			
			\begin{tikzpicture}%[show background grid]
			% Put the graphic inside a node. This makes it easy to place the
			% graphic and to draw on top of it. 
			% The above right option is used to place the lower left corner
			% of the image at the (0,0) coordinate. 
				\uncover<-5>{
					\node [inner sep=0pt,above right, opacity=1.0]  at (0, 1) (objective)
						{
							\only<2->{\includegraphics[height=4.0cm]{design_margins/Stagespace/opt_stagespace_obj.png}}%
						};
					\node [inner sep=0pt,above right, opacity=0.7]  at (2.5, -1.75) (constraint)
						{
							\only<3>{\includegraphics[height=4.0cm]{design_margins/Stagespace/2D_stagespace_res_th_only_no_line.pdf}}%
							\only<4->{\includegraphics[height=4.0cm]{design_margins/Stagespace/opt_stagespace_res.pdf}}%
						};
				}
				\uncover<4->{
					\node[draw, align=center, fill=white, inner sep=2pt,above right, opacity=1]  at (0.35, 2) {
							Solving the problem yields optimal decisions\\% needs another set of braces
							{\uncover<5->{$\downarrow$}}\\%
							{\uncover<5->{optimal designs}}\\%
							{\uncover<5->{$\left\{c^{*},\mathbf{D}^{*}\right\}$}}%
						};
					\uncover<7->{
						\draw (0,1.5) rectangle (8,5) node at (4.,4.75) {Repeat for different requirement arcs};
					}
				}
				% show origin
				% \fill (0,0) circle (2pt);
				% define destination coordinates
			\end{tikzpicture}

		\end{column}

	\end{columns}
\end{frame}
%------------------------------------------------
\subsection{Set-based solution of parametric optimization problem}
%------------------------------------------------
\begin{frame}[t]
	\frametitle{Set-based solution of parametric optimization problem}
	\vspace{-0.5em}
	\begin{nscenter}
		\only<1>{We solve the problem for $\mathit{n = {\color{red} 100}}$ requirement arcs}%
		\only<2>{We solve the problem for $\mathit{n = {\color{red} 1000}}$ requirement arcs}%
		\only<3>{We solve the problem for $\mathit{n = {\color{red} 100000}}$ requirement arcs}%
	\end{nscenter}%

	\begin{columns}[t] % The "c" option specifies centered vertical alignment while the "t" option is used for top vertical alignment
		
		\begin{column}{.42\textwidth} % Left column and width
			\vspace{-1.5em}
			% Optimization problem
			\begin{alertblock}{Objective and constraints}
				\vspace{-1.5em}
				\begin{equation*}
					\begin{aligned}
							& {\text{minimize}}
							& & \text{excess}(c,\mathbf{D};\text{requirement arc})\\
							& \text{subject to}
							& & reliability(c,\mathbf{D};\text{requirement arc})\\
							& \text{over}
							& & \text{The set of feasible design decisions}\\
							& & & \text{(404 possible combinations)}
					\end{aligned}
				\end{equation*}
			\end{alertblock}
			\vspace{-0.5em}
			% Variables
			\begin{exampleblock}{Design variables}
				\vspace{-0.5em}
				\begin{equation*}
					\begin{aligned}
						& & & c \in \left\{"wavy","hatched","tubular"\right\}\\
						& & & \mathbf{D} = \left[\text{decision}_1,\text{decision}_2,\cdots\right]\\
						& & & \text{decision} \in \left\{-1,0,1,2,3,4\right\}\\
					\end{aligned}
				\end{equation*}
			\end{exampleblock}
			\vspace{-0.5em}
			% Parameters
			\begin{block}{Requirement arcs}
				\vspace{-0.5em}
				\begin{equation*}
					\begin{aligned}
						& & & \text{We can randomly generate}\\
						& & & \text{10000 different requirement arcs}
					\end{aligned}
				\end{equation*}
			\end{block}
		\end{column}

		\begin{column}[t]{.5\textwidth} % Right column and width

			\centering
			\tikzstyle{background grid}=[draw, black!50,step=.5cm]
			\begin{tikzpicture}%[show background grid]
			% Put the graphic inside a node. This makes it easy to place the
			% graphic and to draw on top of it. 
			% The above right option is used to place the lower left corner
			% of the image at the (0,0) coordinate. 
				\node [inner sep=0pt,above right, opacity=1] (histogram)
					{
						\only<1>{\includegraphics[height=5.0cm]{design_margins/Histograms/histogram_DOE_E_1.pdf}}%
						\only<2>{\includegraphics[height=5.0cm]{design_margins/Histograms/histogram_DOE_E_10.pdf}}%
						\only<3>{\includegraphics[height=5.0cm]{design_margins/Histograms/histogram_DOE_E_28.pdf}}%
					};
			\end{tikzpicture}

		\end{column}

	\end{columns}
\end{frame}
%------------------------------------------------
\subsection{Tradespace exploration}
%------------------------------------------------
\begin{frame}[t]
	\frametitle{Tradespace exploration}
	\vspace{-0.5em}
	On a tradespace, we plot the feasible designs, %
	\uncover<3->{the {\color{purple}Pareto front}}\uncover<4->{, and {\color{red}the optimal designs}}%

	\begin{columns}[t] % The "c" option specifies centered vertical alignment while the "t" option is used for top vertical alignment
		
		\begin{column}[t]{.42\textwidth} % Left column and width

			% Use a background grid to make it easier to find coordinates
			% When the coordinates have been found, remove the 
			% 'show background grid' option. 
			\tikzstyle{background grid}=[draw, black!50,step=.5cm]
			\begin{tikzpicture}%[show background grid]
				% Put the graphic inside a node. This makes it easy to place the
				% graphic and to draw on top of it. 
				% The above right option is used to place the lower left corner
				% of the image at the (0,0) coordinate.
				\uncover<2->{
					\node [inner sep=0pt,above right, opacity=1.0] (histogram)
						{
							\only<1-2>{\includegraphics[height=5.0cm]{design_margins/tradespace/tradespace_pareto_D.pdf}}%
							\only<3>{\includegraphics[height=5.0cm]{design_margins/tradespace/tradespace_pareto_DP.pdf}}%
							\only<4>{\includegraphics[height=5.0cm]{design_margins/tradespace/tradespace_pareto_DPS.pdf}}%
							\only<5>{\includegraphics[height=5.0cm]{design_margins/tradespace/tradespace_pareto_PS.pdf}}%
						};
					\node[inner sep=0pt,align=flush center,above=\belowcaptionskip of histogram,text width=\linewidth]
						{\vspace{-1em}{\large Tradespace}};
					% show origin
					%\fill (0,0) circle (2pt);
				}%

			\end{tikzpicture}


		\end{column}

		\begin{column}[t]{.5\textwidth} % Right column and width
			% Use a background grid to make it easier to find coordinates
			% When the coordinates have been found, remove the 
			% 'show background grid' option. 
			\tikzstyle{background grid}=[draw, black!50,step=.5cm]
			\begin{tikzpicture}%[show background grid]
				% Put the graphic inside a node. This makes it easy to place the
				% graphic and to draw on top of it. 
				% The above right option is used to place the lower left corner
				% of the image at the (0,0) coordinate. 
				\node [inner sep=0pt,above right, opacity=1.0] (histogram)
					{
						\includegraphics[height=5.0cm]{design_margins/Histograms/histogram_DOE_E_28.pdf}%
					};
				\node[inner sep=0pt,align=flush center,above=\belowcaptionskip of histogram,text width=\linewidth]
					{\vspace{-1em}{\large 
						Histogram (optimal designs)
					}};
				% show origin
				%\fill (0,0) circle (2pt);
				
			\end{tikzpicture}

		\end{column}

	\end{columns}

\end{frame}
%------------------------------------------------
\begin{frame}[t]
	\frametitle{Tradespace exploration: Robustness}
	\vspace{-0.5em}
	We can check the percentage of the 10000 requirements arcs satisfied by each feasible design decision to identify the most {\color{blue} Robust solutions}
	\vspace{-1.0em}
	\begin{columns}[t] % The "c" option specifies centered vertical alignment while the "t" option is used for top vertical alignment
		
		\begin{column}[t]{.42\textwidth} % Left column and width

			% Use a background grid to make it easier to find coordinates
			% When the coordinates have been found, remove the 
			% 'show background grid' option. 
			\tikzstyle{background grid}=[draw, black!50,step=.5cm]
			\begin{tikzpicture}%[show background grid]
				% Put the graphic inside a node. This makes it easy to place the
				% graphic and to draw on top of it. 
				% The above right option is used to place the lower left corner
				% of the image at the (0,0) coordinate.
				\node [inner sep=0pt,above right, opacity=1.0] (histogram)
					{
						\includegraphics[height=5.0cm]{design_margins/tradespace/tradespace_pareto_PSR.pdf}%
					};
				\node[inner sep=0pt,align=flush center,above=\belowcaptionskip of histogram,text width=\linewidth]
					{\vspace{-1em}{\large Tradespace}};
				% show origin
				%\fill (0,0) circle (2pt);
			\end{tikzpicture}


		\end{column}

		\begin{column}[t]{.5\textwidth} % Right column and width
			% Use a background grid to make it easier to find coordinates
			% When the coordinates have been found, remove the 
			% 'show background grid' option. 
			\tikzstyle{background grid}=[draw, black!50,step=.5cm]
			\begin{tikzpicture}%[show background grid]
				% Put the graphic inside a node. This makes it easy to place the
				% graphic and to draw on top of it. 
				% The above right option is used to place the lower left corner
				% of the image at the (0,0) coordinate. 
				\node [inner sep=0pt,above right, opacity=1.0] (histogram)
					{
						\includegraphics[height=5.0cm]{design_margins/Histograms/histogram_DOE_R.pdf}%
					};
				\node[inner sep=0pt,align=flush center,above=\belowcaptionskip of histogram,text width=\linewidth]
					{\vspace{-1em}{\large 
						Histogram (robust designs)
					}};
				% show origin
				%\fill (0,0) circle (2pt);
				
			\end{tikzpicture}

		\end{column}

	\end{columns}

\end{frame}
%------------------------------------------------
\begin{frame}[t]
	\frametitle{Tradespace exploration: Flexibility}
	\vspace{-0.5em}
	The filtered outdegree is the number of design alternatives offered by each decision

	\begin{columns}[t] % The "c" option specifies centered vertical alignment while the "t" option is used for top vertical alignment
		
		\begin{column}[t]{.42\textwidth} % Left column and width

			% Use a background grid to make it easier to find coordinates
			% When the coordinates have been found, remove the 
			% 'show background grid' option. 
			\tikzstyle{background grid}=[draw, black!50,step=.5cm]
			\begin{tikzpicture}%[show background grid]
				% Put the graphic inside a node. This makes it easy to place the
				% graphic and to draw on top of it. 
				% The above right option is used to place the lower left corner
				% of the image at the (0,0) coordinate.
				\node [inner sep=0pt,above right, opacity=1.0] (histogram)
					{
						\includegraphics[height=5.0cm]{design_margins/tradespace/tradespace_pareto_PSRF.pdf}%
					};
				\node[inner sep=0pt,align=flush center,above=\belowcaptionskip of histogram,text width=\linewidth]
					{\vspace{-1em}{\large Tradespace}};
				% show origin
				%\fill (0,0) circle (2pt);
			\end{tikzpicture}


		\end{column}

		\begin{column}[t]{.5\textwidth} % Right column and width
			% Use a background grid to make it easier to find coordinates
			% When the coordinates have been found, remove the 
			% 'show background grid' option. 
				\tikzstyle{background grid}=[draw, black!50,step=.5cm]
				\begin{tikzpicture}%[show background grid]
					% Put the graphic inside a node. This makes it easy to place the
					% graphic and to draw on top of it. 
					% The above right option is used to place the lower left corner
					% of the image at the (0,0) coordinate. 
					\node [inner sep=0pt,above right, opacity=1.0] (histogram)
						{
							\includegraphics[height=5.0cm]{design_margins/Histograms/histogram_DOE_F.pdf}%
						};
					\node[inner sep=0pt,align=flush center,above=\belowcaptionskip of histogram,text width=\linewidth]
						{\vspace{-1em}{\large 
							Histogram (flexible designs)
						}};
					% show origin
					%\fill (0,0) circle (2pt);
					
				\end{tikzpicture}

		\end{column}

	\end{columns}

\end{frame}

%------------------------------------------------
\begin{frame}[t]
	\frametitle{Tradespace exploration: Comparison}
	\vspace{-0.5em}
	We use \textbf{convex hulls} to draw some conclusions about the obtained solution sets

	\begin{columns}[t] % The "c" option specifies centered vertical alignment while the "t" option is used for top vertical alignment
		
		\begin{column}[t]{.42\textwidth} % Left column and width

			% Use a background grid to make it easier to find coordinates
			% When the coordinates have been found, remove the 
			% 'show background grid' option. 
			\tikzstyle{background grid}=[draw, black!50,step=.5cm]
			\begin{tikzpicture}%[show background grid]
				% Put the graphic inside a node. This makes it easy to place the
				% graphic and to draw on top of it. 
				% The above right option is used to place the lower left corner
				% of the image at the (0,0) coordinate.
				\node [inner sep=0pt,above right, opacity=1.0] (histogram)
					{
						\only<1>{\includegraphics[height=5.0cm]{design_margins/tradespace/tradespace_pareto_PSRF_CH.pdf}}%
						\only<2->{\includegraphics[height=5.0cm]{design_margins/tradespace/tradespace_pareto_PSRF_CH_NP.pdf}}%
					};
				\node[inner sep=0pt,align=flush center,above=\belowcaptionskip of histogram,text width=\linewidth]
					{\vspace{-1em}{\large Tradespace}};
				% show origin
				%\fill (0,0) circle (2pt);

				\path (2.49,2.875) coordinate (flexible)
						(3.95,4.2) coordinate (SBD)
						(5.4,4.75) coordinate (robust);

			\end{tikzpicture}


		\end{column}

		\begin{column}[t]{.5\textwidth} % Right column and width
			% Use a background grid to make it easier to find coordinates
			% When the coordinates have been found, remove the 
			% 'show background grid' option. 
			\begin{itemize}
				\uncover<3->{\item \tikz[na] \coordinate (robust_bullet); The \textbf{\color{blue} robust} solutions favor {\color{red} heavier $\uparrow$} designs with {\color{darkgreen} high capabilities $\uparrow$}}\\
				~\\
				\uncover<5->{\item \tikz[na] \coordinate (SBD_bullet); The \textbf{\color{red} optimized} set-based solutions balance weight with capability}\\
				~\\
				\uncover<4->{\item \tikz[na] \coordinate (flexible_bullet); The \textbf{\color{darkgreen} flexible} solutions favor {\color{darkgreen} light $\downarrow$} designs with {\color{red} low capabilities $\downarrow$}}
			\end{itemize}

		\end{column}

	\end{columns}

	% define overlays
	% Note the use of the overlay option. This is required when 
	% you want to access nodes in different pictures.
	\begin{tikzpicture}[overlay]
		\only<3->{\path[->,magenta,thick] (robust) edge [out=0, in=180] (robust_bullet);}
		\only<4->{\path[->,magenta,thick] (flexible) edge [out=0, in=180] (flexible_bullet);}
		\only<5->{\path[->,magenta,thick] (SBD) edge [out=0, in=180] (SBD_bullet);}
	\end{tikzpicture}

\end{frame}
% \include{abm_work}
% \include{machine_learning}
% 
%==============================================================================%
%% Mathematical symbols
\definecolor{design}{HTML}{D0A31D}
\definecolor{spec}{HTML}{FF7575}
\definecolor{margin}{HTML}{CC99FF}
\definecolor{perf}{HTML}{87C061}
\newcommand\design[1]{{\color{design}#1}}
\newcommand\spec[1]{{\color{spec}#1}}
\newcommand\margin[1]{{\color{margin}#1}}
\newcommand\perf[1]{{\color{perf}#1}}
%==============================================================================%
%%%%%%%%%%%%%%%%%%%%%%%%%%%%%%%%%%%%%%%%%%%%%%%%%%%%%%%
%%              Supplementary Content                %%
%%%%%%%%%%%%%%%%%%%%%%%%%%%%%%%%%%%%%%%%%%%%%%%%%%%%%%%
%------------------------------------------------
\section{Supplementary content}
%------------------------------------------------
\begin{frame}[c,noframenumbering]
	\centering
	% \setlength\fboxsep{0pt}
	\begin{titleblock}{}
		~\\%
		{\centering\LARGE Supplementary content: Miscellaneous\\}%
		~\\%
	\end{titleblock}
\end{frame}
%------------------------------------------------
\subsection{Example problem: Design of TRS strut}
%------------------------------------------------
\begin{frame}[t,label=ap2_2]
	\frametitle{Overview of turbine rear frame design}
	\vspace{-1em}
	\only<1-2>{The turbine rear structure (TRS) sits aft of the turbine}%
	\only<3>{We consider the design of a single strut}%
	\only<4>{Variable thermal loads are expected}%
	\only<5-7>{Which cause thermal expansion \uncover<6-7>{and stress on the strut}}%
	\only<8>{The entire problem can be represented compactly as an MDM}%
	~\\
	\begin{columns}[t] % The "c" option specifies centered vertical alignment while the "t" option is used for top vertical alignment
			
		\begin{column}{.45\textwidth} % Left column and width
			\footnotesize
			Design parameters:
			\design{
				\uncover<3->{
				\begin{itemize}\itemsep0em 
					\item \only<8>{$D1$} vane width $w$
					\item \only<8>{$D2$} vane height $h$
					\item \only<8>{$D3$} lean angle $\theta$
					\item \only<8>{$D4$} Material - yield stress $\sigma_y$ (off-the shelf)
				\end{itemize}
			}}%
			Input Specifications:
			\spec{
				\uncover<5->{
				\begin{itemize}\itemsep0em 
					\item \only<8>{$S1$} Nacelle temperature $T_1$ ($\downarrow$)
					\item \only<8>{$S2$} Gas surface temperature $T_2$ ($\uparrow$)
				\end{itemize}
			}}%
			Design margins:
			\margin{
				\uncover<6->{
				\begin{itemize}\itemsep0em 
					\item \only<8>{$E1$} compression $\sigma_y - \sigma_a \ge 0$
					\item \only<8>{$E2$} bending $\sigma_y - \sigma_m \ge 0$
					\item \only<8>{$E3$} buckling $F_\text{buckling} - F_a \ge 0$
				\end{itemize}
			}}%
			Performance parameters:
			\perf{
				\uncover<7->{
				\begin{itemize}\itemsep0em 
					\item \only<8>{$P1$}  Weight
				\end{itemize}
			}}%
		\end{column}
	
		\begin{column}{.4\textwidth} % Left column and width
			\begin{figure}
				\only<1>{
					\includegraphics[width=6.0cm]{design_space_exploration/TRS/engine.png}
					\vspace{-0.2em}
					\caption{Trent 900 aeroengine}
				}%
				\only<2>{
					\includegraphics[width=5.0cm]{design_space_exploration/TRS/rear_view.pdf}
					\vspace{-0.2em}
					\caption{Turbine rear structure}
				}%
				\only<3>{
					\includegraphics[width=7.0cm]{design_space_exploration/TRS/cross_section_view.pdf}
					\vspace{-0.2em}
					\caption{Turbine rear structure: cross-sectional view}
				}%
				\only<4>{
					\includegraphics[width=6.0cm]{design_space_exploration/TRS/expansion_1.pdf}
					\vspace{-0.2em}
					\caption{Thermal expansion}
				}%
				\only<5>{
					\includegraphics[width=6.0cm]{design_space_exploration/TRS/expansion_2.pdf}
					\vspace{-0.2em}
					\caption{Thermal expansion}
				}%
				\only<6-7>{
					\includegraphics[width=7.0cm]{design_space_exploration/TRS/loads.pdf}
					\vspace{-0.2em}
					\caption{Load cases}
				}%
			\end{figure}
		\end{column}
		
	\end{columns}
	
\end{frame}
%------------------------------------------------
\subsection{Implementation}
%------------------------------------------------
\begin{frame}[fragile,t]
	\frametitle{Pseudo-code of MVM}
    %
	\vspace{-1em}
	\begin{algorithm}[H] % must use [H]
		\tiny
		\DontPrintSemicolon % Some LaTeX compilers require you to use \dontprintsemicolon instead
		\KwIn{
			set of possible choices $\mathcal{D}$, chosen design $\mathbf{c}\in\mathcal{D}$,
			the MAN given by analytic functions, models, and surrogates,
			$n_\mathrm{samples}$, Joint PDF $F_\mathbf{S}(\mathbf{s})$
		}
		\KwOut{estimated mean impact and change absorption $\bar{\bar{\mathbf{a}}}$, $\bar{\bar{\mathbf{i}}}$}
		%
		\textbf{[0] Initialization}: 
			set counter $i_\mathrm{sample} \gets 0$, initialize aggregate absorption and impact $\bar{\bar{\mathbf{a}}} \in \mathbb{R}^e$ and $\bar{\bar{\mathbf{i}}} \in \mathbb{R}^e$\;
		%
		\textbf{[1] Monte-carlo simulation}: 
			draw sample input specifications $\mathbf{S}\in\mathbb{R}^n \sim F_\mathbf{S}(\mathbf{s})$\;
		%
		\textbf{[2] Calculate threshold, decided, excess, and performance values}: 
			$\mathbf{t} = \mathbf{f}(\mathbf{S},\mathbf{c})$, 
			$\mathbf{d} = \mathbf{g}(\mathbf{c})$, 
			$\mathbf{e} = \left(\mathbf{t} - \mathbf{d}\right){\mathbf{i}^\mathrm{excess}}^\mathsf{T}$, 
			$\mathbf{p} =  \mathbf{h}(\mathbf{S},\mathbf{c})$\;
		%
		\textbf{[3] Calculate allowable deterioration}\;
		\qquad loop over input specifications \For{$k = 1, 2, ..., n$} {
			\qquad perform a line search along the $k$th direction to identify the root of: 
				$\min{((\mathbf{f}(\mathbf{x},\mathbf{c}) - \mathbf{d}){\mathbf{i}^\mathrm{excess}}^\mathsf{T})} = \mathbf{0}$, $s_k' \gets x^*$, 
				$s_k^\mathrm{limit} \gets s_k'$\;
		}
		%
		\textbf{[4] Calculate margin value metrics}\;
		\qquad initialization tensors and vectors, 
			$\mathbf{I} \in \mathbb{R}^{m\times e} \gets \mathbf{0}$ and $\mathbf{A} \in \mathbb{R}^{n\times e} \gets \mathbf{0}$, 
			$\mathbf{v} \in \mathbb{R}^{n} \gets \mathbf{0}$, $\bar{\mathbf{a}} \in \mathbb{R}^{e} \gets \mathbf{0}$, and $\bar{\mathbf{i}} \in \mathbb{R}^{e} \gets \mathbf{0}$\;
		\qquad loop over margin nodes \For{$i = 1, 2, ..., e$} {  
			\qquad \textbf{[a] Calculate impact on performance}\;
			\qquad $\mathbf{d}^\mathrm{threshold} \gets \mathbf{d}$, $d_i^\mathrm{threshold} \gets t_i$, 
				$\mathbf{c}_i^\mathrm{threshold} = \mathbf{g}^{-1}(\mathbf{d}^\mathrm{threshold})$, 
				$\mathbf{p}_i^\mathrm{threshold} = \mathbf{h}(\mathbf{S},\mathbf{c}_i^\mathrm{threshold})$\;
			\qquad loop over performance parameters \For{$j = 1, 2, ..., m$} {
				\qquad $i_{i,j} \gets (p_j - p_{j,i}^\mathrm{threshold})i_j^p/{p_{j,i}^\mathrm{threshold}}$\;
			}
			\qquad \textbf{[b] Calculate change absorption capability}\;
			\qquad loop over input specifications \For{$k = 1, 2, ..., n$} {
				\qquad calculate deterioration $v_k=(s_k^\mathrm{limit} - s_k)/{s_k}$\;
				\qquad $a_{i,k} \gets (f_i(\mathbf{s}_k^\mathrm{limit},\mathbf{c}) - f_i(\mathbf{S},\mathbf{c}))/(f_i(\mathbf{S},\mathbf{c})v_k)$\;
			}
			\qquad \textbf{[c] Aggregate metrics}: 
				$\bar{a}_i \gets \frac{1}{n}\sum_{k=1}^{n} a_{i,k}$, 
				$\bar{i}_i \gets \frac{1}{m}\sum_{k=1}^{m} i_{i,j}$
		}
		%
		\textbf{[5] sum of monte-carlo samples}: 
			$\bar{\bar{\mathbf{a}}} \gets \bar{\bar{\mathbf{a}}} + \bar{\mathbf{a}}$, 
			$\bar{\bar{\mathbf{i}}} \gets \bar{\bar{\mathbf{i}}} + \bar{\mathbf{i}}$\;
		%
		\textbf{[6] Termination}, 
		\If{$i_\mathrm{sample} < n_\mathrm{samples}$} {
			set $i_\mathrm{sample} \gets i_\mathrm{sample}+1$, go to \textbf{[1]}\;
			otherwise $\bar{\bar{\mathbf{a}}} \gets \bar{\bar{\mathbf{a}}}/n_\mathrm{samples}$ and $\bar{\bar{\mathbf{i}}} \gets \bar{\bar{\mathbf{i}}}/n_\mathrm{samples}$, \textbf{stop}
		}
	\end{algorithm}
\end{frame}
%------------------------------------------------
\begin{frame}[t,noframenumbering]
	\frametitle{UML diagram of \texttt{mvmlib}\ifshowcitations\footpartcite{Alhandawi2022}\fi}
    %
	\begin{figure}
        \centering
		\only<1>{\includegraphics[width=1.0\textwidth]{design_space_exploration/classes_mvmlib.pdf}}%
		\only<2>{\includegraphics[width=0.6\textwidth]{design_space_exploration/packages_mvmlib.pdf}}%
    \end{figure}%
\end{frame}
\addtocounter{footnote}{-1}
%------------------------------------------------
\subsection{Overview of mesh adaptive direct search}
%------------------------------------------------
\begin{frame}[t,label=abm_5]
	\frametitle{Overview of mesh adaptive direct search}
	\tikzstyle{background grid}=[draw, black!50,step=.5cm]
	%
	\uncover<1->{No gradient information available, blackbox is expensive\ifshowcitations\footpartcite{Audet2017}\fi}\\
	%
	\begin{columns}[t] % The "c" option specifies centered vertical alignment while the "t" option is used for top vertical alignment
		\begin{column}{.42\textwidth} % Left column and width
			\vspace{-1.2em}
			% Optimization problem
			\begin{exampleblock}{Objective and constraints}
				\begin{equation*}
					\begin{aligned}
						& \underset{\mathbf{x}}{\text{min}}
						& & f(\mathbf{x})\\
						& \text{subject to}
						& & {c}(\mathbf{x}) \le 0\\
						& \text{where}
						& & \mathbf{x}:\text{ variables}\\
					\end{aligned}
				\end{equation*}
			\end{exampleblock}

			\small
			\begin{itemize}\itemsep0em
				\item SEARCH step and POLL step
				\item<2-> we focus on the POLL step
				\item<3-> incumbent solution: $x_k$
				\item<4-> trail points: $\left\{p_1,p_2,\cdots,p_k\right\}$
				\item<10-> successful trail point $p_7$
			\end{itemize}
		\end{column}
		%
		\begin{column}{.5\textwidth} % Left column and width
			\tikzstyle{background grid}=[draw, black!50,step=.5cm]
			
			\only<2>{\def\frameopt{-2}}%
			\only<3>{\def\frameopt{-1}}%
			\only<4>{\def\frameopt{0}}%
			\only<5>{\def\frameopt{1}}%
			\only<6>{\def\frameopt{2}}%
			\only<7>{\def\frameopt{3}}%
			\only<8>{\def\frameopt{4}}%
			\only<9>{\def\frameopt{5}}%
			\only<10>{\def\frameopt{6}}%
			\only<11>{\def\frameopt{7}}%
			\only<12>{\def\frameopt{8}}%
			\only<13>{\def\frameopt{9}}%

			\tikzstyle{background grid}=[draw, black!50,step=.1cm]
			\hspace*{5em}\raisebox{-1em}{%
				\begin{tikzpicture}[scale=0.45, every node/.style={scale=0.45}, remember picture, overlay] %show background grid, 
					\only<2->{\ifnum\frameopt < 2
    \def\unit{5}
    \def\meshsize{2}
    \def\pollX{0}
    \def\pollY{0}
    \def\pollsize{2}
    %
    \def\pAX{-1}
    \def\pAY{\pollsize/2}
    \def\pBX{\pollsize/2}
    \def\pBY{0}
    \def\pCX{0}
    \def\pCY{-\pollsize/2}
\fi
%
\ifnum\frameopt=-2
    % https://tex.stackexchange.com/a/183633
\newcommand{\Grid}[6]{%
    \def\startX{#1}
    \def\startY{#2}
    \def\maxX{#3}
    \def\maxY{#4}
    \def\size{#5}
    \def\borderthick{#6}
    %
    \draw[line width=\borderthick] (\size*0+\startX,\size*0+\startY) rectangle (\size*\maxX+\startX,\size*-\maxY+\startY);
    \foreach \x in {0,1,...,\maxX}{
        \draw (\size*\x+\startX,\size*0+\startY) -- (\size*\x+\startX,\size*-\maxY+\startY);
    }
    %
    \foreach \y in {0,1,...,\maxY}{
        \draw (\size*0+\startX,\size*-\y+\startY) -- (\size*\maxX+\startX,\size*-\y+\startY);
    }
}
%
\pgfmathsetmacro{\centerX}{\meshsize/2}
\pgfmathsetmacro{\centerY}{-\meshsize/2}
%
\pgfmathsetmacro{\pollcenterX}{\centerX+\pollX}
\pgfmathsetmacro{\pollcenterY}{\centerY+\pollY}
\pgfmathsetmacro{\pollcornerX}{\pollcenterX - \pollsize/2}
\pgfmathsetmacro{\pollcornerY}{\pollcenterY + \pollsize/2}
%
\pgfmathsetmacro{\pAcoordX}{\pollcenterX + \pAX}
\pgfmathsetmacro{\pAcoordY}{\pollcenterY - \pAY}
\pgfmathsetmacro{\pBcoordX}{\pollcenterX + \pBX}
\pgfmathsetmacro{\pBcoordY}{\pollcenterY - \pBY}
\pgfmathsetmacro{\pCcoordX}{\pollcenterX + \pCX}
\pgfmathsetmacro{\pCcoordY}{\pollcenterY - \pCY}
%
\pgfmathsetmacro{\deltameshn}{\pollcenterX-\pollcornerX}
\pgfmathsetmacro{\deltameshd}{\centerX}
\pgfmathsetmacro{\deltamesh}{(\pollcenterX-\pollcornerX)/\centerX}
\pgfmathsetmacro{\deltaframen}{\deltameshn * \deltameshn}
\pgfmathsetmacro{\deltaframed}{\deltameshd * \deltameshd}
\pgfmathsetmacro{\deltaframe}{\deltamesh * \deltamesh}
%
\Grid{0}{0}{\meshsize}{\meshsize}{\unit}{1pt}
\Grid{\pollcornerX*\unit}{\pollcornerY*\unit}{\pollsize}{\pollsize}{\unit}{3pt}
%
\ifdim \deltamesh pt = 1 pt
    \dimline[label style={above=0.5ex,}, line style = {line width=0.7},extension start length=+0.3,extension end length=+0.3] {(\pollcornerX*\unit,1)}{(\pollcenterX*\unit,1)}{\huge $\delta^{k}_\mathrm{poll}=1$};
    
    \dimline[label style={above=0.5ex,}, line style = {line width=0.7},extension start length=+0.3,extension end length=+0.3] {(\meshsize*\unit - 1*\unit,1)}{(\meshsize*\unit,1)}{\huge $\delta^{k}_\mathrm{mesh}=1$};
\else
    \dimline[label style={above=0.5ex,}, line style = {line width=0.7},extension start length=+0.3,extension end length=+0.3] {(\pollcornerX*\unit,0.5)}{(\pollcenterX*\unit,0.5)}{\huge $\delta^{k}_\mathrm{poll}=\addfraction{\deltameshn}{\deltameshd}$};
    
    \dimline[label style={above=0.5ex,}, line style = {line width=0.7},extension start length=+0.3,extension end length=+0.3] {(\meshsize*\unit - 1*\unit,0.5)}{(\meshsize*\unit,0.5)}{\huge $\delta^{k}_\mathrm{mesh}=\addfraction{\deltaframen}{\deltaframed}$};
\fi
\fi
%
\ifnum\frameopt=-1
    % https://tex.stackexchange.com/a/183633
\newcommand{\Grid}[6]{%
    \def\startX{#1}
    \def\startY{#2}
    \def\maxX{#3}
    \def\maxY{#4}
    \def\size{#5}
    \def\borderthick{#6}
    %
    \draw[line width=\borderthick] (\size*0+\startX,\size*0+\startY) rectangle (\size*\maxX+\startX,\size*-\maxY+\startY);
    \foreach \x in {0,1,...,\maxX}{
        \draw (\size*\x+\startX,\size*0+\startY) -- (\size*\x+\startX,\size*-\maxY+\startY);
    }
    %
    \foreach \y in {0,1,...,\maxY}{
        \draw (\size*0+\startX,\size*-\y+\startY) -- (\size*\maxX+\startX,\size*-\y+\startY);
    }
}
%
\pgfmathsetmacro{\centerX}{\meshsize/2}
\pgfmathsetmacro{\centerY}{-\meshsize/2}
%
\pgfmathsetmacro{\pollcenterX}{\centerX+\pollX}
\pgfmathsetmacro{\pollcenterY}{\centerY+\pollY}
\pgfmathsetmacro{\pollcornerX}{\pollcenterX - \pollsize/2}
\pgfmathsetmacro{\pollcornerY}{\pollcenterY + \pollsize/2}
%
\pgfmathsetmacro{\pAcoordX}{\pollcenterX + \pAX}
\pgfmathsetmacro{\pAcoordY}{\pollcenterY - \pAY}
\pgfmathsetmacro{\pBcoordX}{\pollcenterX + \pBX}
\pgfmathsetmacro{\pBcoordY}{\pollcenterY - \pBY}
\pgfmathsetmacro{\pCcoordX}{\pollcenterX + \pCX}
\pgfmathsetmacro{\pCcoordY}{\pollcenterY - \pCY}
%
\pgfmathsetmacro{\deltameshn}{\pollcenterX-\pollcornerX}
\pgfmathsetmacro{\deltameshd}{\centerX}
\pgfmathsetmacro{\deltamesh}{(\pollcenterX-\pollcornerX)/\centerX}
\pgfmathsetmacro{\deltaframen}{\deltameshn * \deltameshn}
\pgfmathsetmacro{\deltaframed}{\deltameshd * \deltameshd}
\pgfmathsetmacro{\deltaframe}{\deltamesh * \deltamesh}
%
\Grid{0}{0}{\meshsize}{\meshsize}{\unit}{1pt}
\Grid{\pollcornerX*\unit}{\pollcornerY*\unit}{\pollsize}{\pollsize}{\unit}{3pt}
%
\ifdim \deltamesh pt = 1 pt
    \dimline[label style={above=0.5ex,}, line style = {line width=0.7},extension start length=+0.3,extension end length=+0.3] {(\pollcornerX*\unit,1)}{(\pollcenterX*\unit,1)}{\huge $\delta^{k}_\mathrm{poll}=1$};
    
    \dimline[label style={above=0.5ex,}, line style = {line width=0.7},extension start length=+0.3,extension end length=+0.3] {(\meshsize*\unit - 1*\unit,1)}{(\meshsize*\unit,1)}{\huge $\delta^{k}_\mathrm{mesh}=1$};
\else
    \dimline[label style={above=0.5ex,}, line style = {line width=0.7},extension start length=+0.3,extension end length=+0.3] {(\pollcornerX*\unit,0.5)}{(\pollcenterX*\unit,0.5)}{\huge $\delta^{k}_\mathrm{poll}=\addfraction{\deltameshn}{\deltameshd}$};
    
    \dimline[label style={above=0.5ex,}, line style = {line width=0.7},extension start length=+0.3,extension end length=+0.3] {(\meshsize*\unit - 1*\unit,0.5)}{(\meshsize*\unit,0.5)}{\huge $\delta^{k}_\mathrm{mesh}=\addfraction{\deltaframen}{\deltaframed}$};
\fi
    \node[circle, minimum size = 1mm, fill=red, label={[fill=white,above right,label distance=0.2cm]\huge $x_k$}] (center) at (\pollcenterX*\unit,\pollcenterY*\unit) {};
\fi
%
\ifnum\frameopt=0
    % https://tex.stackexchange.com/a/183633
\newcommand{\Grid}[6]{%
    \def\startX{#1}
    \def\startY{#2}
    \def\maxX{#3}
    \def\maxY{#4}
    \def\size{#5}
    \def\borderthick{#6}
    %
    \draw[line width=\borderthick] (\size*0+\startX,\size*0+\startY) rectangle (\size*\maxX+\startX,\size*-\maxY+\startY);
    \foreach \x in {0,1,...,\maxX}{
        \draw (\size*\x+\startX,\size*0+\startY) -- (\size*\x+\startX,\size*-\maxY+\startY);
    }
    %
    \foreach \y in {0,1,...,\maxY}{
        \draw (\size*0+\startX,\size*-\y+\startY) -- (\size*\maxX+\startX,\size*-\y+\startY);
    }
}
%
\pgfmathsetmacro{\centerX}{\meshsize/2}
\pgfmathsetmacro{\centerY}{-\meshsize/2}
%
\pgfmathsetmacro{\pollcenterX}{\centerX+\pollX}
\pgfmathsetmacro{\pollcenterY}{\centerY+\pollY}
\pgfmathsetmacro{\pollcornerX}{\pollcenterX - \pollsize/2}
\pgfmathsetmacro{\pollcornerY}{\pollcenterY + \pollsize/2}
%
\pgfmathsetmacro{\pAcoordX}{\pollcenterX + \pAX}
\pgfmathsetmacro{\pAcoordY}{\pollcenterY - \pAY}
\pgfmathsetmacro{\pBcoordX}{\pollcenterX + \pBX}
\pgfmathsetmacro{\pBcoordY}{\pollcenterY - \pBY}
\pgfmathsetmacro{\pCcoordX}{\pollcenterX + \pCX}
\pgfmathsetmacro{\pCcoordY}{\pollcenterY - \pCY}
%
\pgfmathsetmacro{\deltameshn}{\pollcenterX-\pollcornerX}
\pgfmathsetmacro{\deltameshd}{\centerX}
\pgfmathsetmacro{\deltamesh}{(\pollcenterX-\pollcornerX)/\centerX}
\pgfmathsetmacro{\deltaframen}{\deltameshn * \deltameshn}
\pgfmathsetmacro{\deltaframed}{\deltameshd * \deltameshd}
\pgfmathsetmacro{\deltaframe}{\deltamesh * \deltamesh}
%
\Grid{0}{0}{\meshsize}{\meshsize}{\unit}{1pt}
\Grid{\pollcornerX*\unit}{\pollcornerY*\unit}{\pollsize}{\pollsize}{\unit}{3pt}
%
\ifdim \deltamesh pt = 1 pt
    \dimline[label style={above=0.5ex,}, line style = {line width=0.7},extension start length=+0.3,extension end length=+0.3] {(\pollcornerX*\unit,1)}{(\pollcenterX*\unit,1)}{\huge $\delta^{k}_\mathrm{poll}=1$};
    
    \dimline[label style={above=0.5ex,}, line style = {line width=0.7},extension start length=+0.3,extension end length=+0.3] {(\meshsize*\unit - 1*\unit,1)}{(\meshsize*\unit,1)}{\huge $\delta^{k}_\mathrm{mesh}=1$};
\else
    \dimline[label style={above=0.5ex,}, line style = {line width=0.7},extension start length=+0.3,extension end length=+0.3] {(\pollcornerX*\unit,0.5)}{(\pollcenterX*\unit,0.5)}{\huge $\delta^{k}_\mathrm{poll}=\addfraction{\deltameshn}{\deltameshd}$};
    
    \dimline[label style={above=0.5ex,}, line style = {line width=0.7},extension start length=+0.3,extension end length=+0.3] {(\meshsize*\unit - 1*\unit,0.5)}{(\meshsize*\unit,0.5)}{\huge $\delta^{k}_\mathrm{mesh}=\addfraction{\deltaframen}{\deltaframed}$};
\fi
    \node[circle, minimum size = 1mm, fill=red, label={[fill=white,above right,label distance=0.2cm]\huge $x_k$}] (center) at (\pollcenterX*\unit,\pollcenterY*\unit) {};
    \node[circle, minimum size = 1mm, fill=red, label={[fill=white,below,label distance=-0.5cm]\huge $p_1$}] (p1) at (\pAcoordX*\unit,\pAcoordY*\unit) {};
    \node[circle, minimum size = 1mm, fill=red, label={[fill=white,right,label distance=0.5cm]\huge $p_2$}] (p2) at (\pBcoordX*\unit,\pBcoordY*\unit) {};
    \node[circle, minimum size = 1mm, fill=red, label={[fill=white]\huge $p_3$}] (p3) at (\pCcoordX*\unit,\pCcoordY*\unit) {};
\fi
%
\ifnum\frameopt=1
    % https://tex.stackexchange.com/a/183633
\newcommand{\Grid}[6]{%
    \def\startX{#1}
    \def\startY{#2}
    \def\maxX{#3}
    \def\maxY{#4}
    \def\size{#5}
    \def\borderthick{#6}
    %
    \draw[line width=\borderthick] (\size*0+\startX,\size*0+\startY) rectangle (\size*\maxX+\startX,\size*-\maxY+\startY);
    \foreach \x in {0,1,...,\maxX}{
        \draw (\size*\x+\startX,\size*0+\startY) -- (\size*\x+\startX,\size*-\maxY+\startY);
    }
    %
    \foreach \y in {0,1,...,\maxY}{
        \draw (\size*0+\startX,\size*-\y+\startY) -- (\size*\maxX+\startX,\size*-\y+\startY);
    }
}
%
\pgfmathsetmacro{\centerX}{\meshsize/2}
\pgfmathsetmacro{\centerY}{-\meshsize/2}
%
\pgfmathsetmacro{\pollcenterX}{\centerX+\pollX}
\pgfmathsetmacro{\pollcenterY}{\centerY+\pollY}
\pgfmathsetmacro{\pollcornerX}{\pollcenterX - \pollsize/2}
\pgfmathsetmacro{\pollcornerY}{\pollcenterY + \pollsize/2}
%
\pgfmathsetmacro{\pAcoordX}{\pollcenterX + \pAX}
\pgfmathsetmacro{\pAcoordY}{\pollcenterY - \pAY}
\pgfmathsetmacro{\pBcoordX}{\pollcenterX + \pBX}
\pgfmathsetmacro{\pBcoordY}{\pollcenterY - \pBY}
\pgfmathsetmacro{\pCcoordX}{\pollcenterX + \pCX}
\pgfmathsetmacro{\pCcoordY}{\pollcenterY - \pCY}
%
\pgfmathsetmacro{\deltameshn}{\pollcenterX-\pollcornerX}
\pgfmathsetmacro{\deltameshd}{\centerX}
\pgfmathsetmacro{\deltamesh}{(\pollcenterX-\pollcornerX)/\centerX}
\pgfmathsetmacro{\deltaframen}{\deltameshn * \deltameshn}
\pgfmathsetmacro{\deltaframed}{\deltameshd * \deltameshd}
\pgfmathsetmacro{\deltaframe}{\deltamesh * \deltamesh}
%
\Grid{0}{0}{\meshsize}{\meshsize}{\unit}{1pt}
\Grid{\pollcornerX*\unit}{\pollcornerY*\unit}{\pollsize}{\pollsize}{\unit}{3pt}
%
\ifdim \deltamesh pt = 1 pt
    \dimline[label style={above=0.5ex,}, line style = {line width=0.7},extension start length=+0.3,extension end length=+0.3] {(\pollcornerX*\unit,1)}{(\pollcenterX*\unit,1)}{\huge $\delta^{k}_\mathrm{poll}=1$};
    
    \dimline[label style={above=0.5ex,}, line style = {line width=0.7},extension start length=+0.3,extension end length=+0.3] {(\meshsize*\unit - 1*\unit,1)}{(\meshsize*\unit,1)}{\huge $\delta^{k}_\mathrm{mesh}=1$};
\else
    \dimline[label style={above=0.5ex,}, line style = {line width=0.7},extension start length=+0.3,extension end length=+0.3] {(\pollcornerX*\unit,0.5)}{(\pollcenterX*\unit,0.5)}{\huge $\delta^{k}_\mathrm{poll}=\addfraction{\deltameshn}{\deltameshd}$};
    
    \dimline[label style={above=0.5ex,}, line style = {line width=0.7},extension start length=+0.3,extension end length=+0.3] {(\meshsize*\unit - 1*\unit,0.5)}{(\meshsize*\unit,0.5)}{\huge $\delta^{k}_\mathrm{mesh}=\addfraction{\deltaframen}{\deltaframed}$};
\fi
    \node[align=center, minimum width = 100mm, above right] at (0,3) {\huge Poll \textbf{\color{red}failure}};
    \node[star, star points=8, minimum size = 10mm, fill=blue, label={[fill=white,above right,label distance=0.2cm]\huge $x_k$}] (center) at (\pollcenterX*\unit,\pollcenterY*\unit) {};
    \node[circle, minimum size = 1mm, fill=red, label={[fill=white,below,label distance=-0.5cm]\huge $p_1$}] (p1) at (\pAcoordX*\unit,\pAcoordY*\unit) {};
    \node[circle, minimum size = 1mm, fill=red, label={[fill=white,right,label distance=0.5cm]\huge $p_2$}] (p2) at (\pBcoordX*\unit,\pBcoordY*\unit) {};
    \node[circle, minimum size = 1mm, fill=red, label={[fill=white]\huge $p_3$}] (p3) at (\pCcoordX*\unit,\pCcoordY*\unit) {};
\fi
%
\ifnum\frameopt < 4
    \def\unit{1.25}
    \def\meshsize{8}
    \def\pollX{0}
    \def\pollY{0}
    \def\pollsize{4}
    %
    \def\pAX{-\pollsize/2}
    \def\pAY{-1}
    \def\pBX{\pollsize/2}
    \def\pBY{\pollsize/2}
    \def\pCX{\pollsize/2}
    \def\pCY{-1}
\fi
%
\ifnum\frameopt=2
    % https://tex.stackexchange.com/a/183633
\newcommand{\Grid}[6]{%
    \def\startX{#1}
    \def\startY{#2}
    \def\maxX{#3}
    \def\maxY{#4}
    \def\size{#5}
    \def\borderthick{#6}
    %
    \draw[line width=\borderthick] (\size*0+\startX,\size*0+\startY) rectangle (\size*\maxX+\startX,\size*-\maxY+\startY);
    \foreach \x in {0,1,...,\maxX}{
        \draw (\size*\x+\startX,\size*0+\startY) -- (\size*\x+\startX,\size*-\maxY+\startY);
    }
    %
    \foreach \y in {0,1,...,\maxY}{
        \draw (\size*0+\startX,\size*-\y+\startY) -- (\size*\maxX+\startX,\size*-\y+\startY);
    }
}
%
\pgfmathsetmacro{\centerX}{\meshsize/2}
\pgfmathsetmacro{\centerY}{-\meshsize/2}
%
\pgfmathsetmacro{\pollcenterX}{\centerX+\pollX}
\pgfmathsetmacro{\pollcenterY}{\centerY+\pollY}
\pgfmathsetmacro{\pollcornerX}{\pollcenterX - \pollsize/2}
\pgfmathsetmacro{\pollcornerY}{\pollcenterY + \pollsize/2}
%
\pgfmathsetmacro{\pAcoordX}{\pollcenterX + \pAX}
\pgfmathsetmacro{\pAcoordY}{\pollcenterY - \pAY}
\pgfmathsetmacro{\pBcoordX}{\pollcenterX + \pBX}
\pgfmathsetmacro{\pBcoordY}{\pollcenterY - \pBY}
\pgfmathsetmacro{\pCcoordX}{\pollcenterX + \pCX}
\pgfmathsetmacro{\pCcoordY}{\pollcenterY - \pCY}
%
\pgfmathsetmacro{\deltameshn}{\pollcenterX-\pollcornerX}
\pgfmathsetmacro{\deltameshd}{\centerX}
\pgfmathsetmacro{\deltamesh}{(\pollcenterX-\pollcornerX)/\centerX}
\pgfmathsetmacro{\deltaframen}{\deltameshn * \deltameshn}
\pgfmathsetmacro{\deltaframed}{\deltameshd * \deltameshd}
\pgfmathsetmacro{\deltaframe}{\deltamesh * \deltamesh}
%
\Grid{0}{0}{\meshsize}{\meshsize}{\unit}{1pt}
\Grid{\pollcornerX*\unit}{\pollcornerY*\unit}{\pollsize}{\pollsize}{\unit}{3pt}
%
\ifdim \deltamesh pt = 1 pt
    \dimline[label style={above=0.5ex,}, line style = {line width=0.7},extension start length=+0.3,extension end length=+0.3] {(\pollcornerX*\unit,1)}{(\pollcenterX*\unit,1)}{\huge $\delta^{k}_\mathrm{poll}=1$};
    
    \dimline[label style={above=0.5ex,}, line style = {line width=0.7},extension start length=+0.3,extension end length=+0.3] {(\meshsize*\unit - 1*\unit,1)}{(\meshsize*\unit,1)}{\huge $\delta^{k}_\mathrm{mesh}=1$};
\else
    \dimline[label style={above=0.5ex,}, line style = {line width=0.7},extension start length=+0.3,extension end length=+0.3] {(\pollcornerX*\unit,0.5)}{(\pollcenterX*\unit,0.5)}{\huge $\delta^{k}_\mathrm{poll}=\addfraction{\deltameshn}{\deltameshd}$};
    
    \dimline[label style={above=0.5ex,}, line style = {line width=0.7},extension start length=+0.3,extension end length=+0.3] {(\meshsize*\unit - 1*\unit,0.5)}{(\meshsize*\unit,0.5)}{\huge $\delta^{k}_\mathrm{mesh}=\addfraction{\deltaframen}{\deltaframed}$};
\fi
    \node[circle, minimum size = 1mm, fill=red, label={[fill=white,above right,label distance=0.2cm]\huge $x_k$}] (center) at (\pollcenterX*\unit,\pollcenterY*\unit) {};
    \node[circle, minimum size = 1mm, fill=red, label={[fill=white,below,label distance=-0.5cm]\huge $p_4$}] (p1) at (\pAcoordX*\unit,\pAcoordY*\unit) {};
    \node[circle, minimum size = 1mm, fill=red, label={[fill=white,right,label distance=0.5cm]\huge $p_5$}] (p2) at (\pBcoordX*\unit,\pBcoordY*\unit) {};
    \node[circle, minimum size = 1mm, fill=red, label={[fill=white]\huge $p_6$}] (p3) at (\pCcoordX*\unit,\pCcoordY*\unit) {};
\fi
%
\ifnum\frameopt=3
    % https://tex.stackexchange.com/a/183633
\newcommand{\Grid}[6]{%
    \def\startX{#1}
    \def\startY{#2}
    \def\maxX{#3}
    \def\maxY{#4}
    \def\size{#5}
    \def\borderthick{#6}
    %
    \draw[line width=\borderthick] (\size*0+\startX,\size*0+\startY) rectangle (\size*\maxX+\startX,\size*-\maxY+\startY);
    \foreach \x in {0,1,...,\maxX}{
        \draw (\size*\x+\startX,\size*0+\startY) -- (\size*\x+\startX,\size*-\maxY+\startY);
    }
    %
    \foreach \y in {0,1,...,\maxY}{
        \draw (\size*0+\startX,\size*-\y+\startY) -- (\size*\maxX+\startX,\size*-\y+\startY);
    }
}
%
\pgfmathsetmacro{\centerX}{\meshsize/2}
\pgfmathsetmacro{\centerY}{-\meshsize/2}
%
\pgfmathsetmacro{\pollcenterX}{\centerX+\pollX}
\pgfmathsetmacro{\pollcenterY}{\centerY+\pollY}
\pgfmathsetmacro{\pollcornerX}{\pollcenterX - \pollsize/2}
\pgfmathsetmacro{\pollcornerY}{\pollcenterY + \pollsize/2}
%
\pgfmathsetmacro{\pAcoordX}{\pollcenterX + \pAX}
\pgfmathsetmacro{\pAcoordY}{\pollcenterY - \pAY}
\pgfmathsetmacro{\pBcoordX}{\pollcenterX + \pBX}
\pgfmathsetmacro{\pBcoordY}{\pollcenterY - \pBY}
\pgfmathsetmacro{\pCcoordX}{\pollcenterX + \pCX}
\pgfmathsetmacro{\pCcoordY}{\pollcenterY - \pCY}
%
\pgfmathsetmacro{\deltameshn}{\pollcenterX-\pollcornerX}
\pgfmathsetmacro{\deltameshd}{\centerX}
\pgfmathsetmacro{\deltamesh}{(\pollcenterX-\pollcornerX)/\centerX}
\pgfmathsetmacro{\deltaframen}{\deltameshn * \deltameshn}
\pgfmathsetmacro{\deltaframed}{\deltameshd * \deltameshd}
\pgfmathsetmacro{\deltaframe}{\deltamesh * \deltamesh}
%
\Grid{0}{0}{\meshsize}{\meshsize}{\unit}{1pt}
\Grid{\pollcornerX*\unit}{\pollcornerY*\unit}{\pollsize}{\pollsize}{\unit}{3pt}
%
\ifdim \deltamesh pt = 1 pt
    \dimline[label style={above=0.5ex,}, line style = {line width=0.7},extension start length=+0.3,extension end length=+0.3] {(\pollcornerX*\unit,1)}{(\pollcenterX*\unit,1)}{\huge $\delta^{k}_\mathrm{poll}=1$};
    
    \dimline[label style={above=0.5ex,}, line style = {line width=0.7},extension start length=+0.3,extension end length=+0.3] {(\meshsize*\unit - 1*\unit,1)}{(\meshsize*\unit,1)}{\huge $\delta^{k}_\mathrm{mesh}=1$};
\else
    \dimline[label style={above=0.5ex,}, line style = {line width=0.7},extension start length=+0.3,extension end length=+0.3] {(\pollcornerX*\unit,0.5)}{(\pollcenterX*\unit,0.5)}{\huge $\delta^{k}_\mathrm{poll}=\addfraction{\deltameshn}{\deltameshd}$};
    
    \dimline[label style={above=0.5ex,}, line style = {line width=0.7},extension start length=+0.3,extension end length=+0.3] {(\meshsize*\unit - 1*\unit,0.5)}{(\meshsize*\unit,0.5)}{\huge $\delta^{k}_\mathrm{mesh}=\addfraction{\deltaframen}{\deltaframed}$};
\fi
    \node[align=center, minimum width = 100mm, above right] at (0,2) {\huge Poll \textbf{\color{red}failure}};
    \node[star, star points=8, minimum size = 10mm, fill=blue, label={[fill=white,above right,label distance=0.2cm]\huge $x_k$}] (center) at (\pollcenterX*\unit,\pollcenterY*\unit) {};
    \node[circle, minimum size = 1mm, fill=red, label={[fill=white,below,label distance=-0.5cm]\huge $p_4$}] (p1) at (\pAcoordX*\unit,\pAcoordY*\unit) {};
    \node[circle, minimum size = 1mm, fill=red, label={[fill=white,right,label distance=0.5cm]\huge $p_5$}] (p2) at (\pBcoordX*\unit,\pBcoordY*\unit) {};
    \node[circle, minimum size = 1mm, fill=red, label={[fill=white]\huge $p_6$}] (p3) at (\pCcoordX*\unit,\pCcoordY*\unit) {};
\fi
%
\ifnum\frameopt < 7
    \def\unit{0.3125}
    \def\meshsize{32}
    \def\pollX{0}
    \def\pollY{0}
    \def\pollsize{8}
    %
    \def\pAX{-4}
    \def\pAY{\pollsize/2}
    \def\pBX{\pollsize/2}
    \def\pBY{0}
    \def\pCX{0}
    \def\pCY{-\pollsize/2}
\fi
\ifnum\frameopt = 4
    % https://tex.stackexchange.com/a/183633
\newcommand{\Grid}[6]{%
    \def\startX{#1}
    \def\startY{#2}
    \def\maxX{#3}
    \def\maxY{#4}
    \def\size{#5}
    \def\borderthick{#6}
    %
    \draw[line width=\borderthick] (\size*0+\startX,\size*0+\startY) rectangle (\size*\maxX+\startX,\size*-\maxY+\startY);
    \foreach \x in {0,1,...,\maxX}{
        \draw (\size*\x+\startX,\size*0+\startY) -- (\size*\x+\startX,\size*-\maxY+\startY);
    }
    %
    \foreach \y in {0,1,...,\maxY}{
        \draw (\size*0+\startX,\size*-\y+\startY) -- (\size*\maxX+\startX,\size*-\y+\startY);
    }
}
%
\pgfmathsetmacro{\centerX}{\meshsize/2}
\pgfmathsetmacro{\centerY}{-\meshsize/2}
%
\pgfmathsetmacro{\pollcenterX}{\centerX+\pollX}
\pgfmathsetmacro{\pollcenterY}{\centerY+\pollY}
\pgfmathsetmacro{\pollcornerX}{\pollcenterX - \pollsize/2}
\pgfmathsetmacro{\pollcornerY}{\pollcenterY + \pollsize/2}
%
\pgfmathsetmacro{\pAcoordX}{\pollcenterX + \pAX}
\pgfmathsetmacro{\pAcoordY}{\pollcenterY - \pAY}
\pgfmathsetmacro{\pBcoordX}{\pollcenterX + \pBX}
\pgfmathsetmacro{\pBcoordY}{\pollcenterY - \pBY}
\pgfmathsetmacro{\pCcoordX}{\pollcenterX + \pCX}
\pgfmathsetmacro{\pCcoordY}{\pollcenterY - \pCY}
%
\pgfmathsetmacro{\deltameshn}{\pollcenterX-\pollcornerX}
\pgfmathsetmacro{\deltameshd}{\centerX}
\pgfmathsetmacro{\deltamesh}{(\pollcenterX-\pollcornerX)/\centerX}
\pgfmathsetmacro{\deltaframen}{\deltameshn * \deltameshn}
\pgfmathsetmacro{\deltaframed}{\deltameshd * \deltameshd}
\pgfmathsetmacro{\deltaframe}{\deltamesh * \deltamesh}
%
\Grid{0}{0}{\meshsize}{\meshsize}{\unit}{1pt}
\Grid{\pollcornerX*\unit}{\pollcornerY*\unit}{\pollsize}{\pollsize}{\unit}{3pt}
%
\ifdim \deltamesh pt = 1 pt
    \dimline[label style={above=0.5ex,}, line style = {line width=0.7},extension start length=+0.3,extension end length=+0.3] {(\pollcornerX*\unit,1)}{(\pollcenterX*\unit,1)}{\huge $\delta^{k}_\mathrm{poll}=1$};
    
    \dimline[label style={above=0.5ex,}, line style = {line width=0.7},extension start length=+0.3,extension end length=+0.3] {(\meshsize*\unit - 1*\unit,1)}{(\meshsize*\unit,1)}{\huge $\delta^{k}_\mathrm{mesh}=1$};
\else
    \dimline[label style={above=0.5ex,}, line style = {line width=0.7},extension start length=+0.3,extension end length=+0.3] {(\pollcornerX*\unit,0.5)}{(\pollcenterX*\unit,0.5)}{\huge $\delta^{k}_\mathrm{poll}=\addfraction{\deltameshn}{\deltameshd}$};
    
    \dimline[label style={above=0.5ex,}, line style = {line width=0.7},extension start length=+0.3,extension end length=+0.3] {(\meshsize*\unit - 1*\unit,0.5)}{(\meshsize*\unit,0.5)}{\huge $\delta^{k}_\mathrm{mesh}=\addfraction{\deltaframen}{\deltaframed}$};
\fi
    \node[circle, minimum size = 1mm, fill=red, label={[fill=white,above right,label distance=0.2cm]\huge $x_k$}] (center) at (\pollcenterX*\unit,\pollcenterY*\unit) {};
    \node[circle, minimum size = 1mm, fill=red, label={[fill=white,below,label distance=-0.5cm]\huge $p_7$}] (p1) at (\pAcoordX*\unit,\pAcoordY*\unit) {};
    \node[circle, minimum size = 1mm, fill=red, label={[fill=white,right,label distance=0.5cm]\huge $p_8$}] (p2) at (\pBcoordX*\unit,\pBcoordY*\unit) {};
    \node[circle, minimum size = 1mm, fill=red, label={[fill=white]\huge $p_9$}] (p3) at (\pCcoordX*\unit,\pCcoordY*\unit) {};
\fi
%
\ifnum\frameopt = 5
    % https://tex.stackexchange.com/a/183633
\newcommand{\Grid}[6]{%
    \def\startX{#1}
    \def\startY{#2}
    \def\maxX{#3}
    \def\maxY{#4}
    \def\size{#5}
    \def\borderthick{#6}
    %
    \draw[line width=\borderthick] (\size*0+\startX,\size*0+\startY) rectangle (\size*\maxX+\startX,\size*-\maxY+\startY);
    \foreach \x in {0,1,...,\maxX}{
        \draw (\size*\x+\startX,\size*0+\startY) -- (\size*\x+\startX,\size*-\maxY+\startY);
    }
    %
    \foreach \y in {0,1,...,\maxY}{
        \draw (\size*0+\startX,\size*-\y+\startY) -- (\size*\maxX+\startX,\size*-\y+\startY);
    }
}
%
\pgfmathsetmacro{\centerX}{\meshsize/2}
\pgfmathsetmacro{\centerY}{-\meshsize/2}
%
\pgfmathsetmacro{\pollcenterX}{\centerX+\pollX}
\pgfmathsetmacro{\pollcenterY}{\centerY+\pollY}
\pgfmathsetmacro{\pollcornerX}{\pollcenterX - \pollsize/2}
\pgfmathsetmacro{\pollcornerY}{\pollcenterY + \pollsize/2}
%
\pgfmathsetmacro{\pAcoordX}{\pollcenterX + \pAX}
\pgfmathsetmacro{\pAcoordY}{\pollcenterY - \pAY}
\pgfmathsetmacro{\pBcoordX}{\pollcenterX + \pBX}
\pgfmathsetmacro{\pBcoordY}{\pollcenterY - \pBY}
\pgfmathsetmacro{\pCcoordX}{\pollcenterX + \pCX}
\pgfmathsetmacro{\pCcoordY}{\pollcenterY - \pCY}
%
\pgfmathsetmacro{\deltameshn}{\pollcenterX-\pollcornerX}
\pgfmathsetmacro{\deltameshd}{\centerX}
\pgfmathsetmacro{\deltamesh}{(\pollcenterX-\pollcornerX)/\centerX}
\pgfmathsetmacro{\deltaframen}{\deltameshn * \deltameshn}
\pgfmathsetmacro{\deltaframed}{\deltameshd * \deltameshd}
\pgfmathsetmacro{\deltaframe}{\deltamesh * \deltamesh}
%
\Grid{0}{0}{\meshsize}{\meshsize}{\unit}{1pt}
\Grid{\pollcornerX*\unit}{\pollcornerY*\unit}{\pollsize}{\pollsize}{\unit}{3pt}
%
\ifdim \deltamesh pt = 1 pt
    \dimline[label style={above=0.5ex,}, line style = {line width=0.7},extension start length=+0.3,extension end length=+0.3] {(\pollcornerX*\unit,1)}{(\pollcenterX*\unit,1)}{\huge $\delta^{k}_\mathrm{poll}=1$};
    
    \dimline[label style={above=0.5ex,}, line style = {line width=0.7},extension start length=+0.3,extension end length=+0.3] {(\meshsize*\unit - 1*\unit,1)}{(\meshsize*\unit,1)}{\huge $\delta^{k}_\mathrm{mesh}=1$};
\else
    \dimline[label style={above=0.5ex,}, line style = {line width=0.7},extension start length=+0.3,extension end length=+0.3] {(\pollcornerX*\unit,0.5)}{(\pollcenterX*\unit,0.5)}{\huge $\delta^{k}_\mathrm{poll}=\addfraction{\deltameshn}{\deltameshd}$};
    
    \dimline[label style={above=0.5ex,}, line style = {line width=0.7},extension start length=+0.3,extension end length=+0.3] {(\meshsize*\unit - 1*\unit,0.5)}{(\meshsize*\unit,0.5)}{\huge $\delta^{k}_\mathrm{mesh}=\addfraction{\deltaframen}{\deltaframed}$};
\fi
    \node[align=center, minimum width = 100mm, above right] at (0,2) {\huge Poll \textbf{\color{red}failure}};
    \node[star, star points=8, minimum size = 6mm, fill=blue, label={[fill=white,above right,label distance=0.2cm]\huge $x_k$}] (center) at (\pollcenterX*\unit,\pollcenterY*\unit) {};
    \node[circle, minimum size = 1mm, fill=red, label={[fill=white,below,label distance=-0.5cm]\huge $p_7$}] (p1) at (\pAcoordX*\unit,\pAcoordY*\unit) {};
    \node[circle, minimum size = 1mm, fill=red, label={[fill=white,right,label distance=0.5cm]\huge $p_8$}] (p2) at (\pBcoordX*\unit,\pBcoordY*\unit) {};
    \node[circle, minimum size = 1mm, fill=red, label={[fill=white]\huge $p_9$}] (p3) at (\pCcoordX*\unit,\pCcoordY*\unit) {};
\fi
%
\ifnum\frameopt = 6
    % https://tex.stackexchange.com/a/183633
\newcommand{\Grid}[6]{%
    \def\startX{#1}
    \def\startY{#2}
    \def\maxX{#3}
    \def\maxY{#4}
    \def\size{#5}
    \def\borderthick{#6}
    %
    \draw[line width=\borderthick] (\size*0+\startX,\size*0+\startY) rectangle (\size*\maxX+\startX,\size*-\maxY+\startY);
    \foreach \x in {0,1,...,\maxX}{
        \draw (\size*\x+\startX,\size*0+\startY) -- (\size*\x+\startX,\size*-\maxY+\startY);
    }
    %
    \foreach \y in {0,1,...,\maxY}{
        \draw (\size*0+\startX,\size*-\y+\startY) -- (\size*\maxX+\startX,\size*-\y+\startY);
    }
}
%
\pgfmathsetmacro{\centerX}{\meshsize/2}
\pgfmathsetmacro{\centerY}{-\meshsize/2}
%
\pgfmathsetmacro{\pollcenterX}{\centerX+\pollX}
\pgfmathsetmacro{\pollcenterY}{\centerY+\pollY}
\pgfmathsetmacro{\pollcornerX}{\pollcenterX - \pollsize/2}
\pgfmathsetmacro{\pollcornerY}{\pollcenterY + \pollsize/2}
%
\pgfmathsetmacro{\pAcoordX}{\pollcenterX + \pAX}
\pgfmathsetmacro{\pAcoordY}{\pollcenterY - \pAY}
\pgfmathsetmacro{\pBcoordX}{\pollcenterX + \pBX}
\pgfmathsetmacro{\pBcoordY}{\pollcenterY - \pBY}
\pgfmathsetmacro{\pCcoordX}{\pollcenterX + \pCX}
\pgfmathsetmacro{\pCcoordY}{\pollcenterY - \pCY}
%
\pgfmathsetmacro{\deltameshn}{\pollcenterX-\pollcornerX}
\pgfmathsetmacro{\deltameshd}{\centerX}
\pgfmathsetmacro{\deltamesh}{(\pollcenterX-\pollcornerX)/\centerX}
\pgfmathsetmacro{\deltaframen}{\deltameshn * \deltameshn}
\pgfmathsetmacro{\deltaframed}{\deltameshd * \deltameshd}
\pgfmathsetmacro{\deltaframe}{\deltamesh * \deltamesh}
%
\Grid{0}{0}{\meshsize}{\meshsize}{\unit}{1pt}
\Grid{\pollcornerX*\unit}{\pollcornerY*\unit}{\pollsize}{\pollsize}{\unit}{3pt}
%
\ifdim \deltamesh pt = 1 pt
    \dimline[label style={above=0.5ex,}, line style = {line width=0.7},extension start length=+0.3,extension end length=+0.3] {(\pollcornerX*\unit,1)}{(\pollcenterX*\unit,1)}{\huge $\delta^{k}_\mathrm{poll}=1$};
    
    \dimline[label style={above=0.5ex,}, line style = {line width=0.7},extension start length=+0.3,extension end length=+0.3] {(\meshsize*\unit - 1*\unit,1)}{(\meshsize*\unit,1)}{\huge $\delta^{k}_\mathrm{mesh}=1$};
\else
    \dimline[label style={above=0.5ex,}, line style = {line width=0.7},extension start length=+0.3,extension end length=+0.3] {(\pollcornerX*\unit,0.5)}{(\pollcenterX*\unit,0.5)}{\huge $\delta^{k}_\mathrm{poll}=\addfraction{\deltameshn}{\deltameshd}$};
    
    \dimline[label style={above=0.5ex,}, line style = {line width=0.7},extension start length=+0.3,extension end length=+0.3] {(\meshsize*\unit - 1*\unit,0.5)}{(\meshsize*\unit,0.5)}{\huge $\delta^{k}_\mathrm{mesh}=\addfraction{\deltaframen}{\deltaframed}$};
\fi
    \node[align=center, minimum width = 100mm, above right] at (0,2) {\huge Poll \textbf{\color{darkgreen}success}};
    \node[circle, minimum size = 1mm, fill=red, label={[fill=white,above right,label distance=0.2cm]\huge $x_k$}] (center) at (\pollcenterX*\unit,\pollcenterY*\unit) {};
    \node[star, star points=8, minimum size = 6mm, fill=blue, label={[fill=white,below,label distance=-0.5cm]\huge $p_7$}] (p1) at (\pAcoordX*\unit,\pAcoordY*\unit) {};
    \node[circle, minimum size = 1mm, fill=red, label={[fill=white,right,label distance=0.5cm]\huge $p_8$}] (p2) at (\pBcoordX*\unit,\pBcoordY*\unit) {};
    \node[circle, minimum size = 1mm, fill=red, label={[fill=white]\huge $p_9$}] (p3) at (\pCcoordX*\unit,\pCcoordY*\unit) {};
\fi
%
\ifnum\frameopt < 8
    \def\unit{0.625}
    \def\meshsize{16}
    \def\pollX{-2}
    \def\pollY{-2}
    \def\pollsize{8}
    %
    \def\pAX{3}
    \def\pAY{-\pollsize/2}
    \def\pBX{-\pollsize/2}
    \def\pBY{0}
    \def\pCX{2}
    \def\pCY{\pollsize/2}
\fi
%
\ifnum\frameopt = 7
    % https://tex.stackexchange.com/a/183633
\newcommand{\Grid}[6]{%
    \def\startX{#1}
    \def\startY{#2}
    \def\maxX{#3}
    \def\maxY{#4}
    \def\size{#5}
    \def\borderthick{#6}
    %
    \draw[line width=\borderthick] (\size*0+\startX,\size*0+\startY) rectangle (\size*\maxX+\startX,\size*-\maxY+\startY);
    \foreach \x in {0,1,...,\maxX}{
        \draw (\size*\x+\startX,\size*0+\startY) -- (\size*\x+\startX,\size*-\maxY+\startY);
    }
    %
    \foreach \y in {0,1,...,\maxY}{
        \draw (\size*0+\startX,\size*-\y+\startY) -- (\size*\maxX+\startX,\size*-\y+\startY);
    }
}
%
\pgfmathsetmacro{\centerX}{\meshsize/2}
\pgfmathsetmacro{\centerY}{-\meshsize/2}
%
\pgfmathsetmacro{\pollcenterX}{\centerX+\pollX}
\pgfmathsetmacro{\pollcenterY}{\centerY+\pollY}
\pgfmathsetmacro{\pollcornerX}{\pollcenterX - \pollsize/2}
\pgfmathsetmacro{\pollcornerY}{\pollcenterY + \pollsize/2}
%
\pgfmathsetmacro{\pAcoordX}{\pollcenterX + \pAX}
\pgfmathsetmacro{\pAcoordY}{\pollcenterY - \pAY}
\pgfmathsetmacro{\pBcoordX}{\pollcenterX + \pBX}
\pgfmathsetmacro{\pBcoordY}{\pollcenterY - \pBY}
\pgfmathsetmacro{\pCcoordX}{\pollcenterX + \pCX}
\pgfmathsetmacro{\pCcoordY}{\pollcenterY - \pCY}
%
\pgfmathsetmacro{\deltameshn}{\pollcenterX-\pollcornerX}
\pgfmathsetmacro{\deltameshd}{\centerX}
\pgfmathsetmacro{\deltamesh}{(\pollcenterX-\pollcornerX)/\centerX}
\pgfmathsetmacro{\deltaframen}{\deltameshn * \deltameshn}
\pgfmathsetmacro{\deltaframed}{\deltameshd * \deltameshd}
\pgfmathsetmacro{\deltaframe}{\deltamesh * \deltamesh}
%
\Grid{0}{0}{\meshsize}{\meshsize}{\unit}{1pt}
\Grid{\pollcornerX*\unit}{\pollcornerY*\unit}{\pollsize}{\pollsize}{\unit}{3pt}
%
\ifdim \deltamesh pt = 1 pt
    \dimline[label style={above=0.5ex,}, line style = {line width=0.7},extension start length=+0.3,extension end length=+0.3] {(\pollcornerX*\unit,1)}{(\pollcenterX*\unit,1)}{\huge $\delta^{k}_\mathrm{poll}=1$};
    
    \dimline[label style={above=0.5ex,}, line style = {line width=0.7},extension start length=+0.3,extension end length=+0.3] {(\meshsize*\unit - 1*\unit,1)}{(\meshsize*\unit,1)}{\huge $\delta^{k}_\mathrm{mesh}=1$};
\else
    \dimline[label style={above=0.5ex,}, line style = {line width=0.7},extension start length=+0.3,extension end length=+0.3] {(\pollcornerX*\unit,0.5)}{(\pollcenterX*\unit,0.5)}{\huge $\delta^{k}_\mathrm{poll}=\addfraction{\deltameshn}{\deltameshd}$};
    
    \dimline[label style={above=0.5ex,}, line style = {line width=0.7},extension start length=+0.3,extension end length=+0.3] {(\meshsize*\unit - 1*\unit,0.5)}{(\meshsize*\unit,0.5)}{\huge $\delta^{k}_\mathrm{mesh}=\addfraction{\deltaframen}{\deltaframed}$};
\fi
    \node[circle, minimum size = 1mm, fill=red, label={[fill=white,above right,label distance=0.2cm]\huge $x_k$}] (center) at (\pollcenterX*\unit,\pollcenterY*\unit) {};
    \node[circle, minimum size = 1mm, fill=red, label={[fill=white,below,label distance=-0.5cm]\huge $p_{10}$}] (p1) at (\pAcoordX*\unit,\pAcoordY*\unit) {};
    \node[circle, minimum size = 1mm, fill=red, label={[fill=white,right,label distance=0.5cm]\huge $p_{11}$}] (p2) at (\pBcoordX*\unit,\pBcoordY*\unit) {};
    \node[circle, minimum size = 1mm, fill=red, label={[fill=white]\huge $p_{12}$}] (p3) at (\pCcoordX*\unit,\pCcoordY*\unit) {};
\fi
%
\ifnum\frameopt < 9
    \def\unit{0.3125}
    \def\meshsize{32}
    \def\pollX{-4}
    \def\pollY{-4}
    \def\pollsize{8}
    %
    \def\pAX{-2}
    \def\pAY{-\pollsize/2}
    \def\pBX{\pollsize/2}
    \def\pBY{-\pollsize/2}
    \def\pCX{\pollsize/2}
    \def\pCY{\pollsize/2}
\fi
%
\ifnum\frameopt = 8
    % https://tex.stackexchange.com/a/183633
\newcommand{\Grid}[6]{%
    \def\startX{#1}
    \def\startY{#2}
    \def\maxX{#3}
    \def\maxY{#4}
    \def\size{#5}
    \def\borderthick{#6}
    %
    \draw[line width=\borderthick] (\size*0+\startX,\size*0+\startY) rectangle (\size*\maxX+\startX,\size*-\maxY+\startY);
    \foreach \x in {0,1,...,\maxX}{
        \draw (\size*\x+\startX,\size*0+\startY) -- (\size*\x+\startX,\size*-\maxY+\startY);
    }
    %
    \foreach \y in {0,1,...,\maxY}{
        \draw (\size*0+\startX,\size*-\y+\startY) -- (\size*\maxX+\startX,\size*-\y+\startY);
    }
}
%
\pgfmathsetmacro{\centerX}{\meshsize/2}
\pgfmathsetmacro{\centerY}{-\meshsize/2}
%
\pgfmathsetmacro{\pollcenterX}{\centerX+\pollX}
\pgfmathsetmacro{\pollcenterY}{\centerY+\pollY}
\pgfmathsetmacro{\pollcornerX}{\pollcenterX - \pollsize/2}
\pgfmathsetmacro{\pollcornerY}{\pollcenterY + \pollsize/2}
%
\pgfmathsetmacro{\pAcoordX}{\pollcenterX + \pAX}
\pgfmathsetmacro{\pAcoordY}{\pollcenterY - \pAY}
\pgfmathsetmacro{\pBcoordX}{\pollcenterX + \pBX}
\pgfmathsetmacro{\pBcoordY}{\pollcenterY - \pBY}
\pgfmathsetmacro{\pCcoordX}{\pollcenterX + \pCX}
\pgfmathsetmacro{\pCcoordY}{\pollcenterY - \pCY}
%
\pgfmathsetmacro{\deltameshn}{\pollcenterX-\pollcornerX}
\pgfmathsetmacro{\deltameshd}{\centerX}
\pgfmathsetmacro{\deltamesh}{(\pollcenterX-\pollcornerX)/\centerX}
\pgfmathsetmacro{\deltaframen}{\deltameshn * \deltameshn}
\pgfmathsetmacro{\deltaframed}{\deltameshd * \deltameshd}
\pgfmathsetmacro{\deltaframe}{\deltamesh * \deltamesh}
%
\Grid{0}{0}{\meshsize}{\meshsize}{\unit}{1pt}
\Grid{\pollcornerX*\unit}{\pollcornerY*\unit}{\pollsize}{\pollsize}{\unit}{3pt}
%
\ifdim \deltamesh pt = 1 pt
    \dimline[label style={above=0.5ex,}, line style = {line width=0.7},extension start length=+0.3,extension end length=+0.3] {(\pollcornerX*\unit,1)}{(\pollcenterX*\unit,1)}{\huge $\delta^{k}_\mathrm{poll}=1$};
    
    \dimline[label style={above=0.5ex,}, line style = {line width=0.7},extension start length=+0.3,extension end length=+0.3] {(\meshsize*\unit - 1*\unit,1)}{(\meshsize*\unit,1)}{\huge $\delta^{k}_\mathrm{mesh}=1$};
\else
    \dimline[label style={above=0.5ex,}, line style = {line width=0.7},extension start length=+0.3,extension end length=+0.3] {(\pollcornerX*\unit,0.5)}{(\pollcenterX*\unit,0.5)}{\huge $\delta^{k}_\mathrm{poll}=\addfraction{\deltameshn}{\deltameshd}$};
    
    \dimline[label style={above=0.5ex,}, line style = {line width=0.7},extension start length=+0.3,extension end length=+0.3] {(\meshsize*\unit - 1*\unit,0.5)}{(\meshsize*\unit,0.5)}{\huge $\delta^{k}_\mathrm{mesh}=\addfraction{\deltaframen}{\deltaframed}$};
\fi
    \node[align=center, minimum width = 100mm, above right] at (0,2) {\huge Poll \textbf{\color{red}failure}};
    \node[circle, minimum size = 1mm, fill=red, label={[fill=white,above right,label distance=0.2cm]\huge $x_k$}] (center) at (\pollcenterX*\unit,\pollcenterY*\unit) {};
    \node[circle, minimum size = 1mm, fill=red, label={[fill=white,below,label distance=-0.5cm]\huge $p_{13}$}] (p1) at (\pAcoordX*\unit,\pAcoordY*\unit) {};
    \node[circle, minimum size = 1mm, fill=red, label={[fill=white,right,label distance=0.5cm]\huge $p_{14}$}] (p2) at (\pBcoordX*\unit,\pBcoordY*\unit) {};
    \node[circle, minimum size = 1mm, fill=red, label={[fill=white]\huge $p_{15}$}] (p3) at (\pCcoordX*\unit,\pCcoordY*\unit) {};
\fi
%
\ifnum\frameopt < 10
    \def\unit{0.15625}
    \def\meshsize{64}
    \def\pollX{-8}
    \def\pollY{-8}
    \def\pollsize{8}
    %
    \def\pAX{1}
    \def\pAY{\pollsize/2}
    \def\pBX{-\pollsize/2}
    \def\pBY{0}
    \def\pCX{\pollsize/2}
    \def\pCY{-\pollsize/2}
\fi
%
\ifnum\frameopt = 9
    % https://tex.stackexchange.com/a/183633
\newcommand{\Grid}[6]{%
    \def\startX{#1}
    \def\startY{#2}
    \def\maxX{#3}
    \def\maxY{#4}
    \def\size{#5}
    \def\borderthick{#6}
    %
    \draw[line width=\borderthick] (\size*0+\startX,\size*0+\startY) rectangle (\size*\maxX+\startX,\size*-\maxY+\startY);
    \foreach \x in {0,1,...,\maxX}{
        \draw (\size*\x+\startX,\size*0+\startY) -- (\size*\x+\startX,\size*-\maxY+\startY);
    }
    %
    \foreach \y in {0,1,...,\maxY}{
        \draw (\size*0+\startX,\size*-\y+\startY) -- (\size*\maxX+\startX,\size*-\y+\startY);
    }
}
%
\pgfmathsetmacro{\centerX}{\meshsize/2}
\pgfmathsetmacro{\centerY}{-\meshsize/2}
%
\pgfmathsetmacro{\pollcenterX}{\centerX+\pollX}
\pgfmathsetmacro{\pollcenterY}{\centerY+\pollY}
\pgfmathsetmacro{\pollcornerX}{\pollcenterX - \pollsize/2}
\pgfmathsetmacro{\pollcornerY}{\pollcenterY + \pollsize/2}
%
\pgfmathsetmacro{\pAcoordX}{\pollcenterX + \pAX}
\pgfmathsetmacro{\pAcoordY}{\pollcenterY - \pAY}
\pgfmathsetmacro{\pBcoordX}{\pollcenterX + \pBX}
\pgfmathsetmacro{\pBcoordY}{\pollcenterY - \pBY}
\pgfmathsetmacro{\pCcoordX}{\pollcenterX + \pCX}
\pgfmathsetmacro{\pCcoordY}{\pollcenterY - \pCY}
%
\pgfmathsetmacro{\deltameshn}{\pollcenterX-\pollcornerX}
\pgfmathsetmacro{\deltameshd}{\centerX}
\pgfmathsetmacro{\deltamesh}{(\pollcenterX-\pollcornerX)/\centerX}
\pgfmathsetmacro{\deltaframen}{\deltameshn * \deltameshn}
\pgfmathsetmacro{\deltaframed}{\deltameshd * \deltameshd}
\pgfmathsetmacro{\deltaframe}{\deltamesh * \deltamesh}
%
\Grid{0}{0}{\meshsize}{\meshsize}{\unit}{1pt}
\Grid{\pollcornerX*\unit}{\pollcornerY*\unit}{\pollsize}{\pollsize}{\unit}{3pt}
%
\ifdim \deltamesh pt = 1 pt
    \dimline[label style={above=0.5ex,}, line style = {line width=0.7},extension start length=+0.3,extension end length=+0.3] {(\pollcornerX*\unit,1)}{(\pollcenterX*\unit,1)}{\huge $\delta^{k}_\mathrm{poll}=1$};
    
    \dimline[label style={above=0.5ex,}, line style = {line width=0.7},extension start length=+0.3,extension end length=+0.3] {(\meshsize*\unit - 1*\unit,1)}{(\meshsize*\unit,1)}{\huge $\delta^{k}_\mathrm{mesh}=1$};
\else
    \dimline[label style={above=0.5ex,}, line style = {line width=0.7},extension start length=+0.3,extension end length=+0.3] {(\pollcornerX*\unit,0.5)}{(\pollcenterX*\unit,0.5)}{\huge $\delta^{k}_\mathrm{poll}=\addfraction{\deltameshn}{\deltameshd}$};
    
    \dimline[label style={above=0.5ex,}, line style = {line width=0.7},extension start length=+0.3,extension end length=+0.3] {(\meshsize*\unit - 1*\unit,0.5)}{(\meshsize*\unit,0.5)}{\huge $\delta^{k}_\mathrm{mesh}=\addfraction{\deltaframen}{\deltaframed}$};
\fi
    \node[align=center, minimum width = 100mm, above right] at (0,2) {\huge Poll \textbf{\color{red}failure}};
    \node[circle, minimum size = 1mm, fill=red] (center) at (\pollcenterX*\unit,\pollcenterY*\unit) {};
    \node[circle, minimum size = 1mm, fill=red] (p1) at (\pAcoordX*\unit,\pAcoordY*\unit) {};
    \node[circle, minimum size = 1mm, fill=red] (p2) at (\pBcoordX*\unit,\pBcoordY*\unit) {};
    \node[circle, minimum size = 1mm, fill=red] (p3) at (\pCcoordX*\unit,\pCcoordY*\unit) {};
\fi
%
\ifnum\frameopt > -1
    \draw[-, color=red, line width=2pt, opacity=1.0] (center) -- (p1);
    \draw[-, color=red, line width=2pt, opacity=1.0] (center) -- (p2);
    \draw[-, color=red, line width=2pt, opacity=1.0] (center) -- (p3);
\fi}%
				\end{tikzpicture}%
			}%
		\end{column}
	
	\end{columns}
	\vspace{-3em}
\end{frame}
\addtocounter{footnote}{-1}
%------------------------------------------------
\begin{frame}[t,label=abm_6]
	\frametitle{Overview of stochastic mesh adaptive direct search}
	\tikzstyle{background grid}=[draw, black!50,step=.5cm]
	%
	\uncover<1->{No gradient information available, blackbox is expensive and \emphasis{noisy}\ifshowcitations\footpartcite{Audet2020}\fi}%
	\only<9->{{\color{white}\ifshowcitations\footpartcite{Dzahini2022}\fi}}\\
	%
	\begin{columns}[t] % The "c" option specifies centered vertical alignment while the "t" option is used for top vertical alignment
		\begin{column}{.42\textwidth} % Left column and width
			\vspace{-1.2em}
			% Optimization problem
			\begin{exampleblock}{Objective and constraints}
				\vspace{-1.2em}
				\begin{equation*}
					\begin{aligned}
						& \underset{\mathbf{x}}{\text{min}}
						& & f(\mathbf{x}) = \mathbb{E}_{\Theta}\left[{f}_{\Theta}(\mathbf{x})\right]\\
						& \text{subject to}
						& & {c}(\mathbf{x}) = \mathbb{E}_{\Theta}\left[{c}_{\Theta}(\mathbf{x})\right] \le 0\\
						& \text{where}
						& & \mathbf{x}:\text{ variables}~\Theta\mathrm{:realizations}
					\end{aligned}
				\end{equation*}
			\end{exampleblock}
			\only<-8>{
				\small
				\begin{itemize}\itemsep0em
					\item Constructs estimates of objective:
					%
					\vspace{-1em}
					\begin{equation*} \label{eq:samplerate}
						f^{k} = \dfrac{1}{n^k}\sum_{i=1}^{n^k} f_{\Theta_{0,i}}({x}_k)
					\end{equation*}
					\vspace{-1em}
					%
					\item<2-> $n^k$ is the sampling rate
					\item<3-> Tracks uncertainty interval $\mathcal{I}(\delta^k_\text{poll})$ in the estimate
				\end{itemize}
			}%
			\only<9->{~\\~\\\centering Constraint handling using the \textit{progressive barrier} approach \footnotemark[2]}%
		\end{column}
		%
		\begin{column}{.5\textwidth} % Left column and width
			\tikzstyle{background grid}=[draw, black!50,step=.5cm]
			
			\only<-2>{\def\frameopt{2}}%
			\only<3>{\def\frameopts{1}}%
			\only<4>{\def\frameopts{2}}%
			\only<5>{\def\frameopts{3}}%
			\only<6>{\def\frameopts{4}}%
			\only<7>{\def\frameopts{5}}%
			\only<8->{\def\frameopts{6}}%

			\tikzstyle{background grid}=[draw, black!50,step=.1cm]
			\hspace*{5em}\raisebox{-1em}{%
				\begin{tikzpicture}[scale=0.45, every node/.style={scale=0.45}, remember picture, overlay] %show background grid, 
					\only<-2>{\ifnum\frameopt < 2
    \def\unit{5}
    \def\meshsize{2}
    \def\pollX{0}
    \def\pollY{0}
    \def\pollsize{2}
    %
    \def\pAX{-1}
    \def\pAY{\pollsize/2}
    \def\pBX{\pollsize/2}
    \def\pBY{0}
    \def\pCX{0}
    \def\pCY{-\pollsize/2}
\fi
%
\ifnum\frameopt=-2
    % https://tex.stackexchange.com/a/183633
\newcommand{\Grid}[6]{%
    \def\startX{#1}
    \def\startY{#2}
    \def\maxX{#3}
    \def\maxY{#4}
    \def\size{#5}
    \def\borderthick{#6}
    %
    \draw[line width=\borderthick] (\size*0+\startX,\size*0+\startY) rectangle (\size*\maxX+\startX,\size*-\maxY+\startY);
    \foreach \x in {0,1,...,\maxX}{
        \draw (\size*\x+\startX,\size*0+\startY) -- (\size*\x+\startX,\size*-\maxY+\startY);
    }
    %
    \foreach \y in {0,1,...,\maxY}{
        \draw (\size*0+\startX,\size*-\y+\startY) -- (\size*\maxX+\startX,\size*-\y+\startY);
    }
}
%
\pgfmathsetmacro{\centerX}{\meshsize/2}
\pgfmathsetmacro{\centerY}{-\meshsize/2}
%
\pgfmathsetmacro{\pollcenterX}{\centerX+\pollX}
\pgfmathsetmacro{\pollcenterY}{\centerY+\pollY}
\pgfmathsetmacro{\pollcornerX}{\pollcenterX - \pollsize/2}
\pgfmathsetmacro{\pollcornerY}{\pollcenterY + \pollsize/2}
%
\pgfmathsetmacro{\pAcoordX}{\pollcenterX + \pAX}
\pgfmathsetmacro{\pAcoordY}{\pollcenterY - \pAY}
\pgfmathsetmacro{\pBcoordX}{\pollcenterX + \pBX}
\pgfmathsetmacro{\pBcoordY}{\pollcenterY - \pBY}
\pgfmathsetmacro{\pCcoordX}{\pollcenterX + \pCX}
\pgfmathsetmacro{\pCcoordY}{\pollcenterY - \pCY}
%
\pgfmathsetmacro{\deltameshn}{\pollcenterX-\pollcornerX}
\pgfmathsetmacro{\deltameshd}{\centerX}
\pgfmathsetmacro{\deltamesh}{(\pollcenterX-\pollcornerX)/\centerX}
\pgfmathsetmacro{\deltaframen}{\deltameshn * \deltameshn}
\pgfmathsetmacro{\deltaframed}{\deltameshd * \deltameshd}
\pgfmathsetmacro{\deltaframe}{\deltamesh * \deltamesh}
%
\Grid{0}{0}{\meshsize}{\meshsize}{\unit}{1pt}
\Grid{\pollcornerX*\unit}{\pollcornerY*\unit}{\pollsize}{\pollsize}{\unit}{3pt}
%
\ifdim \deltamesh pt = 1 pt
    \dimline[label style={above=0.5ex,}, line style = {line width=0.7},extension start length=+0.3,extension end length=+0.3] {(\pollcornerX*\unit,1)}{(\pollcenterX*\unit,1)}{\huge $\delta^{k}_\mathrm{poll}=1$};
    
    \dimline[label style={above=0.5ex,}, line style = {line width=0.7},extension start length=+0.3,extension end length=+0.3] {(\meshsize*\unit - 1*\unit,1)}{(\meshsize*\unit,1)}{\huge $\delta^{k}_\mathrm{mesh}=1$};
\else
    \dimline[label style={above=0.5ex,}, line style = {line width=0.7},extension start length=+0.3,extension end length=+0.3] {(\pollcornerX*\unit,0.5)}{(\pollcenterX*\unit,0.5)}{\huge $\delta^{k}_\mathrm{poll}=\addfraction{\deltameshn}{\deltameshd}$};
    
    \dimline[label style={above=0.5ex,}, line style = {line width=0.7},extension start length=+0.3,extension end length=+0.3] {(\meshsize*\unit - 1*\unit,0.5)}{(\meshsize*\unit,0.5)}{\huge $\delta^{k}_\mathrm{mesh}=\addfraction{\deltaframen}{\deltaframed}$};
\fi
\fi
%
\ifnum\frameopt=-1
    % https://tex.stackexchange.com/a/183633
\newcommand{\Grid}[6]{%
    \def\startX{#1}
    \def\startY{#2}
    \def\maxX{#3}
    \def\maxY{#4}
    \def\size{#5}
    \def\borderthick{#6}
    %
    \draw[line width=\borderthick] (\size*0+\startX,\size*0+\startY) rectangle (\size*\maxX+\startX,\size*-\maxY+\startY);
    \foreach \x in {0,1,...,\maxX}{
        \draw (\size*\x+\startX,\size*0+\startY) -- (\size*\x+\startX,\size*-\maxY+\startY);
    }
    %
    \foreach \y in {0,1,...,\maxY}{
        \draw (\size*0+\startX,\size*-\y+\startY) -- (\size*\maxX+\startX,\size*-\y+\startY);
    }
}
%
\pgfmathsetmacro{\centerX}{\meshsize/2}
\pgfmathsetmacro{\centerY}{-\meshsize/2}
%
\pgfmathsetmacro{\pollcenterX}{\centerX+\pollX}
\pgfmathsetmacro{\pollcenterY}{\centerY+\pollY}
\pgfmathsetmacro{\pollcornerX}{\pollcenterX - \pollsize/2}
\pgfmathsetmacro{\pollcornerY}{\pollcenterY + \pollsize/2}
%
\pgfmathsetmacro{\pAcoordX}{\pollcenterX + \pAX}
\pgfmathsetmacro{\pAcoordY}{\pollcenterY - \pAY}
\pgfmathsetmacro{\pBcoordX}{\pollcenterX + \pBX}
\pgfmathsetmacro{\pBcoordY}{\pollcenterY - \pBY}
\pgfmathsetmacro{\pCcoordX}{\pollcenterX + \pCX}
\pgfmathsetmacro{\pCcoordY}{\pollcenterY - \pCY}
%
\pgfmathsetmacro{\deltameshn}{\pollcenterX-\pollcornerX}
\pgfmathsetmacro{\deltameshd}{\centerX}
\pgfmathsetmacro{\deltamesh}{(\pollcenterX-\pollcornerX)/\centerX}
\pgfmathsetmacro{\deltaframen}{\deltameshn * \deltameshn}
\pgfmathsetmacro{\deltaframed}{\deltameshd * \deltameshd}
\pgfmathsetmacro{\deltaframe}{\deltamesh * \deltamesh}
%
\Grid{0}{0}{\meshsize}{\meshsize}{\unit}{1pt}
\Grid{\pollcornerX*\unit}{\pollcornerY*\unit}{\pollsize}{\pollsize}{\unit}{3pt}
%
\ifdim \deltamesh pt = 1 pt
    \dimline[label style={above=0.5ex,}, line style = {line width=0.7},extension start length=+0.3,extension end length=+0.3] {(\pollcornerX*\unit,1)}{(\pollcenterX*\unit,1)}{\huge $\delta^{k}_\mathrm{poll}=1$};
    
    \dimline[label style={above=0.5ex,}, line style = {line width=0.7},extension start length=+0.3,extension end length=+0.3] {(\meshsize*\unit - 1*\unit,1)}{(\meshsize*\unit,1)}{\huge $\delta^{k}_\mathrm{mesh}=1$};
\else
    \dimline[label style={above=0.5ex,}, line style = {line width=0.7},extension start length=+0.3,extension end length=+0.3] {(\pollcornerX*\unit,0.5)}{(\pollcenterX*\unit,0.5)}{\huge $\delta^{k}_\mathrm{poll}=\addfraction{\deltameshn}{\deltameshd}$};
    
    \dimline[label style={above=0.5ex,}, line style = {line width=0.7},extension start length=+0.3,extension end length=+0.3] {(\meshsize*\unit - 1*\unit,0.5)}{(\meshsize*\unit,0.5)}{\huge $\delta^{k}_\mathrm{mesh}=\addfraction{\deltaframen}{\deltaframed}$};
\fi
    \node[circle, minimum size = 1mm, fill=red, label={[fill=white,above right,label distance=0.2cm]\huge $x_k$}] (center) at (\pollcenterX*\unit,\pollcenterY*\unit) {};
\fi
%
\ifnum\frameopt=0
    % https://tex.stackexchange.com/a/183633
\newcommand{\Grid}[6]{%
    \def\startX{#1}
    \def\startY{#2}
    \def\maxX{#3}
    \def\maxY{#4}
    \def\size{#5}
    \def\borderthick{#6}
    %
    \draw[line width=\borderthick] (\size*0+\startX,\size*0+\startY) rectangle (\size*\maxX+\startX,\size*-\maxY+\startY);
    \foreach \x in {0,1,...,\maxX}{
        \draw (\size*\x+\startX,\size*0+\startY) -- (\size*\x+\startX,\size*-\maxY+\startY);
    }
    %
    \foreach \y in {0,1,...,\maxY}{
        \draw (\size*0+\startX,\size*-\y+\startY) -- (\size*\maxX+\startX,\size*-\y+\startY);
    }
}
%
\pgfmathsetmacro{\centerX}{\meshsize/2}
\pgfmathsetmacro{\centerY}{-\meshsize/2}
%
\pgfmathsetmacro{\pollcenterX}{\centerX+\pollX}
\pgfmathsetmacro{\pollcenterY}{\centerY+\pollY}
\pgfmathsetmacro{\pollcornerX}{\pollcenterX - \pollsize/2}
\pgfmathsetmacro{\pollcornerY}{\pollcenterY + \pollsize/2}
%
\pgfmathsetmacro{\pAcoordX}{\pollcenterX + \pAX}
\pgfmathsetmacro{\pAcoordY}{\pollcenterY - \pAY}
\pgfmathsetmacro{\pBcoordX}{\pollcenterX + \pBX}
\pgfmathsetmacro{\pBcoordY}{\pollcenterY - \pBY}
\pgfmathsetmacro{\pCcoordX}{\pollcenterX + \pCX}
\pgfmathsetmacro{\pCcoordY}{\pollcenterY - \pCY}
%
\pgfmathsetmacro{\deltameshn}{\pollcenterX-\pollcornerX}
\pgfmathsetmacro{\deltameshd}{\centerX}
\pgfmathsetmacro{\deltamesh}{(\pollcenterX-\pollcornerX)/\centerX}
\pgfmathsetmacro{\deltaframen}{\deltameshn * \deltameshn}
\pgfmathsetmacro{\deltaframed}{\deltameshd * \deltameshd}
\pgfmathsetmacro{\deltaframe}{\deltamesh * \deltamesh}
%
\Grid{0}{0}{\meshsize}{\meshsize}{\unit}{1pt}
\Grid{\pollcornerX*\unit}{\pollcornerY*\unit}{\pollsize}{\pollsize}{\unit}{3pt}
%
\ifdim \deltamesh pt = 1 pt
    \dimline[label style={above=0.5ex,}, line style = {line width=0.7},extension start length=+0.3,extension end length=+0.3] {(\pollcornerX*\unit,1)}{(\pollcenterX*\unit,1)}{\huge $\delta^{k}_\mathrm{poll}=1$};
    
    \dimline[label style={above=0.5ex,}, line style = {line width=0.7},extension start length=+0.3,extension end length=+0.3] {(\meshsize*\unit - 1*\unit,1)}{(\meshsize*\unit,1)}{\huge $\delta^{k}_\mathrm{mesh}=1$};
\else
    \dimline[label style={above=0.5ex,}, line style = {line width=0.7},extension start length=+0.3,extension end length=+0.3] {(\pollcornerX*\unit,0.5)}{(\pollcenterX*\unit,0.5)}{\huge $\delta^{k}_\mathrm{poll}=\addfraction{\deltameshn}{\deltameshd}$};
    
    \dimline[label style={above=0.5ex,}, line style = {line width=0.7},extension start length=+0.3,extension end length=+0.3] {(\meshsize*\unit - 1*\unit,0.5)}{(\meshsize*\unit,0.5)}{\huge $\delta^{k}_\mathrm{mesh}=\addfraction{\deltaframen}{\deltaframed}$};
\fi
    \node[circle, minimum size = 1mm, fill=red, label={[fill=white,above right,label distance=0.2cm]\huge $x_k$}] (center) at (\pollcenterX*\unit,\pollcenterY*\unit) {};
    \node[circle, minimum size = 1mm, fill=red, label={[fill=white,below,label distance=-0.5cm]\huge $p_1$}] (p1) at (\pAcoordX*\unit,\pAcoordY*\unit) {};
    \node[circle, minimum size = 1mm, fill=red, label={[fill=white,right,label distance=0.5cm]\huge $p_2$}] (p2) at (\pBcoordX*\unit,\pBcoordY*\unit) {};
    \node[circle, minimum size = 1mm, fill=red, label={[fill=white]\huge $p_3$}] (p3) at (\pCcoordX*\unit,\pCcoordY*\unit) {};
\fi
%
\ifnum\frameopt=1
    % https://tex.stackexchange.com/a/183633
\newcommand{\Grid}[6]{%
    \def\startX{#1}
    \def\startY{#2}
    \def\maxX{#3}
    \def\maxY{#4}
    \def\size{#5}
    \def\borderthick{#6}
    %
    \draw[line width=\borderthick] (\size*0+\startX,\size*0+\startY) rectangle (\size*\maxX+\startX,\size*-\maxY+\startY);
    \foreach \x in {0,1,...,\maxX}{
        \draw (\size*\x+\startX,\size*0+\startY) -- (\size*\x+\startX,\size*-\maxY+\startY);
    }
    %
    \foreach \y in {0,1,...,\maxY}{
        \draw (\size*0+\startX,\size*-\y+\startY) -- (\size*\maxX+\startX,\size*-\y+\startY);
    }
}
%
\pgfmathsetmacro{\centerX}{\meshsize/2}
\pgfmathsetmacro{\centerY}{-\meshsize/2}
%
\pgfmathsetmacro{\pollcenterX}{\centerX+\pollX}
\pgfmathsetmacro{\pollcenterY}{\centerY+\pollY}
\pgfmathsetmacro{\pollcornerX}{\pollcenterX - \pollsize/2}
\pgfmathsetmacro{\pollcornerY}{\pollcenterY + \pollsize/2}
%
\pgfmathsetmacro{\pAcoordX}{\pollcenterX + \pAX}
\pgfmathsetmacro{\pAcoordY}{\pollcenterY - \pAY}
\pgfmathsetmacro{\pBcoordX}{\pollcenterX + \pBX}
\pgfmathsetmacro{\pBcoordY}{\pollcenterY - \pBY}
\pgfmathsetmacro{\pCcoordX}{\pollcenterX + \pCX}
\pgfmathsetmacro{\pCcoordY}{\pollcenterY - \pCY}
%
\pgfmathsetmacro{\deltameshn}{\pollcenterX-\pollcornerX}
\pgfmathsetmacro{\deltameshd}{\centerX}
\pgfmathsetmacro{\deltamesh}{(\pollcenterX-\pollcornerX)/\centerX}
\pgfmathsetmacro{\deltaframen}{\deltameshn * \deltameshn}
\pgfmathsetmacro{\deltaframed}{\deltameshd * \deltameshd}
\pgfmathsetmacro{\deltaframe}{\deltamesh * \deltamesh}
%
\Grid{0}{0}{\meshsize}{\meshsize}{\unit}{1pt}
\Grid{\pollcornerX*\unit}{\pollcornerY*\unit}{\pollsize}{\pollsize}{\unit}{3pt}
%
\ifdim \deltamesh pt = 1 pt
    \dimline[label style={above=0.5ex,}, line style = {line width=0.7},extension start length=+0.3,extension end length=+0.3] {(\pollcornerX*\unit,1)}{(\pollcenterX*\unit,1)}{\huge $\delta^{k}_\mathrm{poll}=1$};
    
    \dimline[label style={above=0.5ex,}, line style = {line width=0.7},extension start length=+0.3,extension end length=+0.3] {(\meshsize*\unit - 1*\unit,1)}{(\meshsize*\unit,1)}{\huge $\delta^{k}_\mathrm{mesh}=1$};
\else
    \dimline[label style={above=0.5ex,}, line style = {line width=0.7},extension start length=+0.3,extension end length=+0.3] {(\pollcornerX*\unit,0.5)}{(\pollcenterX*\unit,0.5)}{\huge $\delta^{k}_\mathrm{poll}=\addfraction{\deltameshn}{\deltameshd}$};
    
    \dimline[label style={above=0.5ex,}, line style = {line width=0.7},extension start length=+0.3,extension end length=+0.3] {(\meshsize*\unit - 1*\unit,0.5)}{(\meshsize*\unit,0.5)}{\huge $\delta^{k}_\mathrm{mesh}=\addfraction{\deltaframen}{\deltaframed}$};
\fi
    \node[align=center, minimum width = 100mm, above right] at (0,3) {\huge Poll \textbf{\color{red}failure}};
    \node[star, star points=8, minimum size = 10mm, fill=blue, label={[fill=white,above right,label distance=0.2cm]\huge $x_k$}] (center) at (\pollcenterX*\unit,\pollcenterY*\unit) {};
    \node[circle, minimum size = 1mm, fill=red, label={[fill=white,below,label distance=-0.5cm]\huge $p_1$}] (p1) at (\pAcoordX*\unit,\pAcoordY*\unit) {};
    \node[circle, minimum size = 1mm, fill=red, label={[fill=white,right,label distance=0.5cm]\huge $p_2$}] (p2) at (\pBcoordX*\unit,\pBcoordY*\unit) {};
    \node[circle, minimum size = 1mm, fill=red, label={[fill=white]\huge $p_3$}] (p3) at (\pCcoordX*\unit,\pCcoordY*\unit) {};
\fi
%
\ifnum\frameopt < 4
    \def\unit{1.25}
    \def\meshsize{8}
    \def\pollX{0}
    \def\pollY{0}
    \def\pollsize{4}
    %
    \def\pAX{-\pollsize/2}
    \def\pAY{-1}
    \def\pBX{\pollsize/2}
    \def\pBY{\pollsize/2}
    \def\pCX{\pollsize/2}
    \def\pCY{-1}
\fi
%
\ifnum\frameopt=2
    % https://tex.stackexchange.com/a/183633
\newcommand{\Grid}[6]{%
    \def\startX{#1}
    \def\startY{#2}
    \def\maxX{#3}
    \def\maxY{#4}
    \def\size{#5}
    \def\borderthick{#6}
    %
    \draw[line width=\borderthick] (\size*0+\startX,\size*0+\startY) rectangle (\size*\maxX+\startX,\size*-\maxY+\startY);
    \foreach \x in {0,1,...,\maxX}{
        \draw (\size*\x+\startX,\size*0+\startY) -- (\size*\x+\startX,\size*-\maxY+\startY);
    }
    %
    \foreach \y in {0,1,...,\maxY}{
        \draw (\size*0+\startX,\size*-\y+\startY) -- (\size*\maxX+\startX,\size*-\y+\startY);
    }
}
%
\pgfmathsetmacro{\centerX}{\meshsize/2}
\pgfmathsetmacro{\centerY}{-\meshsize/2}
%
\pgfmathsetmacro{\pollcenterX}{\centerX+\pollX}
\pgfmathsetmacro{\pollcenterY}{\centerY+\pollY}
\pgfmathsetmacro{\pollcornerX}{\pollcenterX - \pollsize/2}
\pgfmathsetmacro{\pollcornerY}{\pollcenterY + \pollsize/2}
%
\pgfmathsetmacro{\pAcoordX}{\pollcenterX + \pAX}
\pgfmathsetmacro{\pAcoordY}{\pollcenterY - \pAY}
\pgfmathsetmacro{\pBcoordX}{\pollcenterX + \pBX}
\pgfmathsetmacro{\pBcoordY}{\pollcenterY - \pBY}
\pgfmathsetmacro{\pCcoordX}{\pollcenterX + \pCX}
\pgfmathsetmacro{\pCcoordY}{\pollcenterY - \pCY}
%
\pgfmathsetmacro{\deltameshn}{\pollcenterX-\pollcornerX}
\pgfmathsetmacro{\deltameshd}{\centerX}
\pgfmathsetmacro{\deltamesh}{(\pollcenterX-\pollcornerX)/\centerX}
\pgfmathsetmacro{\deltaframen}{\deltameshn * \deltameshn}
\pgfmathsetmacro{\deltaframed}{\deltameshd * \deltameshd}
\pgfmathsetmacro{\deltaframe}{\deltamesh * \deltamesh}
%
\Grid{0}{0}{\meshsize}{\meshsize}{\unit}{1pt}
\Grid{\pollcornerX*\unit}{\pollcornerY*\unit}{\pollsize}{\pollsize}{\unit}{3pt}
%
\ifdim \deltamesh pt = 1 pt
    \dimline[label style={above=0.5ex,}, line style = {line width=0.7},extension start length=+0.3,extension end length=+0.3] {(\pollcornerX*\unit,1)}{(\pollcenterX*\unit,1)}{\huge $\delta^{k}_\mathrm{poll}=1$};
    
    \dimline[label style={above=0.5ex,}, line style = {line width=0.7},extension start length=+0.3,extension end length=+0.3] {(\meshsize*\unit - 1*\unit,1)}{(\meshsize*\unit,1)}{\huge $\delta^{k}_\mathrm{mesh}=1$};
\else
    \dimline[label style={above=0.5ex,}, line style = {line width=0.7},extension start length=+0.3,extension end length=+0.3] {(\pollcornerX*\unit,0.5)}{(\pollcenterX*\unit,0.5)}{\huge $\delta^{k}_\mathrm{poll}=\addfraction{\deltameshn}{\deltameshd}$};
    
    \dimline[label style={above=0.5ex,}, line style = {line width=0.7},extension start length=+0.3,extension end length=+0.3] {(\meshsize*\unit - 1*\unit,0.5)}{(\meshsize*\unit,0.5)}{\huge $\delta^{k}_\mathrm{mesh}=\addfraction{\deltaframen}{\deltaframed}$};
\fi
    \node[circle, minimum size = 1mm, fill=red, label={[fill=white,above right,label distance=0.2cm]\huge $x_k$}] (center) at (\pollcenterX*\unit,\pollcenterY*\unit) {};
    \node[circle, minimum size = 1mm, fill=red, label={[fill=white,below,label distance=-0.5cm]\huge $p_4$}] (p1) at (\pAcoordX*\unit,\pAcoordY*\unit) {};
    \node[circle, minimum size = 1mm, fill=red, label={[fill=white,right,label distance=0.5cm]\huge $p_5$}] (p2) at (\pBcoordX*\unit,\pBcoordY*\unit) {};
    \node[circle, minimum size = 1mm, fill=red, label={[fill=white]\huge $p_6$}] (p3) at (\pCcoordX*\unit,\pCcoordY*\unit) {};
\fi
%
\ifnum\frameopt=3
    % https://tex.stackexchange.com/a/183633
\newcommand{\Grid}[6]{%
    \def\startX{#1}
    \def\startY{#2}
    \def\maxX{#3}
    \def\maxY{#4}
    \def\size{#5}
    \def\borderthick{#6}
    %
    \draw[line width=\borderthick] (\size*0+\startX,\size*0+\startY) rectangle (\size*\maxX+\startX,\size*-\maxY+\startY);
    \foreach \x in {0,1,...,\maxX}{
        \draw (\size*\x+\startX,\size*0+\startY) -- (\size*\x+\startX,\size*-\maxY+\startY);
    }
    %
    \foreach \y in {0,1,...,\maxY}{
        \draw (\size*0+\startX,\size*-\y+\startY) -- (\size*\maxX+\startX,\size*-\y+\startY);
    }
}
%
\pgfmathsetmacro{\centerX}{\meshsize/2}
\pgfmathsetmacro{\centerY}{-\meshsize/2}
%
\pgfmathsetmacro{\pollcenterX}{\centerX+\pollX}
\pgfmathsetmacro{\pollcenterY}{\centerY+\pollY}
\pgfmathsetmacro{\pollcornerX}{\pollcenterX - \pollsize/2}
\pgfmathsetmacro{\pollcornerY}{\pollcenterY + \pollsize/2}
%
\pgfmathsetmacro{\pAcoordX}{\pollcenterX + \pAX}
\pgfmathsetmacro{\pAcoordY}{\pollcenterY - \pAY}
\pgfmathsetmacro{\pBcoordX}{\pollcenterX + \pBX}
\pgfmathsetmacro{\pBcoordY}{\pollcenterY - \pBY}
\pgfmathsetmacro{\pCcoordX}{\pollcenterX + \pCX}
\pgfmathsetmacro{\pCcoordY}{\pollcenterY - \pCY}
%
\pgfmathsetmacro{\deltameshn}{\pollcenterX-\pollcornerX}
\pgfmathsetmacro{\deltameshd}{\centerX}
\pgfmathsetmacro{\deltamesh}{(\pollcenterX-\pollcornerX)/\centerX}
\pgfmathsetmacro{\deltaframen}{\deltameshn * \deltameshn}
\pgfmathsetmacro{\deltaframed}{\deltameshd * \deltameshd}
\pgfmathsetmacro{\deltaframe}{\deltamesh * \deltamesh}
%
\Grid{0}{0}{\meshsize}{\meshsize}{\unit}{1pt}
\Grid{\pollcornerX*\unit}{\pollcornerY*\unit}{\pollsize}{\pollsize}{\unit}{3pt}
%
\ifdim \deltamesh pt = 1 pt
    \dimline[label style={above=0.5ex,}, line style = {line width=0.7},extension start length=+0.3,extension end length=+0.3] {(\pollcornerX*\unit,1)}{(\pollcenterX*\unit,1)}{\huge $\delta^{k}_\mathrm{poll}=1$};
    
    \dimline[label style={above=0.5ex,}, line style = {line width=0.7},extension start length=+0.3,extension end length=+0.3] {(\meshsize*\unit - 1*\unit,1)}{(\meshsize*\unit,1)}{\huge $\delta^{k}_\mathrm{mesh}=1$};
\else
    \dimline[label style={above=0.5ex,}, line style = {line width=0.7},extension start length=+0.3,extension end length=+0.3] {(\pollcornerX*\unit,0.5)}{(\pollcenterX*\unit,0.5)}{\huge $\delta^{k}_\mathrm{poll}=\addfraction{\deltameshn}{\deltameshd}$};
    
    \dimline[label style={above=0.5ex,}, line style = {line width=0.7},extension start length=+0.3,extension end length=+0.3] {(\meshsize*\unit - 1*\unit,0.5)}{(\meshsize*\unit,0.5)}{\huge $\delta^{k}_\mathrm{mesh}=\addfraction{\deltaframen}{\deltaframed}$};
\fi
    \node[align=center, minimum width = 100mm, above right] at (0,2) {\huge Poll \textbf{\color{red}failure}};
    \node[star, star points=8, minimum size = 10mm, fill=blue, label={[fill=white,above right,label distance=0.2cm]\huge $x_k$}] (center) at (\pollcenterX*\unit,\pollcenterY*\unit) {};
    \node[circle, minimum size = 1mm, fill=red, label={[fill=white,below,label distance=-0.5cm]\huge $p_4$}] (p1) at (\pAcoordX*\unit,\pAcoordY*\unit) {};
    \node[circle, minimum size = 1mm, fill=red, label={[fill=white,right,label distance=0.5cm]\huge $p_5$}] (p2) at (\pBcoordX*\unit,\pBcoordY*\unit) {};
    \node[circle, minimum size = 1mm, fill=red, label={[fill=white]\huge $p_6$}] (p3) at (\pCcoordX*\unit,\pCcoordY*\unit) {};
\fi
%
\ifnum\frameopt < 7
    \def\unit{0.3125}
    \def\meshsize{32}
    \def\pollX{0}
    \def\pollY{0}
    \def\pollsize{8}
    %
    \def\pAX{-4}
    \def\pAY{\pollsize/2}
    \def\pBX{\pollsize/2}
    \def\pBY{0}
    \def\pCX{0}
    \def\pCY{-\pollsize/2}
\fi
\ifnum\frameopt = 4
    % https://tex.stackexchange.com/a/183633
\newcommand{\Grid}[6]{%
    \def\startX{#1}
    \def\startY{#2}
    \def\maxX{#3}
    \def\maxY{#4}
    \def\size{#5}
    \def\borderthick{#6}
    %
    \draw[line width=\borderthick] (\size*0+\startX,\size*0+\startY) rectangle (\size*\maxX+\startX,\size*-\maxY+\startY);
    \foreach \x in {0,1,...,\maxX}{
        \draw (\size*\x+\startX,\size*0+\startY) -- (\size*\x+\startX,\size*-\maxY+\startY);
    }
    %
    \foreach \y in {0,1,...,\maxY}{
        \draw (\size*0+\startX,\size*-\y+\startY) -- (\size*\maxX+\startX,\size*-\y+\startY);
    }
}
%
\pgfmathsetmacro{\centerX}{\meshsize/2}
\pgfmathsetmacro{\centerY}{-\meshsize/2}
%
\pgfmathsetmacro{\pollcenterX}{\centerX+\pollX}
\pgfmathsetmacro{\pollcenterY}{\centerY+\pollY}
\pgfmathsetmacro{\pollcornerX}{\pollcenterX - \pollsize/2}
\pgfmathsetmacro{\pollcornerY}{\pollcenterY + \pollsize/2}
%
\pgfmathsetmacro{\pAcoordX}{\pollcenterX + \pAX}
\pgfmathsetmacro{\pAcoordY}{\pollcenterY - \pAY}
\pgfmathsetmacro{\pBcoordX}{\pollcenterX + \pBX}
\pgfmathsetmacro{\pBcoordY}{\pollcenterY - \pBY}
\pgfmathsetmacro{\pCcoordX}{\pollcenterX + \pCX}
\pgfmathsetmacro{\pCcoordY}{\pollcenterY - \pCY}
%
\pgfmathsetmacro{\deltameshn}{\pollcenterX-\pollcornerX}
\pgfmathsetmacro{\deltameshd}{\centerX}
\pgfmathsetmacro{\deltamesh}{(\pollcenterX-\pollcornerX)/\centerX}
\pgfmathsetmacro{\deltaframen}{\deltameshn * \deltameshn}
\pgfmathsetmacro{\deltaframed}{\deltameshd * \deltameshd}
\pgfmathsetmacro{\deltaframe}{\deltamesh * \deltamesh}
%
\Grid{0}{0}{\meshsize}{\meshsize}{\unit}{1pt}
\Grid{\pollcornerX*\unit}{\pollcornerY*\unit}{\pollsize}{\pollsize}{\unit}{3pt}
%
\ifdim \deltamesh pt = 1 pt
    \dimline[label style={above=0.5ex,}, line style = {line width=0.7},extension start length=+0.3,extension end length=+0.3] {(\pollcornerX*\unit,1)}{(\pollcenterX*\unit,1)}{\huge $\delta^{k}_\mathrm{poll}=1$};
    
    \dimline[label style={above=0.5ex,}, line style = {line width=0.7},extension start length=+0.3,extension end length=+0.3] {(\meshsize*\unit - 1*\unit,1)}{(\meshsize*\unit,1)}{\huge $\delta^{k}_\mathrm{mesh}=1$};
\else
    \dimline[label style={above=0.5ex,}, line style = {line width=0.7},extension start length=+0.3,extension end length=+0.3] {(\pollcornerX*\unit,0.5)}{(\pollcenterX*\unit,0.5)}{\huge $\delta^{k}_\mathrm{poll}=\addfraction{\deltameshn}{\deltameshd}$};
    
    \dimline[label style={above=0.5ex,}, line style = {line width=0.7},extension start length=+0.3,extension end length=+0.3] {(\meshsize*\unit - 1*\unit,0.5)}{(\meshsize*\unit,0.5)}{\huge $\delta^{k}_\mathrm{mesh}=\addfraction{\deltaframen}{\deltaframed}$};
\fi
    \node[circle, minimum size = 1mm, fill=red, label={[fill=white,above right,label distance=0.2cm]\huge $x_k$}] (center) at (\pollcenterX*\unit,\pollcenterY*\unit) {};
    \node[circle, minimum size = 1mm, fill=red, label={[fill=white,below,label distance=-0.5cm]\huge $p_7$}] (p1) at (\pAcoordX*\unit,\pAcoordY*\unit) {};
    \node[circle, minimum size = 1mm, fill=red, label={[fill=white,right,label distance=0.5cm]\huge $p_8$}] (p2) at (\pBcoordX*\unit,\pBcoordY*\unit) {};
    \node[circle, minimum size = 1mm, fill=red, label={[fill=white]\huge $p_9$}] (p3) at (\pCcoordX*\unit,\pCcoordY*\unit) {};
\fi
%
\ifnum\frameopt = 5
    % https://tex.stackexchange.com/a/183633
\newcommand{\Grid}[6]{%
    \def\startX{#1}
    \def\startY{#2}
    \def\maxX{#3}
    \def\maxY{#4}
    \def\size{#5}
    \def\borderthick{#6}
    %
    \draw[line width=\borderthick] (\size*0+\startX,\size*0+\startY) rectangle (\size*\maxX+\startX,\size*-\maxY+\startY);
    \foreach \x in {0,1,...,\maxX}{
        \draw (\size*\x+\startX,\size*0+\startY) -- (\size*\x+\startX,\size*-\maxY+\startY);
    }
    %
    \foreach \y in {0,1,...,\maxY}{
        \draw (\size*0+\startX,\size*-\y+\startY) -- (\size*\maxX+\startX,\size*-\y+\startY);
    }
}
%
\pgfmathsetmacro{\centerX}{\meshsize/2}
\pgfmathsetmacro{\centerY}{-\meshsize/2}
%
\pgfmathsetmacro{\pollcenterX}{\centerX+\pollX}
\pgfmathsetmacro{\pollcenterY}{\centerY+\pollY}
\pgfmathsetmacro{\pollcornerX}{\pollcenterX - \pollsize/2}
\pgfmathsetmacro{\pollcornerY}{\pollcenterY + \pollsize/2}
%
\pgfmathsetmacro{\pAcoordX}{\pollcenterX + \pAX}
\pgfmathsetmacro{\pAcoordY}{\pollcenterY - \pAY}
\pgfmathsetmacro{\pBcoordX}{\pollcenterX + \pBX}
\pgfmathsetmacro{\pBcoordY}{\pollcenterY - \pBY}
\pgfmathsetmacro{\pCcoordX}{\pollcenterX + \pCX}
\pgfmathsetmacro{\pCcoordY}{\pollcenterY - \pCY}
%
\pgfmathsetmacro{\deltameshn}{\pollcenterX-\pollcornerX}
\pgfmathsetmacro{\deltameshd}{\centerX}
\pgfmathsetmacro{\deltamesh}{(\pollcenterX-\pollcornerX)/\centerX}
\pgfmathsetmacro{\deltaframen}{\deltameshn * \deltameshn}
\pgfmathsetmacro{\deltaframed}{\deltameshd * \deltameshd}
\pgfmathsetmacro{\deltaframe}{\deltamesh * \deltamesh}
%
\Grid{0}{0}{\meshsize}{\meshsize}{\unit}{1pt}
\Grid{\pollcornerX*\unit}{\pollcornerY*\unit}{\pollsize}{\pollsize}{\unit}{3pt}
%
\ifdim \deltamesh pt = 1 pt
    \dimline[label style={above=0.5ex,}, line style = {line width=0.7},extension start length=+0.3,extension end length=+0.3] {(\pollcornerX*\unit,1)}{(\pollcenterX*\unit,1)}{\huge $\delta^{k}_\mathrm{poll}=1$};
    
    \dimline[label style={above=0.5ex,}, line style = {line width=0.7},extension start length=+0.3,extension end length=+0.3] {(\meshsize*\unit - 1*\unit,1)}{(\meshsize*\unit,1)}{\huge $\delta^{k}_\mathrm{mesh}=1$};
\else
    \dimline[label style={above=0.5ex,}, line style = {line width=0.7},extension start length=+0.3,extension end length=+0.3] {(\pollcornerX*\unit,0.5)}{(\pollcenterX*\unit,0.5)}{\huge $\delta^{k}_\mathrm{poll}=\addfraction{\deltameshn}{\deltameshd}$};
    
    \dimline[label style={above=0.5ex,}, line style = {line width=0.7},extension start length=+0.3,extension end length=+0.3] {(\meshsize*\unit - 1*\unit,0.5)}{(\meshsize*\unit,0.5)}{\huge $\delta^{k}_\mathrm{mesh}=\addfraction{\deltaframen}{\deltaframed}$};
\fi
    \node[align=center, minimum width = 100mm, above right] at (0,2) {\huge Poll \textbf{\color{red}failure}};
    \node[star, star points=8, minimum size = 6mm, fill=blue, label={[fill=white,above right,label distance=0.2cm]\huge $x_k$}] (center) at (\pollcenterX*\unit,\pollcenterY*\unit) {};
    \node[circle, minimum size = 1mm, fill=red, label={[fill=white,below,label distance=-0.5cm]\huge $p_7$}] (p1) at (\pAcoordX*\unit,\pAcoordY*\unit) {};
    \node[circle, minimum size = 1mm, fill=red, label={[fill=white,right,label distance=0.5cm]\huge $p_8$}] (p2) at (\pBcoordX*\unit,\pBcoordY*\unit) {};
    \node[circle, minimum size = 1mm, fill=red, label={[fill=white]\huge $p_9$}] (p3) at (\pCcoordX*\unit,\pCcoordY*\unit) {};
\fi
%
\ifnum\frameopt = 6
    % https://tex.stackexchange.com/a/183633
\newcommand{\Grid}[6]{%
    \def\startX{#1}
    \def\startY{#2}
    \def\maxX{#3}
    \def\maxY{#4}
    \def\size{#5}
    \def\borderthick{#6}
    %
    \draw[line width=\borderthick] (\size*0+\startX,\size*0+\startY) rectangle (\size*\maxX+\startX,\size*-\maxY+\startY);
    \foreach \x in {0,1,...,\maxX}{
        \draw (\size*\x+\startX,\size*0+\startY) -- (\size*\x+\startX,\size*-\maxY+\startY);
    }
    %
    \foreach \y in {0,1,...,\maxY}{
        \draw (\size*0+\startX,\size*-\y+\startY) -- (\size*\maxX+\startX,\size*-\y+\startY);
    }
}
%
\pgfmathsetmacro{\centerX}{\meshsize/2}
\pgfmathsetmacro{\centerY}{-\meshsize/2}
%
\pgfmathsetmacro{\pollcenterX}{\centerX+\pollX}
\pgfmathsetmacro{\pollcenterY}{\centerY+\pollY}
\pgfmathsetmacro{\pollcornerX}{\pollcenterX - \pollsize/2}
\pgfmathsetmacro{\pollcornerY}{\pollcenterY + \pollsize/2}
%
\pgfmathsetmacro{\pAcoordX}{\pollcenterX + \pAX}
\pgfmathsetmacro{\pAcoordY}{\pollcenterY - \pAY}
\pgfmathsetmacro{\pBcoordX}{\pollcenterX + \pBX}
\pgfmathsetmacro{\pBcoordY}{\pollcenterY - \pBY}
\pgfmathsetmacro{\pCcoordX}{\pollcenterX + \pCX}
\pgfmathsetmacro{\pCcoordY}{\pollcenterY - \pCY}
%
\pgfmathsetmacro{\deltameshn}{\pollcenterX-\pollcornerX}
\pgfmathsetmacro{\deltameshd}{\centerX}
\pgfmathsetmacro{\deltamesh}{(\pollcenterX-\pollcornerX)/\centerX}
\pgfmathsetmacro{\deltaframen}{\deltameshn * \deltameshn}
\pgfmathsetmacro{\deltaframed}{\deltameshd * \deltameshd}
\pgfmathsetmacro{\deltaframe}{\deltamesh * \deltamesh}
%
\Grid{0}{0}{\meshsize}{\meshsize}{\unit}{1pt}
\Grid{\pollcornerX*\unit}{\pollcornerY*\unit}{\pollsize}{\pollsize}{\unit}{3pt}
%
\ifdim \deltamesh pt = 1 pt
    \dimline[label style={above=0.5ex,}, line style = {line width=0.7},extension start length=+0.3,extension end length=+0.3] {(\pollcornerX*\unit,1)}{(\pollcenterX*\unit,1)}{\huge $\delta^{k}_\mathrm{poll}=1$};
    
    \dimline[label style={above=0.5ex,}, line style = {line width=0.7},extension start length=+0.3,extension end length=+0.3] {(\meshsize*\unit - 1*\unit,1)}{(\meshsize*\unit,1)}{\huge $\delta^{k}_\mathrm{mesh}=1$};
\else
    \dimline[label style={above=0.5ex,}, line style = {line width=0.7},extension start length=+0.3,extension end length=+0.3] {(\pollcornerX*\unit,0.5)}{(\pollcenterX*\unit,0.5)}{\huge $\delta^{k}_\mathrm{poll}=\addfraction{\deltameshn}{\deltameshd}$};
    
    \dimline[label style={above=0.5ex,}, line style = {line width=0.7},extension start length=+0.3,extension end length=+0.3] {(\meshsize*\unit - 1*\unit,0.5)}{(\meshsize*\unit,0.5)}{\huge $\delta^{k}_\mathrm{mesh}=\addfraction{\deltaframen}{\deltaframed}$};
\fi
    \node[align=center, minimum width = 100mm, above right] at (0,2) {\huge Poll \textbf{\color{darkgreen}success}};
    \node[circle, minimum size = 1mm, fill=red, label={[fill=white,above right,label distance=0.2cm]\huge $x_k$}] (center) at (\pollcenterX*\unit,\pollcenterY*\unit) {};
    \node[star, star points=8, minimum size = 6mm, fill=blue, label={[fill=white,below,label distance=-0.5cm]\huge $p_7$}] (p1) at (\pAcoordX*\unit,\pAcoordY*\unit) {};
    \node[circle, minimum size = 1mm, fill=red, label={[fill=white,right,label distance=0.5cm]\huge $p_8$}] (p2) at (\pBcoordX*\unit,\pBcoordY*\unit) {};
    \node[circle, minimum size = 1mm, fill=red, label={[fill=white]\huge $p_9$}] (p3) at (\pCcoordX*\unit,\pCcoordY*\unit) {};
\fi
%
\ifnum\frameopt < 8
    \def\unit{0.625}
    \def\meshsize{16}
    \def\pollX{-2}
    \def\pollY{-2}
    \def\pollsize{8}
    %
    \def\pAX{3}
    \def\pAY{-\pollsize/2}
    \def\pBX{-\pollsize/2}
    \def\pBY{0}
    \def\pCX{2}
    \def\pCY{\pollsize/2}
\fi
%
\ifnum\frameopt = 7
    % https://tex.stackexchange.com/a/183633
\newcommand{\Grid}[6]{%
    \def\startX{#1}
    \def\startY{#2}
    \def\maxX{#3}
    \def\maxY{#4}
    \def\size{#5}
    \def\borderthick{#6}
    %
    \draw[line width=\borderthick] (\size*0+\startX,\size*0+\startY) rectangle (\size*\maxX+\startX,\size*-\maxY+\startY);
    \foreach \x in {0,1,...,\maxX}{
        \draw (\size*\x+\startX,\size*0+\startY) -- (\size*\x+\startX,\size*-\maxY+\startY);
    }
    %
    \foreach \y in {0,1,...,\maxY}{
        \draw (\size*0+\startX,\size*-\y+\startY) -- (\size*\maxX+\startX,\size*-\y+\startY);
    }
}
%
\pgfmathsetmacro{\centerX}{\meshsize/2}
\pgfmathsetmacro{\centerY}{-\meshsize/2}
%
\pgfmathsetmacro{\pollcenterX}{\centerX+\pollX}
\pgfmathsetmacro{\pollcenterY}{\centerY+\pollY}
\pgfmathsetmacro{\pollcornerX}{\pollcenterX - \pollsize/2}
\pgfmathsetmacro{\pollcornerY}{\pollcenterY + \pollsize/2}
%
\pgfmathsetmacro{\pAcoordX}{\pollcenterX + \pAX}
\pgfmathsetmacro{\pAcoordY}{\pollcenterY - \pAY}
\pgfmathsetmacro{\pBcoordX}{\pollcenterX + \pBX}
\pgfmathsetmacro{\pBcoordY}{\pollcenterY - \pBY}
\pgfmathsetmacro{\pCcoordX}{\pollcenterX + \pCX}
\pgfmathsetmacro{\pCcoordY}{\pollcenterY - \pCY}
%
\pgfmathsetmacro{\deltameshn}{\pollcenterX-\pollcornerX}
\pgfmathsetmacro{\deltameshd}{\centerX}
\pgfmathsetmacro{\deltamesh}{(\pollcenterX-\pollcornerX)/\centerX}
\pgfmathsetmacro{\deltaframen}{\deltameshn * \deltameshn}
\pgfmathsetmacro{\deltaframed}{\deltameshd * \deltameshd}
\pgfmathsetmacro{\deltaframe}{\deltamesh * \deltamesh}
%
\Grid{0}{0}{\meshsize}{\meshsize}{\unit}{1pt}
\Grid{\pollcornerX*\unit}{\pollcornerY*\unit}{\pollsize}{\pollsize}{\unit}{3pt}
%
\ifdim \deltamesh pt = 1 pt
    \dimline[label style={above=0.5ex,}, line style = {line width=0.7},extension start length=+0.3,extension end length=+0.3] {(\pollcornerX*\unit,1)}{(\pollcenterX*\unit,1)}{\huge $\delta^{k}_\mathrm{poll}=1$};
    
    \dimline[label style={above=0.5ex,}, line style = {line width=0.7},extension start length=+0.3,extension end length=+0.3] {(\meshsize*\unit - 1*\unit,1)}{(\meshsize*\unit,1)}{\huge $\delta^{k}_\mathrm{mesh}=1$};
\else
    \dimline[label style={above=0.5ex,}, line style = {line width=0.7},extension start length=+0.3,extension end length=+0.3] {(\pollcornerX*\unit,0.5)}{(\pollcenterX*\unit,0.5)}{\huge $\delta^{k}_\mathrm{poll}=\addfraction{\deltameshn}{\deltameshd}$};
    
    \dimline[label style={above=0.5ex,}, line style = {line width=0.7},extension start length=+0.3,extension end length=+0.3] {(\meshsize*\unit - 1*\unit,0.5)}{(\meshsize*\unit,0.5)}{\huge $\delta^{k}_\mathrm{mesh}=\addfraction{\deltaframen}{\deltaframed}$};
\fi
    \node[circle, minimum size = 1mm, fill=red, label={[fill=white,above right,label distance=0.2cm]\huge $x_k$}] (center) at (\pollcenterX*\unit,\pollcenterY*\unit) {};
    \node[circle, minimum size = 1mm, fill=red, label={[fill=white,below,label distance=-0.5cm]\huge $p_{10}$}] (p1) at (\pAcoordX*\unit,\pAcoordY*\unit) {};
    \node[circle, minimum size = 1mm, fill=red, label={[fill=white,right,label distance=0.5cm]\huge $p_{11}$}] (p2) at (\pBcoordX*\unit,\pBcoordY*\unit) {};
    \node[circle, minimum size = 1mm, fill=red, label={[fill=white]\huge $p_{12}$}] (p3) at (\pCcoordX*\unit,\pCcoordY*\unit) {};
\fi
%
\ifnum\frameopt < 9
    \def\unit{0.3125}
    \def\meshsize{32}
    \def\pollX{-4}
    \def\pollY{-4}
    \def\pollsize{8}
    %
    \def\pAX{-2}
    \def\pAY{-\pollsize/2}
    \def\pBX{\pollsize/2}
    \def\pBY{-\pollsize/2}
    \def\pCX{\pollsize/2}
    \def\pCY{\pollsize/2}
\fi
%
\ifnum\frameopt = 8
    % https://tex.stackexchange.com/a/183633
\newcommand{\Grid}[6]{%
    \def\startX{#1}
    \def\startY{#2}
    \def\maxX{#3}
    \def\maxY{#4}
    \def\size{#5}
    \def\borderthick{#6}
    %
    \draw[line width=\borderthick] (\size*0+\startX,\size*0+\startY) rectangle (\size*\maxX+\startX,\size*-\maxY+\startY);
    \foreach \x in {0,1,...,\maxX}{
        \draw (\size*\x+\startX,\size*0+\startY) -- (\size*\x+\startX,\size*-\maxY+\startY);
    }
    %
    \foreach \y in {0,1,...,\maxY}{
        \draw (\size*0+\startX,\size*-\y+\startY) -- (\size*\maxX+\startX,\size*-\y+\startY);
    }
}
%
\pgfmathsetmacro{\centerX}{\meshsize/2}
\pgfmathsetmacro{\centerY}{-\meshsize/2}
%
\pgfmathsetmacro{\pollcenterX}{\centerX+\pollX}
\pgfmathsetmacro{\pollcenterY}{\centerY+\pollY}
\pgfmathsetmacro{\pollcornerX}{\pollcenterX - \pollsize/2}
\pgfmathsetmacro{\pollcornerY}{\pollcenterY + \pollsize/2}
%
\pgfmathsetmacro{\pAcoordX}{\pollcenterX + \pAX}
\pgfmathsetmacro{\pAcoordY}{\pollcenterY - \pAY}
\pgfmathsetmacro{\pBcoordX}{\pollcenterX + \pBX}
\pgfmathsetmacro{\pBcoordY}{\pollcenterY - \pBY}
\pgfmathsetmacro{\pCcoordX}{\pollcenterX + \pCX}
\pgfmathsetmacro{\pCcoordY}{\pollcenterY - \pCY}
%
\pgfmathsetmacro{\deltameshn}{\pollcenterX-\pollcornerX}
\pgfmathsetmacro{\deltameshd}{\centerX}
\pgfmathsetmacro{\deltamesh}{(\pollcenterX-\pollcornerX)/\centerX}
\pgfmathsetmacro{\deltaframen}{\deltameshn * \deltameshn}
\pgfmathsetmacro{\deltaframed}{\deltameshd * \deltameshd}
\pgfmathsetmacro{\deltaframe}{\deltamesh * \deltamesh}
%
\Grid{0}{0}{\meshsize}{\meshsize}{\unit}{1pt}
\Grid{\pollcornerX*\unit}{\pollcornerY*\unit}{\pollsize}{\pollsize}{\unit}{3pt}
%
\ifdim \deltamesh pt = 1 pt
    \dimline[label style={above=0.5ex,}, line style = {line width=0.7},extension start length=+0.3,extension end length=+0.3] {(\pollcornerX*\unit,1)}{(\pollcenterX*\unit,1)}{\huge $\delta^{k}_\mathrm{poll}=1$};
    
    \dimline[label style={above=0.5ex,}, line style = {line width=0.7},extension start length=+0.3,extension end length=+0.3] {(\meshsize*\unit - 1*\unit,1)}{(\meshsize*\unit,1)}{\huge $\delta^{k}_\mathrm{mesh}=1$};
\else
    \dimline[label style={above=0.5ex,}, line style = {line width=0.7},extension start length=+0.3,extension end length=+0.3] {(\pollcornerX*\unit,0.5)}{(\pollcenterX*\unit,0.5)}{\huge $\delta^{k}_\mathrm{poll}=\addfraction{\deltameshn}{\deltameshd}$};
    
    \dimline[label style={above=0.5ex,}, line style = {line width=0.7},extension start length=+0.3,extension end length=+0.3] {(\meshsize*\unit - 1*\unit,0.5)}{(\meshsize*\unit,0.5)}{\huge $\delta^{k}_\mathrm{mesh}=\addfraction{\deltaframen}{\deltaframed}$};
\fi
    \node[align=center, minimum width = 100mm, above right] at (0,2) {\huge Poll \textbf{\color{red}failure}};
    \node[circle, minimum size = 1mm, fill=red, label={[fill=white,above right,label distance=0.2cm]\huge $x_k$}] (center) at (\pollcenterX*\unit,\pollcenterY*\unit) {};
    \node[circle, minimum size = 1mm, fill=red, label={[fill=white,below,label distance=-0.5cm]\huge $p_{13}$}] (p1) at (\pAcoordX*\unit,\pAcoordY*\unit) {};
    \node[circle, minimum size = 1mm, fill=red, label={[fill=white,right,label distance=0.5cm]\huge $p_{14}$}] (p2) at (\pBcoordX*\unit,\pBcoordY*\unit) {};
    \node[circle, minimum size = 1mm, fill=red, label={[fill=white]\huge $p_{15}$}] (p3) at (\pCcoordX*\unit,\pCcoordY*\unit) {};
\fi
%
\ifnum\frameopt < 10
    \def\unit{0.15625}
    \def\meshsize{64}
    \def\pollX{-8}
    \def\pollY{-8}
    \def\pollsize{8}
    %
    \def\pAX{1}
    \def\pAY{\pollsize/2}
    \def\pBX{-\pollsize/2}
    \def\pBY{0}
    \def\pCX{\pollsize/2}
    \def\pCY{-\pollsize/2}
\fi
%
\ifnum\frameopt = 9
    % https://tex.stackexchange.com/a/183633
\newcommand{\Grid}[6]{%
    \def\startX{#1}
    \def\startY{#2}
    \def\maxX{#3}
    \def\maxY{#4}
    \def\size{#5}
    \def\borderthick{#6}
    %
    \draw[line width=\borderthick] (\size*0+\startX,\size*0+\startY) rectangle (\size*\maxX+\startX,\size*-\maxY+\startY);
    \foreach \x in {0,1,...,\maxX}{
        \draw (\size*\x+\startX,\size*0+\startY) -- (\size*\x+\startX,\size*-\maxY+\startY);
    }
    %
    \foreach \y in {0,1,...,\maxY}{
        \draw (\size*0+\startX,\size*-\y+\startY) -- (\size*\maxX+\startX,\size*-\y+\startY);
    }
}
%
\pgfmathsetmacro{\centerX}{\meshsize/2}
\pgfmathsetmacro{\centerY}{-\meshsize/2}
%
\pgfmathsetmacro{\pollcenterX}{\centerX+\pollX}
\pgfmathsetmacro{\pollcenterY}{\centerY+\pollY}
\pgfmathsetmacro{\pollcornerX}{\pollcenterX - \pollsize/2}
\pgfmathsetmacro{\pollcornerY}{\pollcenterY + \pollsize/2}
%
\pgfmathsetmacro{\pAcoordX}{\pollcenterX + \pAX}
\pgfmathsetmacro{\pAcoordY}{\pollcenterY - \pAY}
\pgfmathsetmacro{\pBcoordX}{\pollcenterX + \pBX}
\pgfmathsetmacro{\pBcoordY}{\pollcenterY - \pBY}
\pgfmathsetmacro{\pCcoordX}{\pollcenterX + \pCX}
\pgfmathsetmacro{\pCcoordY}{\pollcenterY - \pCY}
%
\pgfmathsetmacro{\deltameshn}{\pollcenterX-\pollcornerX}
\pgfmathsetmacro{\deltameshd}{\centerX}
\pgfmathsetmacro{\deltamesh}{(\pollcenterX-\pollcornerX)/\centerX}
\pgfmathsetmacro{\deltaframen}{\deltameshn * \deltameshn}
\pgfmathsetmacro{\deltaframed}{\deltameshd * \deltameshd}
\pgfmathsetmacro{\deltaframe}{\deltamesh * \deltamesh}
%
\Grid{0}{0}{\meshsize}{\meshsize}{\unit}{1pt}
\Grid{\pollcornerX*\unit}{\pollcornerY*\unit}{\pollsize}{\pollsize}{\unit}{3pt}
%
\ifdim \deltamesh pt = 1 pt
    \dimline[label style={above=0.5ex,}, line style = {line width=0.7},extension start length=+0.3,extension end length=+0.3] {(\pollcornerX*\unit,1)}{(\pollcenterX*\unit,1)}{\huge $\delta^{k}_\mathrm{poll}=1$};
    
    \dimline[label style={above=0.5ex,}, line style = {line width=0.7},extension start length=+0.3,extension end length=+0.3] {(\meshsize*\unit - 1*\unit,1)}{(\meshsize*\unit,1)}{\huge $\delta^{k}_\mathrm{mesh}=1$};
\else
    \dimline[label style={above=0.5ex,}, line style = {line width=0.7},extension start length=+0.3,extension end length=+0.3] {(\pollcornerX*\unit,0.5)}{(\pollcenterX*\unit,0.5)}{\huge $\delta^{k}_\mathrm{poll}=\addfraction{\deltameshn}{\deltameshd}$};
    
    \dimline[label style={above=0.5ex,}, line style = {line width=0.7},extension start length=+0.3,extension end length=+0.3] {(\meshsize*\unit - 1*\unit,0.5)}{(\meshsize*\unit,0.5)}{\huge $\delta^{k}_\mathrm{mesh}=\addfraction{\deltaframen}{\deltaframed}$};
\fi
    \node[align=center, minimum width = 100mm, above right] at (0,2) {\huge Poll \textbf{\color{red}failure}};
    \node[circle, minimum size = 1mm, fill=red] (center) at (\pollcenterX*\unit,\pollcenterY*\unit) {};
    \node[circle, minimum size = 1mm, fill=red] (p1) at (\pAcoordX*\unit,\pAcoordY*\unit) {};
    \node[circle, minimum size = 1mm, fill=red] (p2) at (\pBcoordX*\unit,\pBcoordY*\unit) {};
    \node[circle, minimum size = 1mm, fill=red] (p3) at (\pCcoordX*\unit,\pCcoordY*\unit) {};
\fi
%
\ifnum\frameopt > -1
    \draw[-, color=red, line width=2pt, opacity=1.0] (center) -- (p1);
    \draw[-, color=red, line width=2pt, opacity=1.0] (center) -- (p2);
    \draw[-, color=red, line width=2pt, opacity=1.0] (center) -- (p3);
\fi}%
					\only<3->{\ifnum\frameopts < 2
    \def\unit{1.25}
    \def\meshsize{8}
    \def\pollX{0}
    \def\pollY{0}
    \def\pollsize{4}
    %
    \def\pAX{-\pollsize/2}
    \def\pAY{-1}
    \def\pBX{\pollsize/2}
    \def\pBY{\pollsize/2}
    \def\pCX{\pollsize/2}
    \def\pCY{-1}
\fi
%
\ifnum\frameopts=0
    % https://tex.stackexchange.com/a/183633
\newcommand{\Grid}[6]{%
    \def\startX{#1}
    \def\startY{#2}
    \def\maxX{#3}
    \def\maxY{#4}
    \def\size{#5}
    \def\borderthick{#6}
    %
    \draw[line width=\borderthick] (\size*0+\startX,\size*0+\startY) rectangle (\size*\maxX+\startX,\size*-\maxY+\startY);
    \foreach \x in {0,1,...,\maxX}{
        \draw (\size*\x+\startX,\size*0+\startY) -- (\size*\x+\startX,\size*-\maxY+\startY);
    }
    %
    \foreach \y in {0,1,...,\maxY}{
        \draw (\size*0+\startX,\size*-\y+\startY) -- (\size*\maxX+\startX,\size*-\y+\startY);
    }
}
%
\pgfmathsetmacro{\centerX}{\meshsize/2}
\pgfmathsetmacro{\centerY}{-\meshsize/2}
%
\pgfmathsetmacro{\pollcenterX}{\centerX+\pollX}
\pgfmathsetmacro{\pollcenterY}{\centerY+\pollY}
\pgfmathsetmacro{\pollcornerX}{\pollcenterX - \pollsize/2}
\pgfmathsetmacro{\pollcornerY}{\pollcenterY + \pollsize/2}
%
\pgfmathsetmacro{\pAcoordX}{\pollcenterX + \pAX}
\pgfmathsetmacro{\pAcoordY}{\pollcenterY - \pAY}
\pgfmathsetmacro{\pBcoordX}{\pollcenterX + \pBX}
\pgfmathsetmacro{\pBcoordY}{\pollcenterY - \pBY}
\pgfmathsetmacro{\pCcoordX}{\pollcenterX + \pCX}
\pgfmathsetmacro{\pCcoordY}{\pollcenterY - \pCY}
%
\pgfmathsetmacro{\deltameshn}{\pollcenterX-\pollcornerX}
\pgfmathsetmacro{\deltameshd}{\centerX}
\pgfmathsetmacro{\deltamesh}{(\pollcenterX-\pollcornerX)/\centerX}
\pgfmathsetmacro{\deltaframen}{\deltameshn * \deltameshn}
\pgfmathsetmacro{\deltaframed}{\deltameshd * \deltameshd}
\pgfmathsetmacro{\deltaframe}{\deltamesh * \deltamesh}
%
\Grid{0}{0}{\meshsize}{\meshsize}{\unit}{1pt}
\Grid{\pollcornerX*\unit}{\pollcornerY*\unit}{\pollsize}{\pollsize}{\unit}{3pt}
%
\ifdim \deltamesh pt = 1 pt
    \dimline[label style={above=0.5ex,}, line style = {line width=0.7},extension start length=+0.3,extension end length=+0.3] {(\pollcornerX*\unit,1)}{(\pollcenterX*\unit,1)}{\huge $\delta^{k}_\mathrm{poll}=1$};
    
    \dimline[label style={above=0.5ex,}, line style = {line width=0.7},extension start length=+0.3,extension end length=+0.3] {(\meshsize*\unit - 1*\unit,1)}{(\meshsize*\unit,1)}{\huge $\delta^{k}_\mathrm{mesh}=1$};
\else
    \dimline[label style={above=0.5ex,}, line style = {line width=0.7},extension start length=+0.3,extension end length=+0.3] {(\pollcornerX*\unit,0.5)}{(\pollcenterX*\unit,0.5)}{\huge $\delta^{k}_\mathrm{poll}=\addfraction{\deltameshn}{\deltameshd}$};
    
    \dimline[label style={above=0.5ex,}, line style = {line width=0.7},extension start length=+0.3,extension end length=+0.3] {(\meshsize*\unit - 1*\unit,0.5)}{(\meshsize*\unit,0.5)}{\huge $\delta^{k}_\mathrm{mesh}=\addfraction{\deltaframen}{\deltaframed}$};
\fi
    \node[circle, minimum size = 1mm, fill=red, label={[fill=white,above right,label distance=0.2cm]\huge $x_k$}] (center) at (\pollcenterX*\unit,\pollcenterY*\unit) {};
    \node[circle, minimum size = 1mm, fill=red, label={[fill=white,below,label distance=-0.5cm]\huge $p_4$}] (p1) at (\pAcoordX*\unit,\pAcoordY*\unit) {};
    \node[circle, minimum size = 1mm, fill=red, label={[fill=white,right,label distance=0.5cm]\huge $p_5$}] (p2) at (\pBcoordX*\unit,\pBcoordY*\unit) {};
    \node[circle, minimum size = 1mm, fill=red, label={[fill=white]\huge $p_6$}] (p3) at (\pCcoordX*\unit,\pCcoordY*\unit) {};
\fi
%
\ifnum\frameopts=1
    % https://tex.stackexchange.com/a/183633
\newcommand{\Grid}[6]{%
    \def\startX{#1}
    \def\startY{#2}
    \def\maxX{#3}
    \def\maxY{#4}
    \def\size{#5}
    \def\borderthick{#6}
    %
    \draw[line width=\borderthick] (\size*0+\startX,\size*0+\startY) rectangle (\size*\maxX+\startX,\size*-\maxY+\startY);
    \foreach \x in {0,1,...,\maxX}{
        \draw (\size*\x+\startX,\size*0+\startY) -- (\size*\x+\startX,\size*-\maxY+\startY);
    }
    %
    \foreach \y in {0,1,...,\maxY}{
        \draw (\size*0+\startX,\size*-\y+\startY) -- (\size*\maxX+\startX,\size*-\y+\startY);
    }
}
%
\pgfmathsetmacro{\centerX}{\meshsize/2}
\pgfmathsetmacro{\centerY}{-\meshsize/2}
%
\pgfmathsetmacro{\pollcenterX}{\centerX+\pollX}
\pgfmathsetmacro{\pollcenterY}{\centerY+\pollY}
\pgfmathsetmacro{\pollcornerX}{\pollcenterX - \pollsize/2}
\pgfmathsetmacro{\pollcornerY}{\pollcenterY + \pollsize/2}
%
\pgfmathsetmacro{\pAcoordX}{\pollcenterX + \pAX}
\pgfmathsetmacro{\pAcoordY}{\pollcenterY - \pAY}
\pgfmathsetmacro{\pBcoordX}{\pollcenterX + \pBX}
\pgfmathsetmacro{\pBcoordY}{\pollcenterY - \pBY}
\pgfmathsetmacro{\pCcoordX}{\pollcenterX + \pCX}
\pgfmathsetmacro{\pCcoordY}{\pollcenterY - \pCY}
%
\pgfmathsetmacro{\deltameshn}{\pollcenterX-\pollcornerX}
\pgfmathsetmacro{\deltameshd}{\centerX}
\pgfmathsetmacro{\deltamesh}{(\pollcenterX-\pollcornerX)/\centerX}
\pgfmathsetmacro{\deltaframen}{\deltameshn * \deltameshn}
\pgfmathsetmacro{\deltaframed}{\deltameshd * \deltameshd}
\pgfmathsetmacro{\deltaframe}{\deltamesh * \deltamesh}
%
\Grid{0}{0}{\meshsize}{\meshsize}{\unit}{1pt}
\Grid{\pollcornerX*\unit}{\pollcornerY*\unit}{\pollsize}{\pollsize}{\unit}{3pt}
%
\ifdim \deltamesh pt = 1 pt
    \dimline[label style={above=0.5ex,}, line style = {line width=0.7},extension start length=+0.3,extension end length=+0.3] {(\pollcornerX*\unit,1)}{(\pollcenterX*\unit,1)}{\huge $\delta^{k}_\mathrm{poll}=1$};
    
    \dimline[label style={above=0.5ex,}, line style = {line width=0.7},extension start length=+0.3,extension end length=+0.3] {(\meshsize*\unit - 1*\unit,1)}{(\meshsize*\unit,1)}{\huge $\delta^{k}_\mathrm{mesh}=1$};
\else
    \dimline[label style={above=0.5ex,}, line style = {line width=0.7},extension start length=+0.3,extension end length=+0.3] {(\pollcornerX*\unit,0.5)}{(\pollcenterX*\unit,0.5)}{\huge $\delta^{k}_\mathrm{poll}=\addfraction{\deltameshn}{\deltameshd}$};
    
    \dimline[label style={above=0.5ex,}, line style = {line width=0.7},extension start length=+0.3,extension end length=+0.3] {(\meshsize*\unit - 1*\unit,0.5)}{(\meshsize*\unit,0.5)}{\huge $\delta^{k}_\mathrm{mesh}=\addfraction{\deltaframen}{\deltaframed}$};
\fi
    \node[align=center, minimum width = 100mm, above right] at (0,2) {\huge Poll \textbf{certain} \textbf{\color{red}failure}};
    \node[star, star points=8, minimum size = 6mm, fill=blue, label={[fill=white,above right,label distance=0.2cm]\huge $x_k$}] (center) at (\pollcenterX*\unit,\pollcenterY*\unit) {};
    \node[circle, minimum size = 1mm, fill=red, label={[fill=white,below,label distance=-0.5cm]\huge $p_4$}] (p1) at (\pAcoordX*\unit,\pAcoordY*\unit) {};
    \node[circle, minimum size = 1mm, fill=red, label={[fill=white,right,label distance=0.5cm]\huge $p_5$}] (p2) at (\pBcoordX*\unit,\pBcoordY*\unit) {};
    \node[circle, minimum size = 1mm, fill=red, label={[fill=white]\huge $p_6$}] (p3) at (\pCcoordX*\unit,\pCcoordY*\unit) {};
    \node[circle, minimum size = 1.41*\unit*20mm, fill=blue, opacity=0.5] (uncertainty) at (\pollcenterX*\unit,\pollcenterY*\unit) {};
\fi
%
\ifnum\frameopts < 3
    \def\unit{0.3125}
    \def\meshsize{32}
    \def\pollX{0}
    \def\pollY{0}
    \def\pollsize{8}
    %
    \def\pAX{-\pollsize/2}
    \def\pAY{-3}
    \def\pBX{3}
    \def\pBY{\pollsize/2}
    \def\pCX{\pollsize/2}
    \def\pCY{-\pollsize/2}
\fi
%
\ifnum\frameopts=2
    % https://tex.stackexchange.com/a/183633
\newcommand{\Grid}[6]{%
    \def\startX{#1}
    \def\startY{#2}
    \def\maxX{#3}
    \def\maxY{#4}
    \def\size{#5}
    \def\borderthick{#6}
    %
    \draw[line width=\borderthick] (\size*0+\startX,\size*0+\startY) rectangle (\size*\maxX+\startX,\size*-\maxY+\startY);
    \foreach \x in {0,1,...,\maxX}{
        \draw (\size*\x+\startX,\size*0+\startY) -- (\size*\x+\startX,\size*-\maxY+\startY);
    }
    %
    \foreach \y in {0,1,...,\maxY}{
        \draw (\size*0+\startX,\size*-\y+\startY) -- (\size*\maxX+\startX,\size*-\y+\startY);
    }
}
%
\pgfmathsetmacro{\centerX}{\meshsize/2}
\pgfmathsetmacro{\centerY}{-\meshsize/2}
%
\pgfmathsetmacro{\pollcenterX}{\centerX+\pollX}
\pgfmathsetmacro{\pollcenterY}{\centerY+\pollY}
\pgfmathsetmacro{\pollcornerX}{\pollcenterX - \pollsize/2}
\pgfmathsetmacro{\pollcornerY}{\pollcenterY + \pollsize/2}
%
\pgfmathsetmacro{\pAcoordX}{\pollcenterX + \pAX}
\pgfmathsetmacro{\pAcoordY}{\pollcenterY - \pAY}
\pgfmathsetmacro{\pBcoordX}{\pollcenterX + \pBX}
\pgfmathsetmacro{\pBcoordY}{\pollcenterY - \pBY}
\pgfmathsetmacro{\pCcoordX}{\pollcenterX + \pCX}
\pgfmathsetmacro{\pCcoordY}{\pollcenterY - \pCY}
%
\pgfmathsetmacro{\deltameshn}{\pollcenterX-\pollcornerX}
\pgfmathsetmacro{\deltameshd}{\centerX}
\pgfmathsetmacro{\deltamesh}{(\pollcenterX-\pollcornerX)/\centerX}
\pgfmathsetmacro{\deltaframen}{\deltameshn * \deltameshn}
\pgfmathsetmacro{\deltaframed}{\deltameshd * \deltameshd}
\pgfmathsetmacro{\deltaframe}{\deltamesh * \deltamesh}
%
\Grid{0}{0}{\meshsize}{\meshsize}{\unit}{1pt}
\Grid{\pollcornerX*\unit}{\pollcornerY*\unit}{\pollsize}{\pollsize}{\unit}{3pt}
%
\ifdim \deltamesh pt = 1 pt
    \dimline[label style={above=0.5ex,}, line style = {line width=0.7},extension start length=+0.3,extension end length=+0.3] {(\pollcornerX*\unit,1)}{(\pollcenterX*\unit,1)}{\huge $\delta^{k}_\mathrm{poll}=1$};
    
    \dimline[label style={above=0.5ex,}, line style = {line width=0.7},extension start length=+0.3,extension end length=+0.3] {(\meshsize*\unit - 1*\unit,1)}{(\meshsize*\unit,1)}{\huge $\delta^{k}_\mathrm{mesh}=1$};
\else
    \dimline[label style={above=0.5ex,}, line style = {line width=0.7},extension start length=+0.3,extension end length=+0.3] {(\pollcornerX*\unit,0.5)}{(\pollcenterX*\unit,0.5)}{\huge $\delta^{k}_\mathrm{poll}=\addfraction{\deltameshn}{\deltameshd}$};
    
    \dimline[label style={above=0.5ex,}, line style = {line width=0.7},extension start length=+0.3,extension end length=+0.3] {(\meshsize*\unit - 1*\unit,0.5)}{(\meshsize*\unit,0.5)}{\huge $\delta^{k}_\mathrm{mesh}=\addfraction{\deltaframen}{\deltaframed}$};
\fi
    \node[align=center, minimum width = 100mm, above right] at (0,2) {\huge Poll \textbf{certain} \textbf{\color{red}failure}};
    \node[star, star points=8, minimum size = 6mm, fill=blue, label={[fill=white,above right,label distance=0.2cm]\huge $x_k$}] (center) at (\pollcenterX*\unit,\pollcenterY*\unit) {};
    \node[circle, minimum size = 1mm, fill=red, label={[fill=white,below,label distance=-0.5cm]\huge $p_7$}] (p1) at (\pAcoordX*\unit,\pAcoordY*\unit) {};
    \node[circle, minimum size = 1mm, fill=red, label={[fill=white,right,label distance=0.5cm]\huge $p_8$}] (p2) at (\pBcoordX*\unit,\pBcoordY*\unit) {};
    \node[circle, minimum size = 1mm, fill=red, label={[fill=white]\huge $p_9$}] (p3) at (\pCcoordX*\unit,\pCcoordY*\unit) {};
    \node[circle, minimum size = 1.41*\unit*20mm, fill=blue, opacity=0.5] (uncertainty) at (\pollcenterX*\unit,\pollcenterY*\unit) {};
\fi
\ifnum\frameopts < 5
    \def\unit{1.25}
    \def\meshsize{8}
    \def\pollX{0}
    \def\pollY{0}
    \def\pollsize{4}
    %
    \def\pAX{-\pollsize/2}
    \def\pAY{-1}
    \def\pBX{\pollsize/2}
    \def\pBY{\pollsize/2}
    \def\pCX{\pollsize/2}
    \def\pCY{-1}
\fi
%
\ifnum\frameopts=3
    % https://tex.stackexchange.com/a/183633
\newcommand{\Grid}[6]{%
    \def\startX{#1}
    \def\startY{#2}
    \def\maxX{#3}
    \def\maxY{#4}
    \def\size{#5}
    \def\borderthick{#6}
    %
    \draw[line width=\borderthick] (\size*0+\startX,\size*0+\startY) rectangle (\size*\maxX+\startX,\size*-\maxY+\startY);
    \foreach \x in {0,1,...,\maxX}{
        \draw (\size*\x+\startX,\size*0+\startY) -- (\size*\x+\startX,\size*-\maxY+\startY);
    }
    %
    \foreach \y in {0,1,...,\maxY}{
        \draw (\size*0+\startX,\size*-\y+\startY) -- (\size*\maxX+\startX,\size*-\y+\startY);
    }
}
%
\pgfmathsetmacro{\centerX}{\meshsize/2}
\pgfmathsetmacro{\centerY}{-\meshsize/2}
%
\pgfmathsetmacro{\pollcenterX}{\centerX+\pollX}
\pgfmathsetmacro{\pollcenterY}{\centerY+\pollY}
\pgfmathsetmacro{\pollcornerX}{\pollcenterX - \pollsize/2}
\pgfmathsetmacro{\pollcornerY}{\pollcenterY + \pollsize/2}
%
\pgfmathsetmacro{\pAcoordX}{\pollcenterX + \pAX}
\pgfmathsetmacro{\pAcoordY}{\pollcenterY - \pAY}
\pgfmathsetmacro{\pBcoordX}{\pollcenterX + \pBX}
\pgfmathsetmacro{\pBcoordY}{\pollcenterY - \pBY}
\pgfmathsetmacro{\pCcoordX}{\pollcenterX + \pCX}
\pgfmathsetmacro{\pCcoordY}{\pollcenterY - \pCY}
%
\pgfmathsetmacro{\deltameshn}{\pollcenterX-\pollcornerX}
\pgfmathsetmacro{\deltameshd}{\centerX}
\pgfmathsetmacro{\deltamesh}{(\pollcenterX-\pollcornerX)/\centerX}
\pgfmathsetmacro{\deltaframen}{\deltameshn * \deltameshn}
\pgfmathsetmacro{\deltaframed}{\deltameshd * \deltameshd}
\pgfmathsetmacro{\deltaframe}{\deltamesh * \deltamesh}
%
\Grid{0}{0}{\meshsize}{\meshsize}{\unit}{1pt}
\Grid{\pollcornerX*\unit}{\pollcornerY*\unit}{\pollsize}{\pollsize}{\unit}{3pt}
%
\ifdim \deltamesh pt = 1 pt
    \dimline[label style={above=0.5ex,}, line style = {line width=0.7},extension start length=+0.3,extension end length=+0.3] {(\pollcornerX*\unit,1)}{(\pollcenterX*\unit,1)}{\huge $\delta^{k}_\mathrm{poll}=1$};
    
    \dimline[label style={above=0.5ex,}, line style = {line width=0.7},extension start length=+0.3,extension end length=+0.3] {(\meshsize*\unit - 1*\unit,1)}{(\meshsize*\unit,1)}{\huge $\delta^{k}_\mathrm{mesh}=1$};
\else
    \dimline[label style={above=0.5ex,}, line style = {line width=0.7},extension start length=+0.3,extension end length=+0.3] {(\pollcornerX*\unit,0.5)}{(\pollcenterX*\unit,0.5)}{\huge $\delta^{k}_\mathrm{poll}=\addfraction{\deltameshn}{\deltameshd}$};
    
    \dimline[label style={above=0.5ex,}, line style = {line width=0.7},extension start length=+0.3,extension end length=+0.3] {(\meshsize*\unit - 1*\unit,0.5)}{(\meshsize*\unit,0.5)}{\huge $\delta^{k}_\mathrm{mesh}=\addfraction{\deltaframen}{\deltaframed}$};
\fi
    \node[align=center, minimum width = 100mm, above right] at (0,2) {\huge Poll \textbf{certain} \textbf{\color{red}failure}};
    \node[star, star points=8, minimum size = 6mm, fill=blue, label={[fill=white,above right,label distance=0.2cm]\huge $x_k$}] (center) at (\pollcenterX*\unit,\pollcenterY*\unit) {};
    \node[circle, minimum size = 1mm, fill=red, label={[fill=white,below,label distance=-0.5cm]\huge $p_4$}] (p1) at (\pAcoordX*\unit,\pAcoordY*\unit) {};
    \node[circle, minimum size = 1mm, fill=red, label={[fill=white,right,label distance=0.5cm]\huge $p_5$}] (p2) at (\pBcoordX*\unit,\pBcoordY*\unit) {};
    \node[circle, minimum size = 1mm, fill=red, label={[fill=white]\huge $p_6$}] (p3) at (\pCcoordX*\unit,\pCcoordY*\unit) {};
    \node[circle, minimum size = 1.41*\unit*20mm, fill=blue, opacity=0.5] (uncertainty) at (\pollcenterX*\unit,\pollcenterY*\unit) {};
\fi
%
\ifnum\frameopts=4
    % https://tex.stackexchange.com/a/183633
\newcommand{\Grid}[6]{%
    \def\startX{#1}
    \def\startY{#2}
    \def\maxX{#3}
    \def\maxY{#4}
    \def\size{#5}
    \def\borderthick{#6}
    %
    \draw[line width=\borderthick] (\size*0+\startX,\size*0+\startY) rectangle (\size*\maxX+\startX,\size*-\maxY+\startY);
    \foreach \x in {0,1,...,\maxX}{
        \draw (\size*\x+\startX,\size*0+\startY) -- (\size*\x+\startX,\size*-\maxY+\startY);
    }
    %
    \foreach \y in {0,1,...,\maxY}{
        \draw (\size*0+\startX,\size*-\y+\startY) -- (\size*\maxX+\startX,\size*-\y+\startY);
    }
}
%
\pgfmathsetmacro{\centerX}{\meshsize/2}
\pgfmathsetmacro{\centerY}{-\meshsize/2}
%
\pgfmathsetmacro{\pollcenterX}{\centerX+\pollX}
\pgfmathsetmacro{\pollcenterY}{\centerY+\pollY}
\pgfmathsetmacro{\pollcornerX}{\pollcenterX - \pollsize/2}
\pgfmathsetmacro{\pollcornerY}{\pollcenterY + \pollsize/2}
%
\pgfmathsetmacro{\pAcoordX}{\pollcenterX + \pAX}
\pgfmathsetmacro{\pAcoordY}{\pollcenterY - \pAY}
\pgfmathsetmacro{\pBcoordX}{\pollcenterX + \pBX}
\pgfmathsetmacro{\pBcoordY}{\pollcenterY - \pBY}
\pgfmathsetmacro{\pCcoordX}{\pollcenterX + \pCX}
\pgfmathsetmacro{\pCcoordY}{\pollcenterY - \pCY}
%
\pgfmathsetmacro{\deltameshn}{\pollcenterX-\pollcornerX}
\pgfmathsetmacro{\deltameshd}{\centerX}
\pgfmathsetmacro{\deltamesh}{(\pollcenterX-\pollcornerX)/\centerX}
\pgfmathsetmacro{\deltaframen}{\deltameshn * \deltameshn}
\pgfmathsetmacro{\deltaframed}{\deltameshd * \deltameshd}
\pgfmathsetmacro{\deltaframe}{\deltamesh * \deltamesh}
%
\Grid{0}{0}{\meshsize}{\meshsize}{\unit}{1pt}
\Grid{\pollcornerX*\unit}{\pollcornerY*\unit}{\pollsize}{\pollsize}{\unit}{3pt}
%
\ifdim \deltamesh pt = 1 pt
    \dimline[label style={above=0.5ex,}, line style = {line width=0.7},extension start length=+0.3,extension end length=+0.3] {(\pollcornerX*\unit,1)}{(\pollcenterX*\unit,1)}{\huge $\delta^{k}_\mathrm{poll}=1$};
    
    \dimline[label style={above=0.5ex,}, line style = {line width=0.7},extension start length=+0.3,extension end length=+0.3] {(\meshsize*\unit - 1*\unit,1)}{(\meshsize*\unit,1)}{\huge $\delta^{k}_\mathrm{mesh}=1$};
\else
    \dimline[label style={above=0.5ex,}, line style = {line width=0.7},extension start length=+0.3,extension end length=+0.3] {(\pollcornerX*\unit,0.5)}{(\pollcenterX*\unit,0.5)}{\huge $\delta^{k}_\mathrm{poll}=\addfraction{\deltameshn}{\deltameshd}$};
    
    \dimline[label style={above=0.5ex,}, line style = {line width=0.7},extension start length=+0.3,extension end length=+0.3] {(\meshsize*\unit - 1*\unit,0.5)}{(\meshsize*\unit,0.5)}{\huge $\delta^{k}_\mathrm{mesh}=\addfraction{\deltaframen}{\deltaframed}$};
\fi
    \node[align=center, minimum width = 100mm, above right] at (0,2) {\huge Poll \textbf{uncertain} \textbf{\color{red}failure}};
    \node[star, star points=8, minimum size = 6mm, fill=blue, label={[fill=white,above right,label distance=0.2cm]\huge $x_k$}] (center) at (\pollcenterX*\unit,\pollcenterY*\unit) {};
    \node[circle, minimum size = 1mm, fill=red, label={[fill=white,below,label distance=-0.5cm]\huge $p_4$}] (p1) at (\pAcoordX*\unit,\pAcoordY*\unit) {};
    \node[circle, minimum size = 1mm, fill=red, label={[fill=white,right,label distance=0.5cm]\huge $p_5$}] (p2) at (\pBcoordX*\unit,\pBcoordY*\unit) {};
    \node[circle, minimum size = 1mm, fill=red, label={[fill=white]\huge $p_6$}] (p3) at (\pCcoordX*\unit,\pCcoordY*\unit) {};
    \node[circle, minimum size = 2*1.21*\unit*20mm, fill=blue, opacity=0.5] (uncertainty) at (\pollcenterX*\unit,\pollcenterY*\unit) {};
\fi
%
\ifnum\frameopts < 6
    \def\unit{0.625}
    \def\meshsize{16}
    \def\pollX{0}
    \def\pollY{0}
    \def\pollsize{6}
    %
    \def\pAX{-\pollsize/2}
    \def\pAY{-1}
    \def\pBX{\pollsize/2}
    \def\pBY{\pollsize/2}
    \def\pCX{\pollsize/2}
    \def\pCY{-1}
\fi
%
\ifnum\frameopts=5
    % https://tex.stackexchange.com/a/183633
\newcommand{\Grid}[6]{%
    \def\startX{#1}
    \def\startY{#2}
    \def\maxX{#3}
    \def\maxY{#4}
    \def\size{#5}
    \def\borderthick{#6}
    %
    \draw[line width=\borderthick] (\size*0+\startX,\size*0+\startY) rectangle (\size*\maxX+\startX,\size*-\maxY+\startY);
    \foreach \x in {0,1,...,\maxX}{
        \draw (\size*\x+\startX,\size*0+\startY) -- (\size*\x+\startX,\size*-\maxY+\startY);
    }
    %
    \foreach \y in {0,1,...,\maxY}{
        \draw (\size*0+\startX,\size*-\y+\startY) -- (\size*\maxX+\startX,\size*-\y+\startY);
    }
}
%
\pgfmathsetmacro{\centerX}{\meshsize/2}
\pgfmathsetmacro{\centerY}{-\meshsize/2}
%
\pgfmathsetmacro{\pollcenterX}{\centerX+\pollX}
\pgfmathsetmacro{\pollcenterY}{\centerY+\pollY}
\pgfmathsetmacro{\pollcornerX}{\pollcenterX - \pollsize/2}
\pgfmathsetmacro{\pollcornerY}{\pollcenterY + \pollsize/2}
%
\pgfmathsetmacro{\pAcoordX}{\pollcenterX + \pAX}
\pgfmathsetmacro{\pAcoordY}{\pollcenterY - \pAY}
\pgfmathsetmacro{\pBcoordX}{\pollcenterX + \pBX}
\pgfmathsetmacro{\pBcoordY}{\pollcenterY - \pBY}
\pgfmathsetmacro{\pCcoordX}{\pollcenterX + \pCX}
\pgfmathsetmacro{\pCcoordY}{\pollcenterY - \pCY}
%
\pgfmathsetmacro{\deltameshn}{\pollcenterX-\pollcornerX}
\pgfmathsetmacro{\deltameshd}{\centerX}
\pgfmathsetmacro{\deltamesh}{(\pollcenterX-\pollcornerX)/\centerX}
\pgfmathsetmacro{\deltaframen}{\deltameshn * \deltameshn}
\pgfmathsetmacro{\deltaframed}{\deltameshd * \deltameshd}
\pgfmathsetmacro{\deltaframe}{\deltamesh * \deltamesh}
%
\Grid{0}{0}{\meshsize}{\meshsize}{\unit}{1pt}
\Grid{\pollcornerX*\unit}{\pollcornerY*\unit}{\pollsize}{\pollsize}{\unit}{3pt}
%
\ifdim \deltamesh pt = 1 pt
    \dimline[label style={above=0.5ex,}, line style = {line width=0.7},extension start length=+0.3,extension end length=+0.3] {(\pollcornerX*\unit,1)}{(\pollcenterX*\unit,1)}{\huge $\delta^{k}_\mathrm{poll}=1$};
    
    \dimline[label style={above=0.5ex,}, line style = {line width=0.7},extension start length=+0.3,extension end length=+0.3] {(\meshsize*\unit - 1*\unit,1)}{(\meshsize*\unit,1)}{\huge $\delta^{k}_\mathrm{mesh}=1$};
\else
    \dimline[label style={above=0.5ex,}, line style = {line width=0.7},extension start length=+0.3,extension end length=+0.3] {(\pollcornerX*\unit,0.5)}{(\pollcenterX*\unit,0.5)}{\huge $\delta^{k}_\mathrm{poll}=\addfraction{\deltameshn}{\deltameshd}$};
    
    \dimline[label style={above=0.5ex,}, line style = {line width=0.7},extension start length=+0.3,extension end length=+0.3] {(\meshsize*\unit - 1*\unit,0.5)}{(\meshsize*\unit,0.5)}{\huge $\delta^{k}_\mathrm{mesh}=\addfraction{\deltaframen}{\deltaframed}$};
\fi
    \node[align=center, minimum width = 100mm, above right] at (0,2) {\huge Poll \textbf{uncertain} \textbf{\color{red}failure}};
    \node[star, star points=8, minimum size = 6mm, fill=blue, label={[fill=white,above right,label distance=0.2cm]\huge $x_k$}] (center) at (\pollcenterX*\unit,\pollcenterY*\unit) {};
    \node[circle, minimum size = 1mm, fill=red, label={[fill=white,below,label distance=-0.5cm]\huge $p_7$}] (p1) at (\pAcoordX*\unit,\pAcoordY*\unit) {};
    \node[circle, minimum size = 1mm, fill=red, label={[fill=white,right,label distance=0.5cm]\huge $p_8$}] (p2) at (\pBcoordX*\unit,\pBcoordY*\unit) {};
    \node[circle, minimum size = 1mm, fill=red, label={[fill=white]\huge $p_9$}] (p3) at (\pCcoordX*\unit,\pCcoordY*\unit) {};
    \node[circle, minimum size = 2.5*1.41*\unit*20mm, fill=blue, opacity=0.5] (uncertainty) at (\pollcenterX*\unit,\pollcenterY*\unit) {};
\fi
%
\ifnum\frameopts < 7
    \def\unit{0.3125}
    \def\meshsize{32}
    \def\pollX{0}
    \def\pollY{0}
    \def\pollsize{8}
    %
    \def\pAX{-\pollsize/2}
    \def\pAY{-1}
    \def\pBX{\pollsize/2}
    \def\pBY{\pollsize/2}
    \def\pCX{\pollsize/2}
    \def\pCY{-1}
\fi
%
\ifnum\frameopts=6
    % https://tex.stackexchange.com/a/183633
\newcommand{\Grid}[6]{%
    \def\startX{#1}
    \def\startY{#2}
    \def\maxX{#3}
    \def\maxY{#4}
    \def\size{#5}
    \def\borderthick{#6}
    %
    \draw[line width=\borderthick] (\size*0+\startX,\size*0+\startY) rectangle (\size*\maxX+\startX,\size*-\maxY+\startY);
    \foreach \x in {0,1,...,\maxX}{
        \draw (\size*\x+\startX,\size*0+\startY) -- (\size*\x+\startX,\size*-\maxY+\startY);
    }
    %
    \foreach \y in {0,1,...,\maxY}{
        \draw (\size*0+\startX,\size*-\y+\startY) -- (\size*\maxX+\startX,\size*-\y+\startY);
    }
}
%
\pgfmathsetmacro{\centerX}{\meshsize/2}
\pgfmathsetmacro{\centerY}{-\meshsize/2}
%
\pgfmathsetmacro{\pollcenterX}{\centerX+\pollX}
\pgfmathsetmacro{\pollcenterY}{\centerY+\pollY}
\pgfmathsetmacro{\pollcornerX}{\pollcenterX - \pollsize/2}
\pgfmathsetmacro{\pollcornerY}{\pollcenterY + \pollsize/2}
%
\pgfmathsetmacro{\pAcoordX}{\pollcenterX + \pAX}
\pgfmathsetmacro{\pAcoordY}{\pollcenterY - \pAY}
\pgfmathsetmacro{\pBcoordX}{\pollcenterX + \pBX}
\pgfmathsetmacro{\pBcoordY}{\pollcenterY - \pBY}
\pgfmathsetmacro{\pCcoordX}{\pollcenterX + \pCX}
\pgfmathsetmacro{\pCcoordY}{\pollcenterY - \pCY}
%
\pgfmathsetmacro{\deltameshn}{\pollcenterX-\pollcornerX}
\pgfmathsetmacro{\deltameshd}{\centerX}
\pgfmathsetmacro{\deltamesh}{(\pollcenterX-\pollcornerX)/\centerX}
\pgfmathsetmacro{\deltaframen}{\deltameshn * \deltameshn}
\pgfmathsetmacro{\deltaframed}{\deltameshd * \deltameshd}
\pgfmathsetmacro{\deltaframe}{\deltamesh * \deltamesh}
%
\Grid{0}{0}{\meshsize}{\meshsize}{\unit}{1pt}
\Grid{\pollcornerX*\unit}{\pollcornerY*\unit}{\pollsize}{\pollsize}{\unit}{3pt}
%
\ifdim \deltamesh pt = 1 pt
    \dimline[label style={above=0.5ex,}, line style = {line width=0.7},extension start length=+0.3,extension end length=+0.3] {(\pollcornerX*\unit,1)}{(\pollcenterX*\unit,1)}{\huge $\delta^{k}_\mathrm{poll}=1$};
    
    \dimline[label style={above=0.5ex,}, line style = {line width=0.7},extension start length=+0.3,extension end length=+0.3] {(\meshsize*\unit - 1*\unit,1)}{(\meshsize*\unit,1)}{\huge $\delta^{k}_\mathrm{mesh}=1$};
\else
    \dimline[label style={above=0.5ex,}, line style = {line width=0.7},extension start length=+0.3,extension end length=+0.3] {(\pollcornerX*\unit,0.5)}{(\pollcenterX*\unit,0.5)}{\huge $\delta^{k}_\mathrm{poll}=\addfraction{\deltameshn}{\deltameshd}$};
    
    \dimline[label style={above=0.5ex,}, line style = {line width=0.7},extension start length=+0.3,extension end length=+0.3] {(\meshsize*\unit - 1*\unit,0.5)}{(\meshsize*\unit,0.5)}{\huge $\delta^{k}_\mathrm{mesh}=\addfraction{\deltaframen}{\deltaframed}$};
\fi
    \node[align=center, minimum width = 100mm, above right] at (0,2) {\huge Poll \textbf{certain} \textbf{\color{red}failure}};
    \node[star, star points=8, minimum size = 6mm, fill=blue, label={[fill=white,above right,label distance=0.2cm]\huge $x_k$}] (center) at (\pollcenterX*\unit,\pollcenterY*\unit) {};
    \node[circle, minimum size = 1mm, fill=red] (p1) at (\pAcoordX*\unit,\pAcoordY*\unit) {};
    \node[circle, minimum size = 1mm, fill=red] (p2) at (\pBcoordX*\unit,\pBcoordY*\unit) {};
    \node[circle, minimum size = 1mm, fill=red] (p3) at (\pCcoordX*\unit,\pCcoordY*\unit) {};
    \node[circle, minimum size = 2.5*1.41*\unit*20mm, fill=blue, opacity=0.5] (uncertainty) at (\pollcenterX*\unit,\pollcenterY*\unit) {};
\fi
%
\draw[-, color=red, line width=2pt, opacity=1.0] (center) -- (p1);
\draw[-, color=red, line width=2pt, opacity=1.0] (center) -- (p2);
\draw[-, color=red, line width=2pt, opacity=1.0] (center) -- (p3);
%
\node[align=center, minimum width = 100mm, above right] (ulabel) at (0,-12) {\huge $\left\{x:f^k_s - f^k_0 \in \mathcal{I}\right\}$};
%
\path[stealth-,blue!70,line width = 3pt] (uncertainty.south) edge [out=290, in=70] (ulabel.north);}%
				\end{tikzpicture}%
			}%
		\end{column}
	\end{columns}
	\vspace{-3em}
\end{frame}
\addtocounter{footnote}{-2}
%------------------------------------------------
\subsection{ABM Conclusion}
%------------------------------------------------
\begin{frame}[t,label=abm_7]
	\frametitle{Conclusion}
	\tikzstyle{background grid}=[draw, black!50,step=.5cm]
	%
	\uncover<1->{Formulated and solved public health policy-making problems}%
	\only<6->{{\color{white}\ifshowcitations\footpartcite{Ferguson2020}\fi}}\\
	%
	\begin{columns}[t] % The "c" option specifies centered vertical alignment while the "t" option is used for top vertical alignment
		\begin{column}{.42\textwidth} % Left column and width
			\begin{itemize}
				\item<1-> Identified a public health policy that favored
				\begin{itemize}
					\item<2-> High testing capacity $n_T$
					\item<3-> Large number of essential workers $n_E$
					\item<4-> Modest social distancing $S_D$
				\end{itemize}
				\item<5-> StoMADS outperformed GAs and NOMAD on active constraints
				\item<6-> Applicable to other large-scale agent-based models\footnotemark[1]
			\end{itemize}
		\end{column}
		%
		\begin{column}{.5\textwidth} % Left column and width
			\tikzstyle{background grid}=[draw, black!50,step=.5cm]
			\begin{tikzpicture}[remember picture, overlay]%[show background grid]
				\node [inner sep=0pt,above right, opacity=1.0]  at (0.0\textwidth,-0.75\textheight) (mobility) 
				{
					\only<-5>{\includegraphics[width=0.93\textwidth]{agent_based_models/trajecctory_results/M_compare_opt_2.pdf}}%
				};
				\only<-5>{
					\node[inner sep=0pt,align=flush center,above=\belowcaptionskip of mobility,text width=\linewidth]
					{\vspace{-0em}{
						\large best known solution
					}};
				}%
				% show origin
				% \fill (0,0) circle (2pt);
			\end{tikzpicture}%
			\begin{tikzpicture}[remember picture, overlay] %show background grid, 
				\node [inner sep=0pt,above right, opacity=1.0]  at (0.1\textwidth,-0.66\textheight) (abm) 
					{
						\only<6>{
							\begin{animateinline}[autoplay,width=0.9\textwidth]{8}
								\ifshowanimations
									\multiframe{48}{i=1+15}{%
										\includegraphics{agent_based_models/canada_ABM_anim/map_\i.png}
									}
								\else
									\multiframe{1}{i=721+0}{%
										\includegraphics{agent_based_models/canada_ABM/map_\i.png}
									}
								\fi
							\end{animateinline}%
						}%
						\only<7->{
							\includegraphics[width=0.9\textwidth]{agent_based_models/canada_ABM/map_721.png}
						}
					};
			\end{tikzpicture}%
		\end{column}
	
	\end{columns}
	\vspace{-3em}
\end{frame}
\addtocounter{footnote}{-1}
%------------------------------------------------
\subsection{COVID-19 introduction}
%------------------------------------------------
\begin{frame}[t]
	\frametitle{Background: COVID-19 forecasting targets}
	\tikzstyle{background grid}=[draw, black!50,step=.5cm]
	%
	Forecasting novel epidemics is a multidisciplinary field involving multiple \emph{targets} \ifshowcitations\footpartcite{Wu2021}\fi\\
	\emph{Inputs} for forecasting \emphasis{epidemic size:}\\
    %
    \begin{itemize}
        \item<2-> Growth rate indicators
        \item<3-> \only<5->{\emphasis}{Historical incidence rate data}
        \item<4-> \only<5->{\emphasis}{Serologic assays}\\
        \centering\includegraphics[width=0.23\textwidth]{machine_learning/targets/pcr.png}
    \end{itemize}
    %
	\tikzstyle{background grid}=[draw, black!50,step=.5cm]
	\begin{tikzpicture}[remember picture, overlay] %show background grid, 
		% Put the graphic inside a node. This makes it easy to place the
		% graphic and to draw on top of it. 
		% The above right option is used to place the lower left corner
		% of the image at the (0,0) coordinate. 
		\node [inner sep=0pt,above left, opacity=1.0]  at (0.99\textwidth,-0.0\textheight) (seq2seq) 
			{
				\only<1>{\includegraphics[width=0.45\textwidth]{machine_learning/targets/nowcasting_targets_1.pdf}}%
				\only<2>{\includegraphics[width=0.45\textwidth]{machine_learning/targets/nowcasting_targets_2.pdf}}%
                \only<3->{\includegraphics[width=0.45\textwidth]{machine_learning/targets/nowcasting_targets_3.pdf}}%
			};
		\only<1->{
			\node[inner sep=0pt,align=flush center,below=\belowcaptionskip of seq2seq,text width=\linewidth]
				{\vspace{-1em}{Nowcasting targets}};
		}%
		% show origin
		% \fill (0,0) circle (2pt);
	\end{tikzpicture}%
	%
	\vspace{-3em}
\end{frame}
\addtocounter{footnote}{-1}
%------------------------------------------------
\subsection{Hyperparameter tuning future directions}
%------------------------------------------------
\begin{frame}[t]
	\frametitle{Conclusion and future directions}
	\tikzstyle{background grid}=[draw, black!50,step=.5cm]
	%
	\only<1-5>{Model discovery and development facilitated by hyperparameter optimization}%
	\only<6->{StoMADS can be improved to solve a wide variety of HPO problems}%
	\only<5>{{\color{white}\ifshowcitations\footpartcite{Hay2021}\fi}}%
	\only<7->{{\color{white}\ifshowcitations\footpartcite{Lakhmiri2019}\textsuperscript{,}\footpartcite{Bergstra2011}\fi}}\\
	%
	\tikzstyle{background grid}=[draw, black!50,step=.5cm]
	\begin{tikzpicture}[remember picture, overlay] %show background grid, 
		% Put the graphic  inside a node. This makes it easy to place the
		% graphic and to draw on top of it. 
		% The above right option is used to place the lower left corner
		% of the image at the (0,0) coordinate. 
		\node [inner sep=0pt,above left, opacity=1.0]  at (1.01\textwidth,-0.73\textheight) (prediction) 
			{
				\only<3->{\includegraphics[width=0.5\textwidth]{machine_learning/predictions/model_predictions_unseen_SVR_mean_Ct_daily_cases_G12.pdf}}%
			};
		% show origin
		% \fill (0,0) circle (2pt);
	\end{tikzpicture}%
	
	\begin{columns}[c] % The "c" option specifies centered vertical alignment while the "t" option is used for top vertical alignment
		\begin{column}{.5\textwidth} % Left column and width
			\vspace{-0.0em}
			% Optimization problem
			\begin{itemize}\itemsep0em
			\only<-5>{
				\item<2-5> Cycle Threshold (Ct) is a useful feature for incidence projection
				\item<3-5> Model that generalizes well on unseen data
    			\item<4-5> Simpler models perform well when historical data is limited
				\item<5> Works well on other datasets\footnotemark[1]
			}
			\item<6-> Can be used to meet deployment targets
				\begin{exampleblock}{Objective and constraints}
					\vspace{-1.0em}%
					\begin{equation*}
						\begin{aligned}
							& \underset{\mathbf{x}}{\text{min}}
							& & f(\mathbf{x}) = \mathbb{E}_{\Theta}\left[{f}_{\Theta}(\mathbf{x}) = \mathrm{error}_\mathrm{CV}\right]\\
							& \text{subject to}
							& & {c}(\mathbf{x}) = \mathrm{\small inference~time} - \mathrm{\small threshold} \le 0\\
							& \text{where}
							& & \Theta\mathrm{:realizations}
						\end{aligned}
					\end{equation*}
				\end{exampleblock}
				\item<7-> Should be benchmarked against HyperNOMAD\footnotemark[1], Bayesian optimization\footnotemark[2]
				\item<8-> Mixed variable version is needed
			\end{itemize}
			\only<-5>{\vspace{-10em}}
		\end{column}
		%
		\begin{column}{.5\textwidth} % Left column and width
		\end{column}
	
	\end{columns}
	%
	\vspace{-3em}
\end{frame}
\addtocounter{footnote}{-2}
%------------------------------------------------

%----------------------------------------------------------------------------------------

\end{document} 